%=======================================================================
%  Master of Engineering — Final Year Project (PFE) Report
%
%  Quadrotor Control by NMPC and Reinforcement Learning
%  Learning the Cost Function and Online Adaptation to Disturbances
%
%  Compile with:   latexmk -pdf report.tex
%             or:  pdflatex report.tex  (x3, for TOC / refs / cleveref)
%
%  ---------------------------------------------------------------------
%  HOW THIS FILE RELATES TO THE REPOSITORY
%
%   * Every figure is named ONCE, in the FIGURE FILES block below, and is
%     looked for in report/figures/ first and artifacts/figures/ second.
%     To use your own plot: drop it in report/figures/ under the same name,
%     or point the macro at your own file name. A macro pointing at nothing
%     renders as a labelled frame, so the document always compiles.
%     To regenerate the shipped figures:
%         cd study/notebooks && for n in nb?_*.py; do python "$n"; done
%
%   * Every physical constant is defined ONCE in the macro block below and
%     is taken from artifacts/common/nb1_selftest.csv, which the physics
%     self-test regenerates and asserts at import time. Do not retype a
%     number into the body: add a macro.
%
%   * Results obtained at X500_SCALE=smoke are typeset with \SM{} and carry
%     a dagger. They are pipeline diagnostics, NOT results, per the
%     repository's own scale discipline. Every other quantitative cell is a
%     measured number written in by study/fill_report.py; no cell is blank.
%
%   * Two logos go in report/logos/ as institution.png (left) and
%     school.png (right). Absent, they render as labelled frames.
%=======================================================================
\documentclass[11pt,a4paper,twoside,openright]{report}

%----------------------------------------------------------- encoding
\usepackage[T1]{fontenc}
\usepackage[utf8]{inputenc}
\IfFileExists{french.ldf}{%
	\usepackage[french,english]{babel}%
	\newenvironment{frenchpart}{\begin{otherlanguage}{french}}{\end{otherlanguage}}%
}{%
	\usepackage[english]{babel}%
	\newenvironment{frenchpart}{}{}%
}
\IfFileExists{lmodern.sty}
{\usepackage{lmodern}\usepackage{microtype}}
{\usepackage[expansion=false]{microtype}}

%----------------------------------------------------------- layout
\usepackage[a4paper,inner=28mm,outer=22mm,top=25mm,bottom=28mm,
headheight=22pt]{geometry}
\usepackage{fancyhdr}
\usepackage{emptypage}
\usepackage{setspace}
\onehalfspacing

%----------------------------------------------------------- maths
\usepackage{amsmath,amssymb,amsthm,mathtools}
\usepackage{bm}
\usepackage{siunitx}
\sisetup{detect-all, per-mode=symbol, range-phrase={--}, range-units=single}

%----------------------------------------------------------- floats
\usepackage{graphicx}
%  The graphics search path is set ONCE, in the FIGURE FILES block below.
\usepackage{booktabs}
\usepackage{tabularx}
\usepackage{longtable}
\usepackage{multirow}
\usepackage{array}
\usepackage[font=small,labelfont=bf,margin=6pt]{caption}
\usepackage{subcaption}
\usepackage{float}
\usepackage{rotating}
\usepackage[section]{placeins}
\makeatletter
\AtBeginDocument{\let\old@subsection\subsection
	\renewcommand{\subsection}{\FloatBarrier\old@subsection}}
\makeatother

%----------------------------------------------------------- misc
\usepackage{xcolor}
\usepackage{enumitem}
\usepackage{algorithm}
\usepackage{algpseudocode}
\usepackage{listings}
\usepackage{tikz}
\usetikzlibrary{arrows.meta,positioning,calc,fit,backgrounds,shapes.geometric}
\usepackage[framemethod=TikZ]{mdframed}

%=======================================================================
%  COLOURS  (declared before hyperref, which references NavyDeep)
%=======================================================================
\definecolor{NavyDeep}{HTML}{16324F}
\definecolor{AccentBlue}{HTML}{1C7ED6}
\definecolor{AccentOrange}{HTML}{E8590C}
\definecolor{RuleGrey}{HTML}{B8C0C8}
\definecolor{PanelGrey}{HTML}{F4F6F8}
\definecolor{TextGrey}{HTML}{5A6672}

%----------------------------------------------------------- links
\usepackage[breaklinks]{hyperref}
\hypersetup{
	colorlinks=true,
	linkcolor=NavyDeep,
	citecolor=NavyDeep,
	urlcolor=NavyDeep,
	pdftitle={Quadrotor Control by NMPC and Reinforcement Learning},
	pdfsubject={Learning the Cost Function and Online Adaptation to Disturbances},
}
\usepackage[capitalise,nameinlink]{cleveref}

%=======================================================================
%  METADATA
%=======================================================================
\newcommand{\AuthorName}{BOUZID Amine}
\newcommand{\ReportTitle}{Quadrotor Control by NMPC and Reinforcement Learning}
\newcommand{\ReportSubtitle}{Learning the Cost Function and Online Adaptation to Disturbances}
\newcommand{\SchoolName}{EMINES School of Industrial Management --- UM6P}
\newcommand{\OptionName}{Data Science Industrielle}
\newcommand{\HostOrg}{The UM6P Vanguard Center --- Rabat}
\newcommand{\HostTeam}{Autonomous Systems Cluster}
\newcommand{\AcademicSupervisor}{Abdoulaziz LAOUALI}
\newcommand{\DefenceDate}{24 September 2026}
\newcommand{\InternshipPeriod}{April 2026 -- September 2026}

%=======================================================================
%  SINGLE-SOURCED CONSTANTS
%  Every one is asserted by study/tests/ and regenerated into
%  artifacts/common/nb1_selftest.csv. Change them there, not here.
%=======================================================================
\newcommand{\vMass}{2.064308}            % kg, composite
\newcommand{\vFmax}{8.5486}              % N, one rotor
\newcommand{\vTmax}{34.194}              % N, collective
\newcommand{\vTW}{1.6891}                % thrust-to-weight
\newcommand{\vUhover}{0.728742}          % normalised collective at hover
\newcommand{\vTidle}{0.7694}             % N, idle thrust floor
\newcommand{\vOmMin}{150}                % rad/s, PX4 idle rotor speed
\newcommand{\vOmMax}{1000}               % rad/s
\newcommand{\vAlat}{13.350}              % m/s^2, lateral budget
\newcommand{\vAlatCap}{8.0098}           % m/s^2, alpha * budget
\newcommand{\vDadc}{21.6670}             % 1/s^2, d a_z / d c at hover
\newcommand{\vCondPlant}{62.5}           % cond(M), plant allocation
\newcommand{\vCondCtrl}{20.0}            % cond(M_ctrl), controller belief
\newcommand{\vKmCtrl}{0.05}              % m, allocator moment constant
\newcommand{\vJxx}{0.023832}
\newcommand{\vJyy}{0.023935}
\newcommand{\vJzz}{0.044000}
\newcommand{\vH}{64}                     % RDP window length
\newcommand{\vHsec}{1.28}                % s
\newcommand{\vFrame}{26}                 % RDP frame width
\newcommand{\vSmooth}{32}                % RDP rolling-mean length
\newcommand{\vTether}{0.174}             % m, slung tether
\newcommand{\vTetherT}{0.8369}           % s, pendulum period
\newcommand{\vSeedSpread}{0.1658}        % m, nb3 seed spread
\newcommand{\vQRratio}{7.65}             % best-to-worst on the QR grid

%=======================================================================
%  SEMANTIC MACROS
%=======================================================================
\newcommand{\BL}{\textcolor{black!32}{\rule[0.32ex]{1.35em}{0.6pt}}}
\newcommand{\BLs}{\textcolor{black!32}{\rule[0.32ex]{0.9em}{0.6pt}}}
\newcommand{\todo}[1]{\textcolor{AccentOrange}{\small\sffamily[TODO: #1]}}
% A value measured at X500_SCALE=smoke: a pipeline diagnostic, not a result.
\newcommand{\SM}[1]{\textcolor{TextGrey}{#1\,$^{\dagger}$}}

\newcommand{\R}{\mathbb{R}}
\newcommand{\E}{\mathbb{E}}
\newcommand{\Prob}{\mathbb{P}}
\newcommand{\SO}{\mathrm{SO}}
\newcommand{\vecop}{\operatorname{vec}}
\newcommand{\diag}{\operatorname{diag}}
\newcommand{\blkdiag}{\operatorname{blkdiag}}
\newcommand{\clip}{\operatorname{clip}}
\newcommand{\sgn}{\operatorname{sgn}}
\newcommand{\Tr}{\operatorname{Tr}}
\newcommand{\st}{\quad\text{s.t.}\quad}
\newcommand{\defeq}{\vcentcolon=}
\newcommand{\transp}{^{\!\top}}
\renewcommand{\vec}[1]{\bm{#1}}
\newcommand{\mat}[1]{\bm{#1}}
\newcommand{\ctrlstate}{\vec{x}}
\newcommand{\plantstate}{\vec{s}}
\newcommand{\inp}{\vec{u}}
\newcommand{\err}{\vec{e}}
\newcommand{\dres}{\vec{d}}
\newcommand{\wrench}{\vec{w}}
\newcommand{\quat}{\vec{q}}
\newcommand{\bodyrate}{\vec{\omega}}
\newcommand{\refsym}{\mathrm{ref}}

\DeclareRobustCommand{\ACMPC}{AC\nbd MPC}
\DeclareRobustCommand{\LLTC}{LLTC}
\DeclareRobustCommand{\RDP}{RDP}
\DeclareRobustCommand{\CTBR}{CTBR}
\DeclareRobustCommand{\NMPC}{NMPC}

%=======================================================================
%  THEOREM ENVIRONMENTS
%=======================================================================
\theoremstyle{plain}
\newtheorem{theorem}{Theorem}[chapter]
\newtheorem{lemma}[theorem]{Lemma}
\newtheorem{proposition}[theorem]{Proposition}
\newtheorem{corollary}[theorem]{Corollary}
\theoremstyle{definition}
\newtheorem{definition}[theorem]{Definition}
\newtheorem{assumption}[theorem]{Assumption}
\newtheorem{problem}[theorem]{Problem}
\theoremstyle{remark}
\newtheorem{remark}[theorem]{Remark}

%=======================================================================
%  KEY-FINDING / CAUTION BOXES
%=======================================================================
\newmdenv[
linewidth=0.9pt, linecolor=RuleGrey, backgroundcolor=PanelGrey,
roundcorner=2pt, innertopmargin=8pt, innerbottommargin=8pt,
innerleftmargin=10pt, innerrightmargin=10pt, skipabove=9pt, skipbelow=9pt
]{keybox}

\newmdenv[
linewidth=0pt, leftline=true, linewidth=2.4pt, linecolor=AccentOrange,
backgroundcolor=white, topline=false, bottomline=false, rightline=false,
innertopmargin=4pt, innerbottommargin=4pt, innerleftmargin=10pt,
skipabove=9pt, skipbelow=9pt
]{cautionbox}

% amsmath's \nobreakdash is fragile: a caption carrying it is written to the
% .toc/.lof/.lot through \protected@edef, where it expands into garbage and
% kills the second pass.  This is the same thing with an unbreakable hyphen.
\DeclareRobustCommand{\nbd}{\mbox{-}}

\newcommand{\keyhead}[1]{\textbf{\sffamily\small\textcolor{NavyDeep}{#1}}\par\vspace{2pt}}


%=======================================================================
%  >>>>>  FIGURE FILES -- EDIT THIS BLOCK TO USE YOUR OWN PLOTS  <<<<<
%
%  Every figure in the document is referred to through ONE macro defined
%  here.  To swap a plot for your own, change the file name on the right
%  of the matching line -- nothing else in the document needs touching.
%
%      \renewcommand{\figHeadline}{my_headline_plot.pdf}
%
%  WHERE THE FILE IS LOOKED FOR, in this order (see \graphicspath below):
%
%      report/figures/          <- put YOUR OWN images here
%      artifacts/figures/       <- what the notebooks generate
%      the directory of report.tex, and the compile directory
%
%  So the simplest workflow in TeXstudio is: drop your image into
%  report/figures/ and either name it as below, or point the macro at it.
%  PDF, PNG, JPG and EPS all work; PDF gives the best result for plots.
%
%  A name that resolves to no file on disk renders as a labelled frame
%  rather than breaking the build, so the document always compiles and a
%  missing plot is visible in the PDF.
%=======================================================================

% -- Chapter 3: model, references, feasibility -------------------------
\newcommand{\figFeasibility}     {F1_feasibility.png}
\newcommand{\figSaturationGuard} {F24_saturation_guard.png}

% -- Chapter 5: implementation, flight visualisation -------------------
\newcommand{\figContactSheet}    {F25_contact_sheet.png}
\newcommand{\figGroundTracks}    {F26_ground_tracks.png}
\newcommand{\figWrenchSeries}    {F19_wrench_timeseries.png}

% -- Chapter 6, Increment 0: baselines ---------------------------------
\newcommand{\figLqrFeedforward}  {F2_lqr_feedforward.png}
\newcommand{\figHorizon}         {F3_horizon.png}
\newcommand{\figCorner}          {F4_corner.png}
\newcommand{\figNoise}           {F3b_noise.png}
\newcommand{\figQrHeatmap}       {F5_qr_heatmap.png}
\newcommand{\figMismatchHand}    {F27_model_mismatch_hand.png}
\newcommand{\figMismatchCost}    {F28_cost_mismatch_hand.png}

% -- Chapter 6, Increment 1: learned terminal cost ---------------------
\newcommand{\figLltcFit}         {F6_lltc_fit.png}
\newcommand{\figLltcSpectrum}    {F8_lltc_spectrum.png}
\newcommand{\figLltcLocality}    {F7_lltc_locality.png}

% -- Chapter 6, Increment 2: learned cost map --------------------------
\newcommand{\figRepresentations} {F9_representations.png}
\newcommand{\figExploration}     {F10_exploration.png}
\newcommand{\figMismatchLearned} {F29_model_mismatch_learned.png}

% -- Chapter 6, Increment 3: domain randomisation ----------------------
\newcommand{\figAxisCompetence}  {F13_axis_competence.png}
\newcommand{\figVariances}       {F14_variances.png}
\newcommand{\figCurriculum}      {F15_curriculum.png}

% -- Chapter 6, Increment 4: adaptive scheme ---------------------------
\newcommand{\figTrainingGate}    {F20_training_gate.png}
\newcommand{\figEncoderTrade}    {F16_encoder_trade.png}
\newcommand{\figChannelRtwo}     {F17_channel_r2.png}
\newcommand{\figRdpWindow}       {F30_rdp_window.png}
\newcommand{\figClosedLoop}      {F18_closed_loop.png}

% -- Chapter 6: the comparison, error analysis -------------------------
\newcommand{\figHeadline}        {F21_headline.png}
\newcommand{\figRadar}           {F23_radar.png}
\newcommand{\figByCondition}     {F22_by_condition.png}
\newcommand{\figFragility}       {F31_fragility.png}
\newcommand{\figInterpretability}{F12_interpretability.png}

%  Search path.  report/figures/ comes FIRST so a file you drop there
%  overrides the generated one of the same name without editing anything.
\graphicspath{%
	{figures/}{report/figures/}{../figures/}%
	{artifacts/figures/}{../artifacts/figures/}{report/../artifacts/figures/}%
	{logos/}{report/logos/}%
	{./}}

%=======================================================================
%  FIGURE INCLUSION
%  \realfig{height}{file}{caption}{label}
%  Includes artifacts/figures/<file> when present; otherwise draws a frame
%  naming the artefact and the notebook that produces it, so a missing
%  figure is visible in the PDF instead of breaking the build.
%=======================================================================
\newcommand{\missingfig}[2]{%
	% #1 = reserved height, #2 = file name.  The name is typeset with
	% \detokenize so that an underscore in it stays a printable character
	% instead of opening math mode.
	\begingroup
	\setlength{\fboxsep}{6pt}%
	\setlength{\fboxrule}{0.8pt}%
	\fcolorbox{RuleGrey}{PanelGrey}{%
		\begin{minipage}[c][#1][c]{0.84\linewidth}%
			\centering
			{\sffamily\small\textcolor{NavyDeep}{\textbf{FIGURE NOT BUILT}}}\par
			\vspace{3pt}
			{\ttfamily\footnotesize\textcolor{TextGrey}{\detokenize{artifacts/figures/}\detokenize{#2}}}\par
			\vspace{5pt}
			{\rmfamily\footnotesize\textcolor{TextGrey}{Run the notebooks to generate it.}}%
		\end{minipage}}%
	\endgroup}

%  Resolve a figure name against every directory the build might be run
%  from.  \IfFileExists is relative to the *compile* directory, not to the
%  directory holding this file, so `latexmk report/report.tex' from the
%  repository root and `latexmk report.tex' from report/ need different
%  prefixes.  \figfound is empty when none of them hit.
\makeatletter
\newcommand{\fig@found}{}
\newcommand{\fig@locate}[1]{%
	\def\fig@found{}%
	\IfFileExists{../artifacts/figures/#1}{\def\fig@found{#1}}{%
		\IfFileExists{artifacts/figures/#1}{\def\fig@found{#1}}{%
			\IfFileExists{report/../artifacts/figures/#1}{\def\fig@found{#1}}{%
				\IfFileExists{figures/#1}{\def\fig@found{#1}}{%
					\IfFileExists{#1}{\def\fig@found{#1}}{}}}}}%
}
%  \figslot{width}{height}{file} -- the picture, or the fallback frame.
\newcommand{\figslot}[3]{%
	\fig@locate{#3}%
	\ifx\fig@found\@empty
	\missingfig{#2}{#3}%
	\else
	\includegraphics[width=#1,height=#2,keepaspectratio]{\fig@found}%
	\fi}
\makeatother

\makeatletter
\newcommand{\logo@found}{}
\newcommand{\logoslot}[2]{%
	% #1 = file inside report/logos/, #2 = placeholder caption.
	\def\logo@found{}%
	\IfFileExists{logos/#1}{\def\logo@found{#1}}{%
		\IfFileExists{report/logos/#1}{\def\logo@found{#1}}{}}%
	\ifx\logo@found\@empty
	\begin{tikzpicture}
		\node[draw=RuleGrey, fill=PanelGrey, line width=0.7pt,
		rounded corners=2pt, minimum width=34mm, minimum height=18mm,
		align=center, font=\sffamily\scriptsize, text=TextGrey]
		{#2\\[1pt]\ttfamily\detokenize{logos/}\detokenize{#1}};
	\end{tikzpicture}%
	\else
	\includegraphics[height=20mm,keepaspectratio]{\logo@found}%
	\fi}
\makeatother

\newcommand{\realfig}[4]{%
	\begin{figure}[htbp]\centering
		\figslot{0.94\linewidth}{#1}{#2}
		\caption{#3}
		\label{#4}
\end{figure}}

% two real figures side by side
\newcommand{\realfigtwo}[7]{%
	\begin{figure}[htbp]\centering
		\begin{subfigure}[b]{0.48\linewidth}\centering
			\figslot{\linewidth}{#1}{#2}
			\caption{#5}
		\end{subfigure}\hfill
		\begin{subfigure}[b]{0.48\linewidth}\centering
			\figslot{\linewidth}{#1}{#3}
			\caption{#6}
		\end{subfigure}
		\caption{#4}
		\label{#7}
\end{figure}}

%=======================================================================
%  HEADERS
%=======================================================================
\pagestyle{fancy}
\fancyhf{}
\renewcommand{\headrulewidth}{0.4pt}
\renewcommand{\headrule}{\hbox to\headwidth{\color{RuleGrey}\leaders\hrule height \headrulewidth\hfill}}
\fancyhead[LE]{\footnotesize\sffamily\textcolor{TextGrey}{\leftmark}}
\fancyhead[RO]{\footnotesize\sffamily\textcolor{TextGrey}{\rightmark}}
\fancyfoot[LE,RO]{\footnotesize\sffamily\thepage}
\renewcommand{\chaptermark}[1]{\markboth{\chaptername\ \thechapter.\ #1}{}}
\renewcommand{\sectionmark}[1]{\markright{\thesection\ \ #1}}
\fancypagestyle{plain}{\fancyhf{}\renewcommand{\headrulewidth}{0pt}
	\fancyfoot[LE,RO]{\footnotesize\sffamily\thepage}}

%----------------------------------------------------------- listings
\lstset{
	basicstyle=\ttfamily\footnotesize,
	keywordstyle=\color{NavyDeep}\bfseries,
	commentstyle=\color{TextGrey}\itshape,
	stringstyle=\color{AccentBlue},
	numbers=left, numberstyle=\tiny\color{TextGrey}, numbersep=7pt,
	frame=single, rulecolor=\color{RuleGrey}, framesep=5pt,
	backgroundcolor=\color{PanelGrey},
	breaklines=true, breakatwhitespace=true, showstringspaces=false,
	captionpos=b, xleftmargin=14pt, language=Python
}

\allowdisplaybreaks
\setlength{\parskip}{2pt}

%=======================================================================
\begin{document}
	%=======================================================================

	%----------------------------------------------------------- TITLE PAGE
	\begin{titlepage}
		\centering
		\thispagestyle{empty}

		\begin{tikzpicture}[remember picture, overlay]
			\draw[line width=1.1pt, color=NavyDeep]
			($(current page.north west)+(20mm,-18mm)$) --
			($(current page.north east)+(-18mm,-18mm)$);
			\draw[line width=1.1pt, color=NavyDeep]
			($(current page.south west)+(20mm,20mm)$) --
			($(current page.south east)+(-18mm,20mm)$);
		\end{tikzpicture}

		\vspace*{6mm}

		%------------------------------------------------- LOGOS
		% Left: host institution. Right: school.
		% Drop institution.png and school.png into report/logos/.
		% Either may also be .pdf; change the extension below if so.
		\noindent
		\begin{minipage}[c]{0.34\linewidth}\raggedright
			\logoslot{institution.jpeg}{INSTITUTION LOGO}
		\end{minipage}
		\hfill
		\begin{minipage}[c]{0.28\linewidth}\centering
		\end{minipage}
		\hfill
		\begin{minipage}[c]{0.34\linewidth}\raggedleft
			\logoslot{school.jpeg}{SCHOOL LOGO}
		\end{minipage}

		\vspace{12mm}

		{\sffamily\large\textcolor{TextGrey}{\SchoolName}}\\[2mm]
		{\sffamily\normalsize\textcolor{TextGrey}{\OptionName}}\\[3mm]
		{\sffamily\footnotesize\textcolor{TextGrey}{Master of Engineering ---
				Final Year Project (Projet de Fin d'\'Etudes)}}

		\vspace{18mm}

		{\color{NavyDeep}\rule{0.30\linewidth}{1.4pt}}

		\vspace{9mm}
		{\sffamily\fontsize{24}{29}\selectfont\bfseries\textcolor{NavyDeep}{\ReportTitle}}

		\vspace{6mm}
		{\sffamily\fontsize{13}{17}\selectfont\textcolor{TextGrey}{\ReportSubtitle}}

		\vspace{9mm}
		{\color{NavyDeep}\rule{0.30\linewidth}{1.4pt}}

		\vspace{16mm}

		{\large\textit{\AuthorName}}

		\vspace{14mm}

		\begin{minipage}{0.86\linewidth}\centering
			\small
			\renewcommand{\arraystretch}{1.35}
			\begin{tabular}{r@{\hspace{7mm}}l}
				\textcolor{TextGrey}{Host organisation} & \textit{\HostOrg} \\
				\textcolor{TextGrey}{Team} & \textit{\HostTeam} \\
				\textcolor{TextGrey}{Academic supervisor} & \textit{\AcademicSupervisor} \\
				\textcolor{TextGrey}{Period} & \textit{\InternshipPeriod} \\
			\end{tabular}
		\end{minipage}

		\vfill
		{\sffamily\small\textcolor{TextGrey}{Defended on \DefenceDate}}
		\vspace{14mm}
	\end{titlepage}

	\cleardoublepage

	%----------------------------------------------------------- FRONT MATTER
	\pagenumbering{roman}
	\setcounter{page}{1}

	%============================== ABSTRACT (EN) ==========================
	\chapter*{Abstract}
	\addcontentsline{toc}{chapter}{Abstract / R\'esum\'e}
	\markboth{Abstract}{Abstract}

	\noindent
	Nonlinear model predictive control (NMPC) gives a quadrotor an explicit,
	inspectable optimisation problem to solve at every control period, and a mature
	theory with which to reason about constraints and stability. Its accuracy,
	however, is bounded by the fidelity of the model it predicts with, and by a
	stage cost that a human must choose by hand. Reinforcement learning (RL) removes
	the hand-tuning but replaces the optimisation problem by an opaque policy that is
	difficult to certify. This work develops and evaluates a hybrid in which an NMPC
	scheme is retained as the \emph{structure} of the policy and RL is used to learn
	its \emph{cost}, following the result of Gros and Zanon that a predictive
	controller built on an inexact model can still realise the optimal policy of the
	true system once its stage and terminal costs are adapted.

	Four increments are studied on a Holybro X500 quadrotor, simulated by a
	\num{23}-state rigid-body plant containing the inner rate loop, the rotor mixer
	and the PX4 idle-thrust floor, and deployed in a PX4/Gazebo ROS\,2 workspace:
	(i)~a learned quadratic terminal cost intended to give a one-step horizon the
	behaviour of a long one; (ii)~a fully learned cost map over a differentiable
	one-step MPC, trained by PPO, which removes every hand-chosen stage weight;
	(iii)~domain randomisation, and a measurement of what it costs; and (iv)~an
	adaptive scheme in which a recurrent residual-dynamics predictor reconstructs
	the \num{6}-dimensional external wrench from onboard signals and feeds it to the
	model inside the optimiser. Controllers are compared against LQR and classical
	NMPC on a common action space, a common cost, and a common actuator envelope,
	under nominal, in-distribution and out-of-distribution disturbance suites.

	The report also documents a systematic audit of the experimental pipeline
	conducted against the two source papers and against the manufacturer's own
	vehicle description. It identified defects that invalidated earlier conclusions,
	among them a dimensional inconsistency in the disturbance channel supplied to
	the predictive model, an actuator map that ignored the autopilot's idle floor,
	and exploration noise that never reached the action it was supposed to perturb.
	Negative and non-reproducing results are reported as such.

	\vspace{4mm}
	\noindent\textbf{Keywords:} nonlinear model predictive control; reinforcement
	learning; differentiable optimisation; iLQR; quadrotor; residual dynamics
	learning; domain randomisation; PX4 Autopilot; JAX.

	%============================== R\'ESUM\'E (FR) ========================
	\begin{frenchpart}
		\chapter*{R\'esum\'e}
		\markboth{R\'esum\'e}{R\'esum\'e}

		\noindent
		La commande pr\'edictive non lin\'eaire (NMPC) dote un quadrirotor d'un probl\`eme
		d'optimisation explicite, r\'esolu \`a chaque p\'eriode d'\'echantillonnage, et
		d'une th\'eorie mature permettant de raisonner sur les contraintes et la
		stabilit\'e. Sa pr\'ecision reste toutefois born\'ee par la fid\'elit\'e du
		mod\`ele de pr\'ediction et par un co\^ut d'\'etape que l'ing\'enieur doit
		r\'egler \`a la main. L'apprentissage par renforcement (RL) supprime ce
		r\'eglage manuel, mais remplace le probl\`eme d'optimisation par une politique
		opaque, difficile \`a certifier. Ce travail d\'eveloppe et \'evalue une approche
		hybride~: le sch\'ema NMPC est conserv\'e comme \emph{structure} de la
		politique, et le RL sert \`a en apprendre le \emph{co\^ut}, en s'appuyant sur le
		r\'esultat de Gros et Zanon selon lequel un contr\^oleur pr\'edictif fond\'e sur
		un mod\`ele inexact peut n\'eanmoins r\'ealiser la politique optimale du
		syst\`eme r\'eel d\`es lors que ses co\^uts d'\'etape et terminal sont adapt\'es.

		Quatre incr\'ements sont \'etudi\'es sur un quadrirotor Holybro X500, simul\'e
		par un mod\`ele rigide \`a \num{23} \'etats incluant la boucle interne de
		vitesse angulaire, le m\'elangeur et le plancher de pouss\'ee au ralenti de
		PX4, puis d\'eploy\'e dans un espace de travail ROS\,2 PX4/Gazebo~: (i)~un
		co\^ut terminal quadratique appris, destin\'e \`a conf\'erer \`a un horizon
		d'un pas le comportement d'un horizon long~; (ii)~une carte de co\^ut
		enti\`erement apprise au-dessus d'une MPC diff\'erentiable \`a un pas,
		entra\^in\'ee par PPO, qui supprime tout poids r\'egl\'e \`a la main~;
		(iii)~la randomisation de domaine, et la mesure de son co\^ut~; (iv)~un
		sch\'ema adaptatif dans lequel un pr\'edicteur r\'ecurrent de dynamique
		r\'esiduelle reconstruit le torseur ext\'erieur \`a \num{6} composantes \`a
		partir des seuls signaux embarqu\'es et l'injecte dans le mod\`ele interne de
		l'optimiseur. Les contr\^oleurs sont compar\'es \`a un LQR et \`a une NMPC
		classique, sur le m\^eme espace d'action, le m\^eme co\^ut et la m\^eme
		enveloppe d'actionneurs, sous des perturbations nominales, en distribution et
		hors distribution.

		Le rapport documente enfin un audit syst\'ematique de la cha\^ine
		exp\'erimentale, men\'e face aux deux articles sources et \`a la description
		constructeur du v\'ehicule. Il a mis en \'evidence plusieurs d\'efauts
		invalidant des conclusions ant\'erieures, dont une incoh\'erence
		dimensionnelle du canal de perturbation fourni au mod\`ele pr\'edictif, une
		carte d'actionneurs ignorant le plancher de ralenti de l'autopilote, et un
		bruit d'exploration n'atteignant jamais l'action qu'il devait perturber. Les
		r\'esultats n\'egatifs ou non reproduits sont rapport\'es comme tels.

		\vspace{4mm}
		\noindent\textbf{Mots-cl\'es~:} commande pr\'edictive non lin\'eaire~;
		apprentissage par renforcement~; optimisation diff\'erentiable~; iLQR~;
		quadrirotor~; apprentissage de dynamique r\'esiduelle~; randomisation de
		domaine~; PX4~; JAX.
	\end{frenchpart}

	\cleardoublepage

	%----------------------------------------------------------- ACKNOWLEDGEMENTS
	\chapter*{Acknowledgements}
	\addcontentsline{toc}{chapter}{Acknowledgements}
	\markboth{Acknowledgements}{Acknowledgements}
	I thank \HostOrg{} for hosting this final-year project, and the
	\HostTeam{} for the working conditions that made it possible --- in
	particular the access to high-end GPU workstations without which the
	training campaigns in \cref{ch:results} could not have been run at all, and
	the flight and marine platforms that kept the questions in this report
	anchored to hardware rather than to simulation alone.

	I am grateful to my academic supervisor, \AcademicSupervisor, for steady
	guidance throughout, and for pressing on the questions that were easiest to
	leave unasked. Several of the corrections recorded in \cref{app:audit}
	began as a supervisor's question about a number that did not look right.

	I thank the members of the team for the technical discussions that shaped
	this work, for their patience with a project that changed direction more
	than once, and for the reviews that caught in conversation what no test
	would have caught in code.

	Finally, I thank \SchoolName{} and the teaching staff of the
	\OptionName{} option for the training that preceded this project, and my
	family and friends for their support over its six months.

	\cleardoublepage

	%----------------------------------------------------------- CONTENTS
	\tableofcontents
	\cleardoublepage
	\listoffigures
	\cleardoublepage
	\listoftables
	\cleardoublepage

	%----------------------------------------------------------- NOMENCLATURE
	\chapter*{Nomenclature}
	\addcontentsline{toc}{chapter}{Nomenclature}
	\markboth{Nomenclature}{Nomenclature}

	\vspace{-3mm}
	\subsection*{Acronyms}
	\begin{longtable}{@{}p{26mm}p{112mm}@{}}
		\toprule
		\ACMPC  & Actor--critic model predictive control \\
		CTBR    & Collective thrust and body rates (the action space) \\
		DR      & Domain randomisation \\
		ENU/FLU & East--North--Up (world) / Forward--Left--Up (body) frames \\
		GAE     & Generalised advantage estimation \\
		iLQR    & Iterative linear--quadratic regulator \\
		LLTC    & Learned local time-varying terminal cost \\
		LQR     & Linear--quadratic regulator \\
		MDP     & Markov decision process \\
		MPC     & Model predictive control \\
		NED/FRD & North--East--Down / Forward--Right--Down frames (PX4 convention) \\
		NMPC    & Nonlinear model predictive control \\
		OOD     & Out of distribution \\
		PPO     & Proximal policy optimisation \\
		PWM     & Normalised motor command \\
		RDP     & Residual dynamics predictor \\
		RL      & Reinforcement learning \\
		RMSE    & Root-mean-square error \\
		SITL    & Software in the loop \\
		\bottomrule
	\end{longtable}

	\vspace{-4mm}
	\subsection*{Principal symbols}
	\begin{longtable}{@{}p{26mm}p{88mm}p{22mm}@{}}
		\toprule
		\textbf{Symbol} & \textbf{Meaning} & \textbf{Units} \\
		\midrule
		\endhead
		$\ctrlstate\in\R^{17}$ & Control-model state $[\vec p,\vec v,\quat,\bodyrate,\vec\Omega]$ & --- \\
		$\plantstate\in\R^{23}$ & Plant state (rigid body $+$ actuators $+$ filters) & --- \\
		$\inp\in\R^{4}$ & CTBR command $[c,\bar\omega_x,\bar\omega_y,\bar\omega_z]$ & --- \\
		$\err\in\R^{16}$ & Tracking error in the reduced attitude parametrisation & --- \\
		$\dres\in\R^{6}$ & Model residual $[\vec a_\mathrm{res},\vec\alpha_\mathrm{res}]$ & \si{\metre\per\square\second}, \si{\radian\per\square\second} \\
		$\wrench\in\R^{6}$ & External wrench $[\vec F_\mathrm{ext},\vec\tau_\mathrm{ext}]$ & \si{\newton}, \si{\newton\metre} \\
		$m$ & Composite vehicle mass & \si{\kilogram} \\
		$T_{\max}$ & Maximum collective thrust & \si{\newton} \\
		$\Omega_{\min},\Omega_{\max}$ & Idle and maximum rotor speed & \si{\radian\per\second} \\
		$\Omega_\mathrm{sc}$ & Rotor-speed scale in the error coordinates & \si{\radian\per\second} \\
		$\mat J$ & Composite inertia tensor about the CG & \si{\kilogram\square\metre} \\
		$\mat R(\quat)$ & Body-to-world rotation matrix & --- \\
		$\mat S_k,\vec c_k$ & Learned quadratic / linear cost terms at stage $k$ & --- \\
		$\mat P_\mathrm{term}$ & Terminal cost matrix & --- \\
		$N$ & Prediction horizon (stages) & --- \\
		$H$ & Residual-predictor window, $\vH$ frames & --- \\
		$\Delta t_c$ & Outer-loop control period, \num{0.02} & \si{\second} \\
		$\gamma,\lambda$ & Discount factor, GAE trace decay & --- \\
		$\theta$ & Learned parameters of the cost map / predictor & --- \\
		\bottomrule
	\end{longtable}

	\cleardoublepage
	\pagenumbering{arabic}
	\setcounter{page}{1}
	%=======================================================================
	\chapter{Introduction}
	\label{ch:intro}
	%=======================================================================

	\section{General context}

	An autonomous multirotor is a small, fast, weakly damped system with four
	actuators, six degrees of freedom and a thrust-to-weight ratio rarely above two.
	It must reject wind, tolerate a payload whose mass and position are not known in
	advance, and do all of this inside a control period of a few tens of
	milliseconds. Two families of methods dominate the literature, and they fail in
	opposite ways.

	\textbf{Model predictive control} formulates an explicit finite-horizon optimal
	control problem at every sampling instant, applies the first element of its
	solution, and repeats. Its appeal in aerial robotics is not only performance but
	\emph{legibility}: constraints appear as constraints, the cost is written down,
	and a mature body of theory addresses recursive feasibility and closed-loop
	stability~\cite{rawlings2017mpc,mayne2000constrained}. Its weakness is equally
	explicit. The quality of the resulting policy is bounded by the fidelity of the
	prediction model, and by a stage cost --- a set of weights $Q$ and $R$ --- that
	somebody must choose. Neither bound is addressed by the theory: it assumes the
	model and takes the cost as given.

	\textbf{Reinforcement learning} removes the hand-tuning. A policy is optimised
	directly against a task-level reward, and recent results in aggressive flight
	demonstrate that the resulting controllers can exceed carefully engineered
	baselines~\cite{hwangbo2017quadrl,kaufmann2023champion}. The price is that the
	policy is a neural network. There is no optimisation problem to inspect, no
	place where a constraint is written, and the certification story is an open
	research area~\cite{garcia2015safe}.

	The question this project addresses is whether the two can be combined in a way
	that keeps what is valuable in each: \emph{retain the predictive controller as
		the structure of the policy, and learn its cost}.

	\section{The scientific premise}

	That combination is not merely a heuristic. Gros and Zanon~\cite{gros2020data}
	prove that an economic NMPC scheme built on a \emph{wrong} model can nevertheless
	deliver the optimal policy of the true system, provided its stage cost, terminal
	cost and constraints are adapted. Informally, model error can be absorbed into
	the cost. The practical consequence they draw is the one this project exploits:
	an NMPC scheme may be used as a \emph{function approximator} inside RL, taking
	the place normally occupied by a deep network, and all the usual RL machinery ---
	temporal-difference learning, policy gradients --- can be used to tune it.

	This inverts the usual framing of ``learning-based MPC''. The model is not what
	is learned. The cost is.

	\section{Problem statement}

	\begin{problem}[Informal]
		Given a quadrotor whose dynamic model is only approximately known and whose
		operating conditions include unmeasured wind, mass and inertia variation, rotor
		degradation and payloads attached at unknown positions, design a controller
		which
		\begin{enumerate}[nosep,label=(\roman*)]
			\item retains an explicit, short-horizon optimal control problem in the loop,
			so that constraints and actuator limits are represented structurally;
			\item requires no hand-chosen stage-cost weights;
			\item adapts online to the disturbance actually acting on the vehicle; and
			\item executes within a \SI{20}{\milli\second} control period on a
			single-vehicle deployment.
		\end{enumerate}
	\end{problem}

	\section{Objectives and contributions}

	The work is organised as four increments, each of which isolates one claim and
	tests it against a baseline that differs in exactly one respect.

	\begin{enumerate}[leftmargin=*]
		\item \textbf{Establish the problem and the cost of an exact solution.}
		Quantify how tracking accuracy and computation scale with the prediction
		horizon; verify that the reference trajectories are inside the vehicle's
		actuator envelope; and measure the sensitivity of the closed loop to the
		choice of $Q$ and $R$. That last measurement --- the best-to-worst spread over
		the admissible grid (\cref{sec:res-baselines}) --- is the motivation for
		everything that follows.
		\item \textbf{Learn a terminal cost (\LLTC).} A quadratic terminal cost
		$\mat P_\theta(\err_0)=\mat L\mat L\transp+\varepsilon\mat I$ is regressed
		against the realised cost-to-go, with the intention of giving a one-step
		horizon the closed-loop behaviour of a long one --- the classical
		observation that a finite-horizon problem recovers the infinite-horizon
		optimum when the terminal cost is the true value
		function~\cite{lowrey2019polo,bhardwaj2021blending}.
		\item \textbf{Learn the whole cost map (\ACMPC).} Following Romero et
		al.~\cite{romero2024acmpc}, a network maps the observation to the full
		stage cost $(\mat S_k,\vec c_k)$ of a differentiable one-step MPC, trained
		end to end by PPO through the solver. Every hand-chosen stage weight is
		removed.
		\item \textbf{Adapt to the disturbance (\RDP).} A recurrent predictor
		reconstructs the \num{6}-dimensional external wrench from a rolling window
		of states and \emph{actuator commands}, and the estimate is injected into
		the model inside the optimiser. This is compared against domain
		randomisation, which buys robustness without the ability to specialise.
	\end{enumerate}

	Alongside these, two contributions are methodological rather than algorithmic,
	and are reported because they materially changed the conclusions:

	\begin{enumerate}[leftmargin=*,resume]
		\item \textbf{A systematic audit of the experimental pipeline}
		(\cref{sec:audit-defects}, \cref{app:audit}), conducted against the two
		source papers and against the manufacturer's own vehicle description, which
		identified defects that silently changed results --- among them an actuator
		map that ignored the autopilot's idle floor, exploration noise that never
		reached the action, and a dimensional error at the disturbance interface.
		\item \textbf{A sim-to-real path} in the form of a PX4/Gazebo ROS\,2 workspace
		that computes the same wrench from onboard signals only, so that the
		predictor trained in simulation can be evaluated against a physics engine
		it has never seen, with a real flight-control stack in the loop.
	\end{enumerate}

	\section{Structure of this report}

	\Cref{ch:context} situates the work. \Cref{ch:sota} reviews the relevant
	literature and develops the mathematical framework: the vehicle model, the error
	coordinates, the differentiable optimisation layer and the RL formulation.
	\Cref{ch:data} describes the data --- which, in a simulation-driven study, means
	the environment, the disturbance distributions and the sampling design.
	\Cref{ch:method} presents the methodology increment by increment, together with
	the results borrowed from the literature that justify each choice.
	\Cref{ch:impl} covers implementation, tooling and the experimental protocol.
	\Cref{ch:results} reports and analyses the results, including the negative ones.
	\Cref{ch:conclusion} concludes.

	\subsection{The argument, in one page}
	\label{sec:argument}

	The increments of \cref{ch:method} are not a menu of things that were tried.
	They are a chain, and each link exists because the previous one failed in a
	specific, measurable way. Stated once here so that every later chapter can be
	read against it:

	\begin{enumerate}[leftmargin=*]
		\item \textbf{The problem.} A quadrotor must track an aggressive reference on
		an embedded computer, under a plant that is never exactly the one it plans
		with. Three things are therefore being asked at once: accuracy, robustness and
		a solve that fits inside \SI{20}{\milli\second}.

		\item \textbf{The classical answer works, and pays for it three times.}
		LQR and NMPC solve the nominal problem well (\cref{sec:baselines}). But the
		online solve is (i) \emph{sensitive to model mismatch} --- it plans several
		steps against a model that is wrong, so the error compounds where a fixed gain
		merely reacts; (ii) \emph{expensive}, with latency growing steeply in the
		horizon $N$; and (iii) \emph{sensitive to its cost weights}, which nothing in
		the theory determines and a human must choose. \Cref{sec:res-sensitivity}
		measures all three.

		\item \textbf{So: shorten the horizon without losing what it bought.} A
		long horizon buys optimality and a margin against mismatch. Shortening it to
		$N{=}1$ removes the latency problem outright; the question is whether the
		lost optimality can be put back into the \emph{cost} rather than the
		\emph{horizon}.

		\item \textbf{Increment~1 --- put it in the terminal cost. Not enough.}
		A learned local terminal cost (\cref{sec:lltc}) is the minimal version of that
		idea: one matrix at the end of a one-step problem. \Cref{sec:res-lltc} reports
		whether it recovers long-horizon behaviour.

		\item \textbf{Increment~2 --- put it in the whole cost map.} AC-MPC
		(\cref{sec:acmpc}) learns the stage cost at every stage, which subsumes the
		terminal cost and, in doing so, removes the hand-chosen weights that
		point~2(iii) identified as a liability.

		\item \textbf{Increment~3 --- generalisation by randomisation, at a price.}
		A cost map learned on one plant is a cost map for that plant. Domain
		randomisation (\cref{sec:dr}) buys robustness across a family, but a policy
		that cannot \emph{observe} which member of the family it is flying can only
		find one compromise acceptable everywhere. Generality is bought with
		specificity.

		\item \textbf{Increment~4 --- give it the missing observation.} If the
		conservatism of Increment~3 is the price of not knowing the plant, then
		estimating the disturbance online (\cref{sec:adaptive}) should recover the
		specificity without giving up the generality.

		\item \textbf{The endpoint: an enhanced cost and an enhanced model.} The two
		learned objects in \cref{fig:architecture} are exactly these. The cost map
		$\theta$ is the enhanced \emph{cost}, carrying what the horizon used to
		provide; the residual predictor $\psi$ is the enhanced \emph{model}, carrying
		what domain randomisation had to average over. Everything between them ---
		the dynamics, the iLQR, the actuator map --- is unchanged, which is what makes
		each contribution attributable.
	\end{enumerate}

	Whether each link holds is an empirical question, and \cref{ch:results} answers
	the ones the campaign has reached. Where it does not hold, the chain is reported
	as broken rather than repaired in prose.

	%=======================================================================
	\chapter{Host organisation and problem context}
	\label{ch:context}
	%=======================================================================

	\section{The host organisation}

	\HostOrg{} is a research centre organised around three clusters: large
	language models, autonomous agents, and robotics. The three are not
	administrative neighbours but a deliberate arrangement: the centre's premise
	is that the learning methods developed in the first two are the ones that
	will eventually be asked to close physical control loops in the third, and
	that the interesting problems live where they meet.

	This project sat in the \HostTeam, whose mission is to build a
	\emph{general control system framework} --- one architecture that can be
	retargeted across vehicles and across media rather than a controller
	rewritten for each platform. That mission is what fixes the shape of this
	report. A framework claim is only worth making if the same structure
	survives a change of plant, which is why \cref{ch:method} is organised as
	four increments each adding one mechanism, why the control model is
	separated from the plant model throughout (\cref{sec:ctrlmodel}), and
	why \cref{ch:impl} keeps the study half and the robotics half in separate
	codebases that meet only through a serialised model file. A result that
	held only for one airframe, or only inside one simulator, would not answer
	the question the team is asking.

	\begin{itemize}[nosep]
		\item \textbf{Organisation:} \HostOrg{} --- clusters for large language
		models, agents, and robotics.
		\item \textbf{Team:} \HostTeam.
		\item \textbf{Mission relevant to this project:} to build a general
		control system framework, retargetable across vehicles rather than
		rebuilt per platform.
		\item \textbf{Existing assets the project built on:} high-end
		CUDA-capable workstations for training; unmanned surface vessel and
		drone platforms for flight and on-water work; and an established PX4
		flight-stack codebase, which is why \cref{ch:impl} targets PX4 and
		Gazebo rather than a generic simulator and why the constants of
		\cref{tab:vehicle} are taken from the airframe description the team
		already flies.
	\end{itemize}

	\section{The operational problem}

	The engineering problem behind the scientific one can be stated in a single
	sentence: \emph{the controller that is tuned on the bench is not the controller
		that flies}. Three mechanisms cause this, and they have different remedies.

	\begin{description}[leftmargin=*,style=nextline]
		\item[Parametric drift.] Mass, inertia, rotor efficiency and motor time
		constants differ between the model and the airframe, and change over the life
		of the vehicle as batteries, propellers and payload mounts are swapped. A
		controller tuned for the nominal parameters degrades gracefully, but a
		\emph{predictive} controller degrades faster, because its predictions are
		wrong over the whole horizon rather than only at the current instant.
		\item[Unmodelled external effects.] Wind, ground effect, aerodynamic
		interactions and slung loads produce forces and moments that no rigid-body
		model contains. These are not parametric: no re-identification of $m$ or
		$\mat J$ removes them.
		\item[Cost-design opacity.] Even with a perfect model, the closed-loop behaviour
		of an MPC is determined by $Q$ and $R$, and the mapping from those weights to
		the behaviour an operator actually wants --- ``follow the path, do not
		saturate, do not oscillate'' --- is neither explicit nor stable across
		operating points. In practice this is resolved by an engineer with a
		stopwatch.
	\end{description}

	Each increment in \cref{ch:method} attacks one of these. \LLTC{} and \ACMPC{}
	attack cost-design opacity; domain randomisation attacks parametric drift by
	averaging over it; the \RDP{} attacks unmodelled external effects by estimating
	them.

	\section{Why a quadrotor, and why this one}

	The platform is a Holybro X500 in the configuration distributed with the PX4
	Gazebo simulation assets. It is a useful test article for three reasons that are
	specific rather than generic.

	\begin{enumerate}[leftmargin=*]
		\item \textbf{The actuator envelope binds.} Every parameter in
		\cref{tab:vehicle} is read from the manufacturer's SDF description. The
		resulting thrust-to-weight ratio is \vTW, which leaves a lateral
		acceleration budget of
		\begin{equation}
			a_\mathrm{lat}^{\max}=\sqrt{\left(\frac{T_{\max}}{m}\right)^{\!2}-g^2}
			=\SI{13.350}{\metre\per\square\second},
			\label{eq:latbudget}
		\end{equation}
		so an aggressive reference is genuinely infeasible rather than merely
		difficult. Any experiment that ignores \cref{eq:latbudget} ends up measuring
		the reference rather than the controller --- a failure mode that occurred in
		this project and is documented in \cref{sec:feasibility}.
		\item \textbf{The inner loop is part of the plant.} The action space is
		\emph{collective thrust and body rates} (\CTBR), the interface PX4 exposes.
		The rate PID, the mixer and the rotor first-order lag therefore sit
		\emph{below} the controller under study and must be simulated, not abstracted
		away. This is what makes the motor commands observable, which
		\cref{sec:rdp-theory} shows is essential for estimating a standing moment.
		\item \textbf{A deployment path exists.} PX4 SITL with Gazebo is a
		high-fidelity, freely available target, and the same ROS\,2 node can address a
		real airframe. Sim-to-real is therefore a question of engineering rather than
		of scope.
	\end{enumerate}

	\begin{table}[htbp]
		\centering
		\caption{Vehicle parameters, all derived from the manufacturer SDF description
			and the \texttt{gz\_x500} airframe configuration. Derived quantities are
			recomputed and asserted by the physics self-test at import time; the values
			below are those printed in \texttt{artifacts/common/nb1\_selftest.csv}.}
		\label{tab:vehicle}
		\small
		\begin{tabularx}{\linewidth}{@{}l X r l@{}}
			\toprule
			\textbf{Symbol} & \textbf{Quantity} & \textbf{Value} & \textbf{Unit} \\
			\midrule
			\multicolumn{4}{@{}l}{\itshape Read from the SDF} \\
			$m_\mathrm{body}$ & Airframe mass & 2.0000 & \si{\kilogram} \\
			$m_\mathrm{rotor}$ & Mass of one rotor assembly & 0.016077 & \si{\kilogram} \\
			$\mat J_\mathrm{body}$ & Airframe inertia, diagonal & $(0.02167,0.02167,0.0400)$ & \si{\kilogram\square\metre} \\
			$\vec r_i$ & Rotor positions ($\pm$\num{0.174}, $\pm$\num{0.174}, \num{0.06}) & --- & \si{\metre} \\
			$K_T$ & Rotor thrust constant & $8.5486\times10^{-6}$ & \si{\newton\second\squared} \\
			$c_\mathrm{drag}$ & SDF rotor drag coefficient & $8.06428\times10^{-5}$ & --- \\
			$c_\mathrm{roll}$ & SDF rolling-moment coefficient & $1\times10^{-6}$ & --- \\
			$\Omega_{\max}$ & Maximum rotor speed & \vOmMax & \si{\radian\per\second} \\
			$\Omega_{\min}$ & Idle rotor speed (PX4 \texttt{SIM\_GZ\_EC\_MIN}) & \vOmMin & \si{\radian\per\second} \\
			$\tau_\uparrow,\tau_\downarrow$ & Rotor spin-up / spin-down time constants & 0.0125 / 0.0250 & \si{\second} \\
			$\vec D$ & Translational drag coefficients (declared addition) & $(0.30,0.30,0.35)$ & \si{\newton\second\per\metre} \\
			\midrule
			\multicolumn{4}{@{}l}{\itshape Derived} \\
			$m$ & Composite mass, $m_\mathrm{body}+4m_\mathrm{rotor}$ & \vMass & \si{\kilogram} \\
			$\mat J$ & Composite inertia about the CG & $(\vJxx,\vJyy,\vJzz)$ & \si{\kilogram\square\metre} \\
			$f_{\max}$ & Maximum thrust of one rotor, $K_T\Omega_{\max}^2$ & \vFmax & \si{\newton} \\
			$T_{\max}$ & Maximum collective thrust & \vTmax & \si{\newton} \\
			$T_\mathrm{idle}$ & \textbf{Idle thrust floor}, $T_{\max}(\Omega_{\min}/\Omega_{\max})^2$ & \vTidle & \si{\newton} \\
			$T/W$ & Thrust-to-weight ratio & \vTW & --- \\
			$u_\mathrm{hover}$ & Hover collective command & \textbf{\vUhover} & --- \\
			$\partial a_z/\partial c$ & Collective authority at hover & \vDadc & \si{\per\square\second} \\
			$a_\mathrm{lat}^{\max}$ & Lateral acceleration budget, \cref{eq:latbudget} & \vAlat & \si{\metre\per\square\second} \\
			\midrule
			\multicolumn{4}{@{}l}{\itshape Allocation} \\
			$k_m^\mathrm{ctrl}$ & Allocator moment constant (\texttt{CA\_ROTORn\_KM}) & \vKmCtrl & \si{\metre} \\
			$\mathrm{cond}(\mat M)$ & Conditioning of the true allocation & \vCondPlant & --- \\
			$\mathrm{cond}(\mat M_\mathrm{ctrl})$ & Conditioning of the allocator's belief & \vCondCtrl & --- \\
			\midrule
			\multicolumn{4}{@{}l}{\itshape Inner loop and timing} \\
			$\vec K_r$ & Rate-loop proportional gain & $(14,14,8)$ & \si{\per\second} \\
			$\vec K_i$ & Rate-loop integral gain & $(4,4,2)$ & --- \\
			$\vec K_d$ & Rate-loop derivative gain & $(0.28,0.28,0.10)$ & --- \\
			$f_\mathrm{gyro}$ & Gyro low-pass cutoff & 40 & \si{\hertz} \\
			$\Delta t_c$ & Outer-loop control period & 0.02 & \si{\second} \\
			$n_\mathrm{sub}$ & Physics sub-steps per control period & 10 & --- \\
			\bottomrule
		\end{tabularx}
	\end{table}

	\begin{keybox}
		\keyhead{Three consequences of taking the airframe file literally}
		The values above differ from a textbook quadrotor model in three ways that each
		changed a result, and all three were found by the audit of
		\cref{sec:audit-defects}.

		\textbf{The rotors never stop.} PX4 holds an idle speed
		$\Omega_{\min}=\vOmMin\,\si{\radian\per\second}$, so a zero collective command
		still produces $T_\mathrm{idle}=\vTidle\,\si{\newton}$, or \SI{3.8}{\percent} of
		vehicle weight. The command-to-thrust map is therefore affine in rotor speed,
		not a pure square, which moves the hover command from \num{0.769} to
		\textbf{\vUhover} and the collective authority from
		\SI{25.49}{\per\square\second} to \vDadc\,\si{\per\square\second}.

		\textbf{The allocator is wrong on purpose.} The mixer inverts a matrix built
		with $k_m^\mathrm{ctrl}=\vKmCtrl\,\si{\metre}$, while the plant produces yaw
		moments through the SDF's rotor-drag coefficient. The two disagree by design,
		exactly as they do on the real vehicle: $\mathrm{cond}(\mat M)=\vCondPlant$
		against $\mathrm{cond}(\mat M_\mathrm{ctrl})=\vCondCtrl$, and only about a
		third of a commanded yaw moment is realised. Modelling the allocator as exact
		would have deleted this error source.

		\textbf{Part of the command box is unreachable.} The autopilot clips body-rate
		commands at \SI{220}{\degree\per\second}, so of the nominal
		$\bodyrate_{\max}=(10,10,4)\,\si{\radian\per\second}$ only
		$(0.384,0.384,0.960)$ of the normalised box can actually be flown. Optimising
		over the full box plans in a region the plant cannot reach.
	\end{keybox}

	\section{Scope and non-goals}

	To keep the comparison meaningful, several things were deliberately excluded.

	\begin{itemize}[leftmargin=*]
		\item \textbf{No perception.} State estimation is assumed; measurement noise is
		injected explicitly as a studied variable (\cref{sec:noise}), but no
		estimator is designed.
		\item \textbf{No trajectory planning.} References are analytic and fixed. The
		subject is tracking, not planning.
		\item \textbf{No hardware flight.} All results are simulated; the PX4/Gazebo
		workspace is the sim-to-real bridge and its experiments are specified in
		\cref{sec:ros-experiments}, but flights on a physical airframe are future
		work.
		\item \textbf{No formal safety certificate.} Robust MPC and safe-RL
		formulations~\cite{zanon2021safe,berkenkamp2017safe} are discussed as the
		natural continuation, not implemented.
	\end{itemize}
	%=======================================================================
	\chapter{State of the art and mathematical framework}
	\label{ch:sota}
	%=======================================================================

	This chapter has two jobs. \Cref{sec:sota} positions the work against the
	literature and states what each prior approach does and does not provide.
	\Crefrange{sec:model}{sec:rl-formulation} then develop the mathematics precisely
	enough that every design choice in \cref{ch:method} is decidable, and every
	discrepancy between an equation and its implementation is detectable. This
	precision is not decoration: \cref{ch:results} reports several defects that were
	found exactly by confronting code with the equation it was supposed to realise.

	%-----------------------------------------------------------------------
	\section{Related work}
	\label{sec:sota}
	%-----------------------------------------------------------------------

	\subsection{Predictive control of multirotors}

	Classical NMPC for quadrotors rests on differential flatness: position and yaw
	are flat outputs, so a smooth reference determines the full state and the
	nominal input in closed form~\cite{mellinger2011minimum}. Faessler et
	al.~\cite{faessler2018differential} extend this to rotor drag, which matters at
	the speeds considered here. Stability and recursive feasibility for constrained
	nonlinear MPC are treated comprehensively in~\cite{mayne2000constrained,
		rawlings2017mpc}; the soft state constraints used in this work are the exact
	$\ell_1$ relaxation of Scokaert and Rawlings~\cite{scokaert1999feasibility},
	which is exact --- i.e.\ the relaxed and unrelaxed solutions coincide whenever a
	feasible trajectory exists --- provided the penalty weight exceeds the norm of
	the optimal multipliers.

	\textbf{Limitation.} All of this presumes the model. Performance degrades with
	model error in a way the theory does not quantify for the multirotor case, and
	the cost weights remain a free design.

	\subsection{Learning the model: residual dynamics}

	The dominant response in aerial robotics is to learn the model error. Shi et
	al.~\cite{shi2019neurallander} learn ground-effect residuals; Torrente et
	al.~\cite{torrente2021gpmpc} place a Gaussian process residual inside an MPC;
	Saviolo et al.~\cite{saviolo2022physics} learn a temporal residual with physical
	priors. The closest published analogue to the predictor used here is
	Neural-Fly~\cite{oconnell2022neuralfly}, which separates a slowly learned basis
	from rapidly adapted linear coefficients and demonstrates agile flight in strong
	wind.

	\textbf{Limitation.} Fitting the model to data does not, in general, make the
	optimiser produce the optimal policy, because the objective being optimised is
	still the wrong one. This is precisely the observation that motivated the
	\emph{modifier adaptation} literature in real-time
	optimisation~\cite{marchetti2009modifier}, where the cost is corrected instead
	of the model.

	\subsection{Learning the policy}

	Model-free RL applied to quadrotors dates to Hwangbo et
	al.~\cite{hwangbo2017quadrl}, and reaches its most visible result in
	champion-level drone racing~\cite{kaufmann2023champion}. The algorithmic
	machinery used in this work --- PPO~\cite{schulman2017ppo},
	GAE~\cite{schulman2016gae} --- comes from this
	line, as does deterministic policy gradient~\cite{silver2014dpg}.

	\textbf{Limitation.} The resulting policy is a network. Constraints are
	represented only through reward shaping; there is nothing to inspect, and safe
	RL remains open~\cite{garcia2015safe,berkenkamp2017safe}.

	\subsection{Learning inside the optimiser: the hybrid}

	Three threads converge on the architecture used here.

	\begin{enumerate}[leftmargin=*]
		\item \textbf{Differentiable optimal control.} Amos et
		al.~\cite{amos2018differentiable} show that an iLQR/LQR solver can be treated
		as a differentiable layer, so a cost can be trained by gradients that pass
		\emph{through} the solve. The underlying sensitivity result is classical:
		under LICQ and second-order sufficiency, the derivative of the optimal value
		with respect to a parameter is the derivative of the Lagrangian at the
		primal--dual solution~\cite{buskens2001sensitivity}, so the gradient is a
		by-product of solving.
		\item \textbf{Learned terminal values.} POLO~\cite{lowrey2019polo} and
		Bhardwaj et al.~\cite{bhardwaj2021blending} replace the terminal cost of a
		short-horizon MPC with a learned value function, recovering long-horizon
		behaviour at short-horizon cost. \LLTC{} is a specialisation of this to a
		quadratic terminal form, which keeps the inner problem a quadratic programme.
		\item \textbf{Learned costs, with theory.} Gros and
		Zanon~\cite{gros2020data} give the formal justification. Their Theorem~1
		states that for an MDP with true transition kernel $\Prob[s^+\mid s,a]$,
		optimal action-value $Q_\star$ and value $V_\star$, and an \emph{approximate}
		model $\hat\Prob[\hat s^+\mid s,a]$, the modified stage cost
		\begin{equation}
			\hat L(s,a)=Q_\star(s,a)-\gamma\,\E\!\left[V_\star(\hat s^+)\mid s,a\right]
			\label{eq:groszanon}
		\end{equation}
		together with the terminal cost $V_\star$ makes the finite-horizon problem
		under the \emph{wrong} model reproduce $\pi_\star$, $V_\star$ and $Q_\star$ of
		the true system. Zanon and Gros~\cite{zanon2021safe} extend the construction to
		robust MPC with safety guarantees.
	\end{enumerate}

	The practical reading of \cref{eq:groszanon} is the one adopted throughout this
	report: \emph{$\hat L$ is not computable, because it requires $Q_\star$ and the
		true kernel; so learn it instead}. \ACMPC{} is exactly this, with the cost map
	playing the role of $\hat L$ and PPO playing the role of the learning rule.

	\subsection{Actor--critic MPC}

	Romero et al.~\cite{romero2024acmpc} instantiate the idea for quadrotors: a
	network emits the stage cost of a short-horizon differentiable MPC, trained by
	an actor--critic method against a task-level reward. They report two findings
	this work attempts to reproduce --- that a one-step horizon suffices when the
	cost is learned, and that richer cost parametrisations (Cholesky, full) fail to
	learn within the diagonal's sample budget. \Cref{sec:res-acmpc} reports the
	outcome of both reproduction attempts, including where they did not reproduce.

	\subsection{Robustness by randomisation}

	Domain randomisation~\cite{tobin2017domain,peng2018simtoreal} trains over a
	distribution of plants so the policy is acceptable on all of them. It is simple
	and effective, and its limitation is structural rather than practical: a policy
	that cannot observe which plant it is on cannot specialise, only compromise.
	That observation is the bridge from \cref{sec:dr} to \cref{sec:adaptive}.

	\begin{keybox}
		\keyhead{Positioning}
		This work does not learn the model and does not learn the policy. It keeps the
		optimal control problem and learns \emph{the cost function inside it}, with
		\cref{eq:groszanon} as the licence, differentiable iLQR as the mechanism, and a
		recurrent wrench estimator supplying the one piece of information a
		disturbance-blind cost cannot infer.
	\end{keybox}

	%-----------------------------------------------------------------------
	\section{Vehicle model}
	\label{sec:model}
	%-----------------------------------------------------------------------

	\subsection{Frames and conventions}

	Two frames are used. The world frame is \textbf{ENU} with gravity $-g\vec e_3$;
	the body frame is \textbf{FLU}. Attitude is a Hamilton, scalar-first unit
	quaternion $\quat=[q_w,q_x,q_y,q_z]$ mapping body to world, with
	$\mat R(\quat)\in\SO(3)$ the corresponding rotation matrix and
	$\vec z_b(\quat)=\mat R(\quat)\vec e_3$ its third column. PX4 uses NED/FRD, so
	the ROS\,2 layer converts on entry and exit; \cref{sec:ros} states that mapping
	and the tests that verify it.

	\begin{remark}[Why the frame table is stated once and enforced]
		A sign error in a frame conversion presents as a control problem, which is the
		most expensive possible misdiagnosis. In this project a documented instance cost
		several days: a world-frame outer loop appeared to circle and then reverse, and
		the cause was a body-rate reconstruction dividing by wall-clock time rather than
		sample time. \Cref{sec:ros} describes the resulting invariant tests.
	\end{remark}

	\subsection{The actuator map}
	\label{sec:actuator}

	Because PX4 holds the rotors at an idle speed $\Omega_{\min}$, the normalised
	command $c\in[0,1]$ maps \emph{affinely} to rotor speed and only then
	quadratically to thrust:
	\begin{equation}
		\Omega(c)=\Omega_{\min}+c\,(\Omega_{\max}-\Omega_{\min}),
		\qquad
		T(c)=T_{\max}\!\left(r+c\,(1-r)\right)^{2},
		\qquad r=\frac{\Omega_{\min}}{\Omega_{\max}}=0.15 .
		\label{eq:thrustmap}
	\end{equation}

	\begin{keybox}
		\keyhead{The idle floor is load-bearing}
		Setting $r=0$ --- the textbook $T=T_{\max}c^2$ --- is wrong on this airframe in
		three compounding ways. It places hover at $c=\num{0.769}$ instead of
		\vUhover; it overstates collective authority at hover,
		\begin{equation}
			\left.\frac{\partial a_z}{\partial c}\right|_\mathrm{hover}
			=\frac{2T_{\max}(1-r)\big(r+u_\mathrm{hover}(1-r)\big)}{m}
			=\SI{21.667}{\per\square\second},
			\label{eq:Bcol}
		\end{equation}
		against \SI{25.49}{\per\square\second} for the pure square, an error of
		\SI{17.6}{\percent}; and it asserts that the vehicle can produce zero thrust,
		when in fact it cannot produce less than $T_\mathrm{idle}=\vTidle\,\si{\newton}$.
		All three mis-scale the LQR gain, the Riccati terminal matrix
		$\mat P_T$, and every terminal cost in the study. A unit test asserts
		\cref{eq:Bcol} against a finite difference of \cref{eq:fc-trans}.
	\end{keybox}

	\subsection{Control model (\texorpdfstring{$17$}{17} states)}
	\label{sec:ctrlmodel}

	The model solved \emph{inside} the optimiser carries the rigid body, the
	autopilot's rate loop and the actuator:
	\begin{equation}
		\ctrlstate=[\underbrace{\vec p}_{3},\underbrace{\vec v}_{3},
		\underbrace{\quat}_{4},\underbrace{\bodyrate}_{3},
		\underbrace{\vec\Omega}_{4}]\in\R^{17},
		\qquad
		\inp=[c,\bar\omega_x,\bar\omega_y,\bar\omega_z]\in\R^{4}.
		\label{eq:ctrlstate}
	\end{equation}
	The input remains the CTBR command the autopilot actually exposes, so nothing
	about the deployment interface is assumed away; what the two extra blocks buy
	is that the rate loop and the actuator are \emph{predicted} rather than taken
	to be instantaneous.

	Every constant below is \emph{nominal}. A controller does not know the true
	inertia, rotor constant or motor time constant, and the mismatch between these
	and the plant's is precisely the model error the adaptive scheme of
	\cref{sec:adaptive} is built to estimate; using the true values here would
	delete the object of study.

	\paragraph{Commanded rate, inner loop and mixer.}
	\begin{align}
		\bodyrate_c &= \clip\!\big(\diag(\bodyrate_{\max})\,\bar{\bodyrate},\
		-\bodyrate_\mathrm{lim},\ \bodyrate_\mathrm{lim}\big),
		&\bodyrate_{\max}&=(10,10,4)\ \si{\radian\per\second},
		\label{eq:unpack}\\[2pt]
		\vec\tau_\mathrm{cmd} &= \mat J_\mathrm{nom}\big(\vec K_r\odot
		(\bodyrate_c-\bodyrate)\big)+\bodyrate\times\mat J_\mathrm{nom}\bodyrate,
		&\bodyrate_\mathrm{lim}&=\SI{220}{\degree\per\second},
		\label{eq:fc-inner}\\[2pt]
		\vec f_\mathrm{cmd} &= \clip\!\left(\mat M_\mathrm{ctrl}^{-1}
		\begin{bmatrix}T(c)\\ \vec\tau_\mathrm{cmd}\end{bmatrix},\ 0,\ f_{\max}\right),
		&\vec\Omega_\mathrm{cmd}&=\clip\!\big(\sqrt{\vec f_\mathrm{cmd}/K_T},\
		\Omega_{\min},\ \Omega_{\max}\big).
		\label{eq:fc-mixer}
	\end{align}

	\paragraph{Rigid body and actuator.}
	\begin{align}
		\dot{\vec p} &= \vec v,
		&\dot{\vec v} &= \frac{K_T\sum_i\Omega_i^{2}}{m}\vec z_b(\quat)-g\vec e_3
		+\vec a_\mathrm{res},
		\label{eq:fc-trans}\\[2pt]
		\dot{\quat} &= \tfrac12\,\quat\otimes
		\begin{bmatrix}0\\ \bodyrate\end{bmatrix},
		&\dot{\bodyrate} &= \mat J_\mathrm{nom}^{-1}\!\left(
		\mat M_\tau^\mathrm{nom}K_T\vec\Omega^{2}
		-\bodyrate\times\mat J_\mathrm{nom}\bodyrate\right)+\vec\alpha_\mathrm{res},
		\label{eq:fc-rot}\\[2pt]
		\dot{\vec\Omega} &= \frac{\vec\Omega_\mathrm{cmd}-\vec\Omega}{\tau_m},
		&\tau_m&=\begin{cases}\tau_\uparrow & \vec\Omega_\mathrm{cmd}\ge\vec\Omega\\
			\tau_\downarrow & \text{otherwise.}\end{cases}
		\label{eq:fc-act}
	\end{align}

	Here $\dres=[\vec a_\mathrm{res},\vec\alpha_\mathrm{res}]\in\R^6$ is an
	\emph{optional model correction}, with units \si{\metre\per\square\second} and
	\si{\radian\per\square\second}. This is the slot the adaptive scheme of
	\cref{sec:adaptive} writes into; \cref{sec:disturbance} shows that, because
	\cref{eq:fc-rot} carries a torque input, an external wrench maps into it
	exactly.

	\begin{proposition}[The model is a cascade]
		\label{prop:cascade}
		In \crefrange{eq:fc-trans}{eq:fc-act} the command $\inp$ enters only through
		$\vec\Omega_\mathrm{cmd}$. Hence $\partial\dot{\ctrlstate}/\partial\inp$ is
		supported on the rotor block alone: the input reaches the body rate only
		through $\vec\Omega$, and the attitude only through $\bodyrate$.
	\end{proposition}

	\begin{proof}
		$\inp$ appears in \crefrange{eq:unpack}{eq:fc-mixer} and nowhere else.
		\Cref{eq:fc-trans,eq:fc-rot} depend on $\ctrlstate$ alone, and
		\cref{eq:fc-act} is the only equation in which $\vec\Omega_\mathrm{cmd}$
		occurs.
	\end{proof}

	\begin{proposition}[Steady-state gains]
		\label{prop:dcgain}
		Hold $\inp$ constant until $\dot{\vec\Omega}=\vec 0$ and
		$\dot{\bodyrate}=\vec 0$. Then the model reproduces the two gains that
		characterise the vehicle:
		\begin{equation}
			\bodyrate_\infty=\diag(\bodyrate_{\max})\,\bar{\bodyrate},
			\qquad
			\left.\frac{\partial a_z}{\partial c}\right|_\mathrm{hover}
			=\frac{2T_{\max}(1-r)\big(r+u_\mathrm{hover}(1-r)\big)}{m},
			\label{eq:dcgain}
		\end{equation}
		the second being \cref{eq:Bcol} exactly.
	\end{proposition}

	\begin{proof}
		At $\dot{\vec\Omega}=\vec 0$ we have $\vec\Omega=\vec\Omega_\mathrm{cmd}$, so
		by \cref{eq:fc-mixer} the realised wrench is
		$\mat M_\mathrm{ctrl}K_T\vec\Omega^{2}=[T(c),\vec\tau_\mathrm{cmd}]\transp$
		whenever the clips are inactive. Substituting the moment part into
		\cref{eq:fc-rot} and setting $\dot{\bodyrate}=\vec 0$ leaves
		$\vec K_r\odot(\bodyrate_c-\bodyrate)=\vec 0$, i.e.\ $\bodyrate=\bodyrate_c$,
		which is \cref{eq:unpack}. For the collective, a symmetric change of $c$
		leaves the moment rows unchanged, so $K_T\sum_i\Omega_i^{2}=T(c)$ and the map
		$c\mapsto a_z$ is \cref{eq:thrustmap} divided by $m$, whose derivative at
		hover is \cref{eq:Bcol}.
	\end{proof}

	\begin{remark}[Why the inner loop is purely proportional]
		The autopilot of \cref{eq:ratepid} is a PID, and its integrator drives
		$\bodyrate_f\to\bodyrate_c$, so its steady-state rate gain is exactly
		$\diag(\bodyrate_{\max})$. Reproducing $\vec K_d$ while omitting $\vec K_i$
		would instead give
		$\bodyrate_\infty=\big(\vec K_r\oslash(\vec K_r+\vec K_d)\big)\odot
		\diag(\bodyrate_{\max})\bar{\bodyrate}$ --- a \SI{2}{\percent} droop on roll
		and pitch. That is a \emph{standing} gain error, not a transient: it
		integrates directly into attitude error over the horizon, which is worse than
		a slightly fast transient. Carrying $\vec I_\omega$ as three further states
		would remove it, at the cost of a mode the horizon barely resolves. Dropping
		the pair is therefore the choice that preserves \cref{prop:dcgain}, and it is
		the reason \cref{eq:fc-inner} has no derivative or integral term.
	\end{remark}

	\paragraph{Integration.} Classical fourth-order Runge--Kutta over
	$\Delta t_c=\SI{0.02}{\second}$ in $n^{c}_\mathrm{sub}=2$ sub-steps, with the
	quaternion renormalised and the rotor speeds kept physical after each. The
	sub-stepping is dictated by the actuator rather than by accuracy in the slow
	states: the rotor pole $\tau_\uparrow=\SI{12.5}{\milli\second}$ is \emph{faster}
	than the control period, so $\Delta t_c/\tau_\uparrow=1.6$ and a single RK4 step
	resolves the fastest mode of \cref{eq:fc-act} only marginally.

	\begin{remark}[The input box is the reachable set]
		Because of the \SI{220}{\degree\per\second} clip in \cref{eq:unpack}, the
		optimiser's box is not $[-1,1]^3$ on the rate channels but
		$\pm(0.384,0.384,0.960)$: the fraction of each axis the plant can actually
		deliver. Enforcing the clip \emph{inside} the dynamics instead would leave zero
		derivative beyond the knee and make $\mat Q_{uu}$ singular in the backward pass;
		shrinking the box achieves the same physics while keeping the solve
		well-conditioned. This is defect~D\nbd 7 of \cref{sec:audit-defects}.
	\end{remark}

	\subsection{Plant model (\texorpdfstring{$23$}{23} states)}

	The simulated plant is deliberately richer than the control model, so that model
	error exists by construction rather than by injection:
	\[
	\plantstate=[\underbrace{\vec p}_{3},\underbrace{\quat}_{4},\underbrace{\vec v}_{3},
	\underbrace{\bodyrate}_{3},\underbrace{\vec\Omega}_{4},
	\underbrace{\vec I_\omega}_{3},\underbrace{\bodyrate_f}_{3}]\in\R^{23},
	\]
	with $\vec\Omega$ the four rotor speeds, $\vec I_\omega$ the rate-loop integral
	state and $\bodyrate_f$ the gyro-filtered rate. It is integrated by RK4 at
	$\Delta t_c/n_\mathrm{sub}=\SI{2}{\milli\second}$.

	\paragraph{Inner rate loop and mixer.}
	\begin{align}
		\vec\tau_\mathrm{cmd} &=
		\mat J\!\left(\vec K_r(\bodyrate_c-\bodyrate_f)+\vec K_i\vec I_\omega
		-\vec K_d\bodyrate_f\right)+\bodyrate\times\mat J\bodyrate,
		\label{eq:ratepid}\\
		\vec f_\mathrm{cmd} &= \clip\!\left(\mat M_\mathrm{ctrl}^{-1}
		\begin{bmatrix}T_c\\ \vec\tau_\mathrm{cmd}\end{bmatrix},\,
		T_\mathrm{idle}/4,\,f_{\max}\right),
		\qquad
		\vec\Omega_\mathrm{cmd}=\sqrt{\vec f_\mathrm{cmd}/K_T},
		\label{eq:mixer}
	\end{align}
	with the allocation matrix built from the arms $\vec a_i=\vec r_i-\mathrm{cg}$
	and the spin directions $\sigma_i=\pm1$:
	\begin{equation}
		\mat M=\begin{bmatrix}\vec 1\transp\\ \mat M_\tau\end{bmatrix},
		\qquad
		\mat M_\tau=\begin{bmatrix}
			a_{1,y} & a_{2,y} & a_{3,y} & a_{4,y}\\
			-a_{1,x} & -a_{2,x} & -a_{3,x} & -a_{4,x}\\
			-\sigma_1 k_m & -\sigma_2 k_m & -\sigma_3 k_m & -\sigma_4 k_m
		\end{bmatrix}.
		\label{eq:mtau}
	\end{equation}

	\begin{cautionbox}
		\textbf{Two different $k_m$, and that is the point.} The \emph{allocator}
		inverts \cref{eq:mtau} with $k_m^\mathrm{ctrl}=\vKmCtrl\,\si{\metre}$, the value
		PX4 carries in \texttt{CA\_ROTORn\_KM}. The \emph{plant} generates yaw moments
		from the SDF's rotor-drag coefficient $8.06428\times10^{-5}$, which corresponds
		to a far smaller effective moment arm. The controller therefore believes it has
		roughly three times the yaw authority it actually has, and about a third of a
		commanded yaw moment is realised. This is not a modelling error to be removed;
		it is a faithful reproduction of the disagreement present on the real vehicle,
		and deleting it by using one constant for both would have made the yaw channel
		artificially easy.
	\end{cautionbox}

	\paragraph{Rigid-body and actuator dynamics.}
	\begin{align}
		\dot{\vec v} &= \frac{\sum_i K_T\Omega_i^2}{m}\vec z_b-g\vec e_3
		-\frac{\vec D\odot(\vec v-\vec v_w)}{m}
		-\frac{c_\mathrm{drag}\sum_i\Omega_i\,\vec v_\perp^{\,b}}{m}
		+\frac{\mat R\,\vec F_\mathrm{ext}^{\,b}}{m},
		\label{eq:fp-trans}\\
		\dot{\bodyrate} &= \mat J^{-1}\!\left(
		\mat M_\tau K_T\vec\Omega^2+\vec\tau_\mathrm{roll}+\vec\tau_\mathrm{ext}
		-\bodyrate\times\mat J\bodyrate\right),
		\label{eq:fp-rot}\\
		\dot{\vec\Omega} &= \frac{\vec\Omega_\mathrm{cmd}-\vec\Omega}{\tau_m},\quad
		\tau_m=\begin{cases}\tau_\uparrow & \vec\Omega_\mathrm{cmd}\ge\vec\Omega\\
			\tau_\downarrow & \text{otherwise,}\end{cases}
		\qquad
		\dot{\bodyrate}_f = 2\pi f_\mathrm{gyro}(\bodyrate-\bodyrate_f),
		\label{eq:fp-act}
	\end{align}
	where $\vec\tau_\mathrm{roll}$ is the SDF rolling moment and the rotor-drag term
	acts on the in-plane body velocity $\vec v_\perp^{\,b}$. Both were absent from
	the first implementation and are defect~A\nbd 5 of
	\cref{sec:audit-defects}.

	\paragraph{Wind.} Per episode, deterministic in $t$, with a horizontal mean
	direction $\hat{\vec d}$ at random azimuth and an in-plane normal
	$\hat{\vec d}^{\perp}$:
	\begin{equation}
		\vec v_w(t)=\bar w\Big[\big(1+0.35\sin(2\pi\cdot0.23\,t+\phi)\big)\hat{\vec d}
		+0.35\sin(2\pi\cdot0.7\,t)\,\hat{\vec d}^{\perp}\Big].
		\label{eq:wind}
	\end{equation}

	\subsection{Payload model}

	A payload of mass $m_p$ at body position $\vec r_p$ shifts the centre of gravity,
	adds inertia by the parallel-axis theorem, \emph{and moves the rotor arms},
	which is what converts an off-centre load into a coupled moment:
	\begin{align}
		\mathrm{cg} &= \frac{m_b\,\mathrm{cg}_0+m_p\vec r_p}{m_b+m_p},
		\label{eq:cg}\\
		\mat J_t &= \lambda_J\mat J_\mathrm{nom}
		+m_b\big(\|\vec\delta_b\|^2\mat I-\vec\delta_b\vec\delta_b\transp\big)
		+m_p\big(\|\vec\delta_p\|^2\mat I-\vec\delta_p\vec\delta_p\transp\big),
		\label{eq:Jpayload}
	\end{align}
	with $\vec\delta_b=\mathrm{cg}_0-\mathrm{cg}$, $\vec\delta_p=\vec r_p-\mathrm{cg}$,
	after which the \emph{plant's} $\mat M_\tau$ is rebuilt from
	$\vec a_i=\vec r_i-\mathrm{cg}$ while the \emph{allocator} keeps its nominal
	geometry --- the separation of plant truth from controller belief that
	defect~M\nbd 1 required.

	%-----------------------------------------------------------------------
	\section{Error coordinates}
	\label{sec:error}
	%-----------------------------------------------------------------------

	Tracking is expressed in a reduced $16$-dimensional error, which removes the
	quaternion's redundancy while remaining a local diffeomorphism:
	\begin{equation}
		\err=\begin{bmatrix}\vec p-\vec p_\refsym\\ \vec v-\vec v_\refsym\\
			2\,\sgn(q_{e,w})\,\vec q_{e,v}\\
			\bodyrate-\bodyrate_\refsym\\
			(\vec\Omega-\vec\Omega_\refsym)/\Omega_\mathrm{sc}
		\end{bmatrix}\in\R^{16},
		\qquad \quat_e=\quat_\refsym^{*}\otimes\quat.
		\label{eq:err}
	\end{equation}
	The sign selects the shorter rotation. For a tilt $\vartheta$ about a unit axis,
	$\|\err_{7:9}\|=2\sin(\vartheta/2)$, so the parametrisation is exact for
	$\vartheta<\pi$ and first-order equivalent to a rotation vector. The rate and
	rotor blocks are plain vector differences --- neither lives on a manifold ---
	so \cref{lem:inverse} inverts the whole of \cref{eq:err}, not merely its
	attitude part.

	\begin{lemma}[Exact inverse of the attitude error]
		\label{lem:inverse}
		Let $\vec\delta=2\sgn(q_{e,w})\vec q_{e,v}$ with $\|\vec\delta\|\le2$. The unique
		$\quat_e$ in the upper hemisphere consistent with $\vec\delta$ is
		\begin{equation}
			\quat_e=\big[\cos(n/2),\ \hat{\vec\delta}\sin(n/2)\big],
			\qquad n = 2\arcsin\!\big(\|\vec\delta\|/2\big),
			\qquad \hat{\vec\delta}=\vec\delta/\|\vec\delta\|.
			\label{eq:einv}
		\end{equation}
	\end{lemma}

	\begin{proof}
		Write $\quat_e=[\cos(\vartheta/2),\hat{\vec n}\sin(\vartheta/2)]$ with
		$\vartheta\in[0,\pi]$ and $\hat{\vec n}$ a unit axis, which is the unique
		hemisphere-normalised representation. Then
		$\vec\delta=2\sin(\vartheta/2)\hat{\vec n}$, so
		$\hat{\vec\delta}=\hat{\vec n}$ and
		$\|\vec\delta\|=2\sin(\vartheta/2)$. Since $\vartheta/2\in[0,\pi/2]$, the sine
		is injective there and $\vartheta=2\arcsin(\|\vec\delta\|/2)$.
	\end{proof}

	\begin{remark}[Why \cref{lem:inverse} is stated explicitly]
		The first-order approximation $n\approx\|\vec\delta\|$ is exact to
		$\mathcal O(\vartheta^3)$ and therefore invisible at small tilt, but the error
		grows quickly: $\num{-0.04}\si{\degree}$ at \SI{15}{\degree},
		$\num{-2.70}\si{\degree}$ at \SI{60}{\degree} and $\num{-20.8}\si{\degree}$ at
		\SI{120}{\degree}. The implementation used the approximation; this is
		defect~E\nbd 1 in \cref{sec:audit-defects}, low severity but exactly
		quantifiable, and a one-line correction.
	\end{remark}

	\begin{remark}[The rotor scale is a change of coordinates, not a weight]
		\label{rem:omsc}
		$\Omega_\mathrm{sc}=\Omega_{\max}$ enters \cref{eq:err} and its inverse
		symmetrically, so it is a similarity transform of the error space: no
		trajectory, gain or optimal cost depends on the value chosen. Its purpose is
		conditioning. Rotor speeds are $\mathcal O(10^{3})\,\si{\radian\per\second}$
		while the remaining blocks are $\mathcal O(10^{-1})$, and the learned stage
		weights of \cref{sec:acmpc} already span five decades; leaving the rotor block
		unscaled would make the conditioning of $\mat Q_{uu}$ a function of the units
		in which rotor speed happens to be expressed.
	\end{remark}

	\begin{remark}[The reference rotor speed, and why $\err=\vec 0$ is exact at trim]
		\label{rem:omref}
		Differential flatness fixes $\vec p_\refsym$, $\vec v_\refsym$ and
		$\quat_\refsym$, and a centred difference of $\quat_\refsym$ fixes
		$\bodyrate_\refsym$. The reference rotor speed follows by inverting the
		nominal allocator at the reference wrench,
		\begin{equation}
			\vec\tau_\refsym=\mat J_\mathrm{nom}\dot{\bodyrate}_\refsym
			+\bodyrate_\refsym\times\mat J_\mathrm{nom}\bodyrate_\refsym,
			\qquad
			\vec\Omega_\refsym=\clip\!\left(\sqrt{\mat M_\mathrm{ctrl}^{-1}
				\begin{bmatrix}T_\refsym\\ \vec\tau_\refsym\end{bmatrix}\Big/K_T},\
			\Omega_{\min},\ \Omega_{\max}\right),
			\label{eq:omref}
		\end{equation}
		with $\dot{\bodyrate}_\refsym$ a centred difference of $\bodyrate_\refsym$ ---
		so no closed-form jerk or snap of the reference path is required. At hover
		$\vec\tau_\refsym=\vec 0$ and $T_\refsym=mg$, giving
		$\vec\Omega_\refsym=\sqrt{mg/4K_T}$, which is exactly the trimmed rotor speed.
		Hence $\err$ vanishes identically on the trimmed hover state, and $\err=\vec 0$
		is a genuine equilibrium of \cref{eq:edyn} rather than an approximate one.
	\end{remark}

	%-----------------------------------------------------------------------
	\section{References and feasibility}
	\label{sec:feasibility}
	%-----------------------------------------------------------------------

	Six reference families are used, all of the form $\theta=\omega t+\phi_0$ about
	a centre, with radius $R\sim\mathcal U(0.5,2.0)\,\si{\metre}$ and
	$\phi_0\sim\mathcal U(0,2\pi)$.

	\begin{table}[htbp]
		\centering
		\caption{Reference families. $\kappa_v,\kappa_a$ are the peak speed and peak
			acceleration at $R=1$, $\omega=1$; actual peaks are $R\omega\kappa_v$ and
			$R\omega^2\kappa_a$. The $\kappa_a$ column is asserted by the physics
			self-test.}
		\label{tab:paths}
		\small
		\begin{tabularx}{\linewidth}{@{}l X r r@{}}
			\toprule
			\textbf{Kind} & \textbf{$\vec p_\refsym(\theta)$} & $\kappa_v$ & $\kappa_a$ \\
			\midrule
			hover  & $\vec c$ & 0 & 0 \\
			step   & $\vec c+\vec\Delta$, $\vec\Delta$ fixed per episode & 0 & 0 \\
			circle & $\vec c+R(\cos\theta,\ \sin\theta,\ 0)$ & 1.0000 & 1.0000 \\
			fig8   & $\vec c+R(\sin\theta,\ \sin\theta\cos\theta,\ 0)$ (Gerono lemniscate) & 1.4142 & 2.1250 \\
			helix  & circle $+\,0.25\,\mathrm{spd}\cdot t\,\vec e_3$ & 1.0308$^{\dagger}$ & 1.0000 \\
			square & $\vec c+R\,r(\theta)(\cos\theta,\sin\theta,0)$, $r=(\cos^8\theta+\sin^8\theta)^{-1/8}$ & 1.4160 & 9.0779 \\
			\bottomrule
			\multicolumn{4}{@{}l@{}}{\footnotesize $^{\dagger}$ valid only at $\omega=\mathrm{spd}/R$; see defect~P\nbd 1, \cref{app:audit}.}
		\end{tabularx}
	\end{table}

	\begin{proposition}[Peak acceleration of the lemniscate]
		\label{prop:fig8}
		For $\vec p_\refsym=R(\sin\theta,\sin\theta\cos\theta,0)$ with $\theta=\omega t$,
		$\max_\theta\|\ddot{\vec p}_\refsym\|=R\omega^2\sqrt{4.5156}=2.1250\,R\omega^2$.
	\end{proposition}

	\begin{proof}
		Differentiating twice,
		\[
		\ddot{\vec p}_\refsym=R\omega^2(-\sin\theta,-2\sin2\theta,0),
		\qquad
		\frac{\|\ddot{\vec p}_\refsym\|^2}{(R\omega^2)^2}=\sin^2\theta+4\sin^2 2\theta.
		\]
		With $u=\sin^2\theta\in[0,1]$ and $\sin^2 2\theta=4u(1-u)$ this is
		$g(u)=17u-16u^2$, a downward parabola maximised at $u^{*}=17/32$ with
		$g(u^{*})=289/64=4.5156$.
	\end{proof}

	\paragraph{Feasibility capping.} Rather than set $\omega=\mathrm{spd}/R$, the
	angular rate is capped so that the demand stays inside a fraction
	$\alpha=0.60$ of the envelope \cref{eq:latbudget}:
	\begin{equation}
		\omega=\min\!\left(\frac{\mathrm{spd}}{\max(R,0.3)},\
		\sqrt{\frac{\alpha\,a_\mathrm{lat}^{\max}}{R\,\kappa_a}}\right).
		\label{eq:omegacap}
	\end{equation}

	\begin{cautionbox}
		\textbf{Why \cref{eq:omegacap} exists.} Without it, the superellipse at
		$R=\SI{0.5}{\metre}$ and \SI{1.5}{\metre\per\second} demands
		\SI{40.85}{\metre\per\square\second} --- \num{3.06} times the whole envelope
		and \num{5.10} times the design cap, as
		\texttt{artifacts/common/path\_feasibility.csv} records. Every controller then
		fails by the same large margin, and the \texttt{square} column measures the
		reference's infeasibility rather than any property of the controller. The
		measured outcome of the capping is that \texttt{square} is clipped on
		\SI{100}{\percent} of radii and \texttt{fig8} on \SI{6.2}{\percent}, while
		\texttt{circle} and \texttt{helix} are never clipped.
	\end{cautionbox}

	\paragraph{Reference attitude.} Yaw is held at zero and the reference attitude is
	the unique tilt-only rotation aligning $\vec e_3$ with the required thrust
	direction --- the standard flatness construction~\cite{mellinger2011minimum}:
	\begin{equation}
		\vec z_b^{\refsym}=\frac{\vec a_\refsym+g\vec e_3}{\|\vec a_\refsym+g\vec e_3\|},
		\qquad
		\quat_\refsym=\exp\!\left(\tfrac12\,
		\angle(\vec e_3,\vec z_b^{\refsym})\
		\widehat{\vec e_3\times\vec z_b^{\refsym}}\right).
		\label{eq:qref}
	\end{equation}

	\realfig{62mm}{\figFeasibility}
	{Reference feasibility. Peak demanded $\|\vec a_\refsym\|$ against radius for
		each path family, before and after the capping \cref{eq:omegacap}, against the
		envelope $a_\mathrm{lat}^{\max}=\vAlat\,\si{\metre\per\square\second}$ and the
		design fraction $\alpha a_\mathrm{lat}^{\max}=\vAlatCap\,\si{\metre\per\square\second}$.
		Curves above the envelope cannot be tracked by any controller and would make the
		corresponding benchmark column meaningless.}
	{fig:feasibility}

	%-----------------------------------------------------------------------
	\section{The differentiable predictive control layer}
	\label{sec:mpc-layer}
	%-----------------------------------------------------------------------

	\subsection{Error dynamics along a moving reference}

	Stage $k$ of the horizon predicts the state at $t+k\Delta t_c$ and must therefore
	be scored against $\ctrlstate_{\refsym,k}$ at that time, not against
	$\ctrlstate_\refsym(t)$:
	\begin{equation}
		\err_{k+1}=\mathrm{err}\Big(
		\Phi\big(\iota(\err_k,\ctrlstate_{\refsym,k}),\ \inp_\refsym+\delta\inp_k,\ \dres\big),\
		\ctrlstate_{\refsym,k+1}\Big),
		\label{eq:edyn}
	\end{equation}
	where $\Phi$ is one RK4 step of \crefrange{eq:fc-trans}{eq:fc-rot}, $\iota$ is
	the inverse map of \cref{lem:inverse}, and $\mathrm{err}$ is \cref{eq:err}.

	\begin{cautionbox}
		\textbf{This is not optional.} If the reference is frozen at
		$\ctrlstate_\refsym(t)$ over the horizon, the solver plans to come to rest on the
		current waypoint. \Cref{tab:res-horizon} is built to expose this: with the
		preview supplied, tracking error should fall with $N$ and saturate; genuinely
		frozen, it should \emph{rise} with $N$. An assertion fails if the two families
		of columns ever coincide, which is how the original implementation's silent
		failure was detected.
	\end{cautionbox}

	\subsection{The optimal control problem}

	With $N_\tau=N_e+N_u=20$, stage weights $\mat S_k\in\mathbb S^{20}_{+}$ and
	linear terms $\vec c_k\in\R^{20}$:
	\begin{equation}
		\begin{aligned}
			\min_{\delta\inp_{0:N-1}}\quad
			& \sum_{k=0}^{N-1}\left(
			\tfrac12 \begin{bmatrix}\err_k\\ \delta\inp_k\end{bmatrix}\transp
			\mat S_k \begin{bmatrix}\err_k\\ \delta\inp_k\end{bmatrix}
			+\vec c_k\transp\begin{bmatrix}\err_k\\ \delta\inp_k\end{bmatrix}\right)
			+\tfrac12\err_N\transp \mat P_\mathrm{term}\,\err_N\\[2pt]
			\st & \err_{k+1}=\text{\cref{eq:edyn}},\qquad
			\delta\inp_k\in[\,\underline{\vec u}-\inp_\refsym,\ \overline{\vec u}-\inp_\refsym\,].
		\end{aligned}
		\label{eq:ocp}
	\end{equation}

	State constraints, where used, are relaxed exactly in the sense
	of~\cite{scokaert1999feasibility}:
	\begin{equation}
		h_\theta(\ctrlstate_k,\inp_k)\le\vec\sigma_k,\qquad
		\text{penalty } \vec w\transp\vec\sigma_k,\quad \vec\sigma_k\ge0 .
		\label{eq:l1relax}
	\end{equation}

	\begin{remark}[Why the relaxation is necessary for learning, not merely convenient]
		Without it the value function is $+\infty$ on any constraint violation, and an
		infinite value is useless in a temporal-difference update. The relaxation keeps
		the value finite so the learning signal remains defined. Stated generally: any
		form of learning is meaningless if infinite penalties are assigned to violating
		limitations~\cite{gros2020data}.
	\end{remark}

	\subsection{iLQR and differentiation}

	\Cref{eq:ocp} is solved by iLQR~\cite{li2004ilqr,tassa2012synthesis}: a forward
	rollout, a backward Riccati recursion on the linearised dynamics with
	Levenberg--Marquardt regularisation $\mu$, and a backtracking line search, with
	the input box handled by clamping inside the rollout.

	\begin{algorithm}[htbp]
		\caption{iLQR solve used as a differentiable layer}
		\label{alg:ilqr}
		\begin{algorithmic}[1]
			\Require $\err_0$, reference preview $\ctrlstate_{\refsym,0:N}$, cost
			$(\mat S,\vec c,\mat P_\mathrm{term})$, residual $\dres$, iterations $n_\mathrm{it}$,
			differentiated tail $n_\mathrm{diff}$
			\State initialise $\delta\inp_{0:N-1}\gets \vec 0$ (hover)
			\For{$i=1,\dots,n_\mathrm{it}$}
			\State roll out \cref{eq:edyn} to obtain $\err_{0:N}$; evaluate $J$
			\State linearise: $\mat A_k=\partial\err_{k+1}/\partial\err_k$,\ \
			$\mat B_k=\partial\err_{k+1}/\partial\delta\inp_k$
			\State backward pass: $\mat V_N\gets\mat P_\mathrm{term}$; recurse
			$\mat Q_{uu},\mat Q_{ux}$; gains $\mat K_k,\vec k_k$
			\State line search over $\alpha\in\{1,\tfrac12,\dots\}$; accept on sufficient
			decrease, else increase $\mu$
			\If{$i>n_\mathrm{it}-n_\mathrm{diff}$} \State keep the tape (differentiated)
			\Else \State \textbf{stop-gradient} (fixed point treated as constant)
			\EndIf
			\EndFor
			\State \Return $\delta\inp_0$, $J$, trajectories $\err_{0:N},\delta\inp_{0:N-1}$
		\end{algorithmic}
	\end{algorithm}

	Only the last $n_\mathrm{diff}$ iterations are differentiated
	($n_\mathrm{diff}=2$ during training, $1$ at evaluation), which trades the exact
	implicit gradient of~\cite{amos2018differentiable} against memory. The
	justification for treating the earlier iterations as constant is the sensitivity
	result of~\cite{buskens2001sensitivity}: at a solution satisfying LICQ and
	second-order sufficiency,
	\begin{equation}
		\nabla_\theta Q_\theta(s,a)=\nabla_\theta\mathcal L_\theta(s,\vec y^{\star}),
		\label{eq:sensitivity}
	\end{equation}
	so the gradient depends on the parameters only through the Lagrangian evaluated
	at the solution, not on the path taken to reach it.

	\begin{remark}[Non-differentiability at weakly active constraints]
		\Cref{eq:sensitivity} is well defined whenever no inequality is weakly active;
		at a weakly active constraint only subgradients exist. In practice this is
		avoided by an interior-point-like treatment, and the clamp derivative convention
		is pinned by a custom JVP to match the reference implementation --- a measured
		$0.297$ absolute discrepancy in $\mat B$ arises from getting it wrong, which is
		enough to send the line search down a different descent path.
	\end{remark}

	%-----------------------------------------------------------------------
	\section{Disturbance representations}
	\label{sec:disturbance}
	%-----------------------------------------------------------------------

	Two distinct $6$-vectors describe ``the disturbance'' and they are \emph{not}
	interchangeable. Confusing them is the highest-severity defect found in this
	project, so both are defined here.

	\begin{definition}[Model residual]
		\label{def:residual}
		\begin{equation}
			\dres=\begin{bmatrix}
				\vec a_\mathrm{plant}-\vec a_\mathrm{model}\\
				\dot{\bodyrate}_\mathrm{plant}-\dot{\bodyrate}_\mathrm{model}
			\end{bmatrix}\in\R^6,
			\label{eq:dres}
		\end{equation}
		where $\vec a_\mathrm{model}$ and $\dot{\bodyrate}_\mathrm{model}$ are the
		right-hand sides of \cref{eq:fc-trans,eq:fc-rot} evaluated with
		$\dres=\vec 0$. Units $[\si{\metre\per\square\second},\
		\si{\radian\per\square\second}]$. This is exactly the slot $\dres$ occupies in
		\crefrange{eq:fc-trans}{eq:fc-rot}: both blocks are what the nominal model
		fails to account for, in the coordinates the model itself uses.
	\end{definition}

	\begin{definition}[External wrench]
		\label{def:wrench}
		\begin{align}
			\vec F_\mathrm{ext} &= m\,\vec a_\mathrm{true}
			-\big(T\vec z_b-mg\vec e_3\big) && \text{(world, \si{\newton})},
			\label{eq:Fext}\\
			\vec\tau_\mathrm{ext} &= \mat J_\mathrm{nom}\dot{\bodyrate}_\mathrm{true}
			+\bodyrate\times\mat J_\mathrm{nom}\bodyrate
			-\mat M_\tau^\mathrm{nom}K_T\vec\Omega^2 && \text{(body, \si{\newton\metre})}.
			\label{eq:tauext}
		\end{align}
	\end{definition}

	\begin{proposition}[Conversion]
		\label{prop:convert}
		\Cref{def:residual,def:wrench} are two coordinate representations of the same
		disturbance, related by the invertible, constant, block-diagonal map
		\begin{equation}
			\dres=\begin{bmatrix}m^{-1}\mat I_3 & \mat 0\\[2pt]
				\mat 0 & \mat J_\mathrm{nom}^{-1}\end{bmatrix}\wrench,
			\qquad\text{that is}\qquad
			\vec a_\mathrm{res}=\frac{\vec F_\mathrm{ext}}{m},
			\quad
			\vec\alpha_\mathrm{res}=\mat J_\mathrm{nom}^{-1}\vec\tau_\mathrm{ext}.
			\label{eq:convert}
		\end{equation}
		Both blocks are exact, and the map is invertible because $m>0$ and
		$\mat J_\mathrm{nom}\succ0$.
	\end{proposition}

	\begin{proof}
		The force block is \cref{eq:Fext} divided by $m$, which is the definition of
		$\vec a_\mathrm{res}$ in \cref{eq:fc-trans}. For the moment block,
		\cref{eq:tauext} defines $\vec\tau_\mathrm{ext}$ as the moment the nominal
		rigid body fails to resolve, while \cref{eq:fc-rot} maps any resolved moment
		to angular acceleration through $\mat J_\mathrm{nom}^{-1}$ and carries
		$\vec\alpha_\mathrm{res}$ additively on the same line. An unmodelled moment
		$\vec\tau_\mathrm{ext}$ is therefore exactly an unmodelled angular
		acceleration $\mat J_\mathrm{nom}^{-1}\vec\tau_\mathrm{ext}$.
	\end{proof}

	\begin{remark}[Exactness is structural]
		\label{rem:exactness}
		\Cref{eq:convert} holds because \cref{eq:fc-rot} has a torque input. A control
		model whose deepest rotational state is attitude has no image for the moment
		block at all: a moment can then only be represented through the rate offset it
		eventually produces, and that offset is a property of the autopilot's
		closed-loop response --- it depends on how long the moment has been acting
		relative to the rate loop's integral time constant, so the same physical
		disturbance has different representations at different times. Carrying
		$\bodyrate$ as a state removes that dependence. One consequence is that the
		estimator of \cref{sec:adaptive} can be trained against the wrench --- what a
		force--torque sensor measures, and what \cref{def:wrench} is expressed in ---
		and its prediction converted for the optimiser with no modelling step in
		between.
	\end{remark}

	\begin{remark}[What the exact conversion does \emph{not} remove]
		\Cref{prop:convert} is exact as a change of coordinates between the two
		descriptions of the disturbance. It says nothing about whether the nominal
		$m$ and $\mat J_\mathrm{nom}$ used in it are the true ones --- by construction
		they are not, and the difference is part of what the estimator must learn. Nor
		does it make the \emph{prediction} exact: \cref{eq:fc-inner} is a proportional
		reduction of a PID rate loop, so a moment that the autopilot's integrator has
		already absorbed is still represented here as a persistent angular
		acceleration.
	\end{remark}

	\begin{cautionbox}
		\textbf{Numerical illustration.} With the study's randomisation magnitudes
		$F_\mathrm{DR}=0.32\,mg=\SI{6.478}{\newton}$ and
		$\tau_\mathrm{DR}=\SI{0.3572}{\newton\metre}$: feeding the raw wrench into the
		residual slot injects \SI{6.478}{\metre\per\square\second} where
		\SI{3.138}{\metre\per\square\second} is correct --- a factor $m=\vMass$ --- and
		\SI{0.3572}{\radian\per\square\second} where \cref{eq:convert} gives
		$\tau_\mathrm{DR}/J_{xx}=\SI{14.99}{\radian\per\square\second}$, short by the
		factor $\mat J_\mathrm{nom}$. The two errors point in opposite directions ---
		the force block too large, the moment block too small --- which is exactly why
		neither is caught by a sanity check on magnitude. See defect~D\nbd 1.
	\end{cautionbox}

	%-----------------------------------------------------------------------
	\section{Reinforcement learning formulation}
	\label{sec:rl-formulation}
	%-----------------------------------------------------------------------

	\subsection{Markov decision process}

	The closed loop is an MDP $(\mathcal S,\mathcal A,\Prob,r,\gamma)$ with state
	the plant state augmented by the observation history, action the CTBR command,
	and $\Prob$ induced by \crefrange{eq:fp-trans}{eq:fp-act} together with the
	episode randomisation of \cref{ch:data}. The Bellman equations are standard:
	\begin{equation}
		Q_\star(s,a)=r(s,a)+\gamma\,\E\!\left[V_\star(s^+)\mid s,a\right],
		\qquad V_\star(s)=\max_a Q_\star(s,a).
		\label{eq:bellman}
	\end{equation}

	\subsection{Observation}

	The policy is \emph{disturbance-blind} by default, receiving
	$n_o=47$ features:
	\begin{equation}
		\vec o=\Big[\underbrace{\err}_{16}\ \Big\|\
		\underbrace{\{\vec p_\refsym(t{+}5h\Delta t_c)-\vec p,\ \vec v_\refsym(t{+}5h\Delta t_c)\}_{h=1}^{3}}_{18}\ \Big\|\
		\underbrace{\clip(\textstyle\int\!\err_p,\pm5)}_{3}\ \Big\|\
		\underbrace{\overline{\delta\inp}}_{4}\ \Big\|\
		\underbrace{\overline{\err_v}}_{3}\ \Big\|\
		\underbrace{\overline{\bodyrate}}_{3}\Big],
		\label{eq:obs}
	\end{equation}
	extended by $\dres\in\R^6$ when the oracle channel is enabled. The preview
	stride is five control steps, i.e.\ lookaheads of $0.1/0.2/\SI{0.3}{\second}$.
	The observation is normalised by a running mean and variance whose variance is
	floored, for the reason given in \cref{sec:audit-defects}.

	\subsection{Reward}

	Two modes are used. The interpretable quadratic reward, following
	\cite{romero2024acmpc}, makes the value function directly comparable with the
	MPC cost:
	\begin{equation}
		r_\mathrm{quad}=-\frac{\err\transp\mat Q\err+\delta\inp\transp\mat R\,\delta\inp}{100}
		-0.02\|\bodyrate\|^2-0.05\|\Delta\inp\|^2-c_\mathrm{crash}\mathbb 1_\mathrm{crash},
		\label{eq:rquad}
	\end{equation}
	and a shaped progress reward for the tracking curriculum,
	\begin{equation}
		r_\mathrm{prog}=(d_k-d_{k+1})-0.6\,d_{k+1}^2-0.02\|\bodyrate\|^2
		-0.05\|\Delta\inp\|^2+0.05-c_\mathrm{crash}\mathbb 1_\mathrm{crash},
		\label{eq:rprog}
	\end{equation}
	with $d_k=\|\vec p_k-\vec p_{\refsym,k}\|$ and $c_\mathrm{crash}=5$.

	\begin{definition}[Termination]
		\label{def:crash}
		An episode terminates when $p_z\le\SI{0.02}{\metre}$, or
		$\|\vec p-\vec p_\refsym\|>\SI{3}{\metre}$, or the state is non-finite, or
		$\|\bodyrate\|>\SI{25}{\radian\per\second}$.
	\end{definition}

	\begin{cautionbox}
		\textbf{The \SI{3}{\metre} bound is a measurement artefact waiting to happen.}
		Under a sufficiently severe disturbance every controller terminates at the bound,
		so every reported RMSE converges to the same value. An evaluation in which all
		controllers agree to within a fraction of a metre is reporting the cap, not the
		tracking. \Cref{sec:moderation} shows this occurred and how it was corrected.
	\end{cautionbox}

	\subsection{Policy optimisation}

	The policy is optimised by PPO~\cite{schulman2017ppo} with
	GAE($\lambda$)~\cite{schulman2016gae}. The clipped surrogate is
	\begin{equation}
		\mathcal L^\mathrm{CLIP}(\theta)=\E_t\!\left[
		\min\!\big(\rho_t(\theta)\hat A_t,\
		\clip(\rho_t(\theta),1-\epsilon,1+\epsilon)\hat A_t\big)\right],
		\quad
		\rho_t(\theta)=\frac{\pi_\theta(a_t\mid s_t)}{\pi_{\theta_\mathrm{old}}(a_t\mid s_t)},
		\label{eq:ppo}
	\end{equation}
	\begin{equation}
		\hat A_t=\sum_{l=0}^{\infty}(\gamma\lambda)^l\delta_{t+l},
		\qquad \delta_t=r_t+\gamma V(s_{t+1})-V(s_t).
		\label{eq:gae}
	\end{equation}

	\begin{cautionbox}
		\textbf{Where the exploration noise must go.} \Cref{eq:ppo} is only meaningful
		if $\pi_\theta$ is the distribution the action was actually drawn from. Sampling
		the noise anywhere other than the action --- for instance perturbing the cost
		map and then solving deterministically --- makes $\rho_t$ the ratio of two
		densities that never generated anything, and the clipped surrogate optimises a
		quantity unrelated to the rollout. The action is therefore drawn as
		$a=\mu_\theta(s)+\exp(\log\sigma)\,\varepsilon$ with $\mu_\theta$ the MPC
		solution, and $\log\sigma$ is itself learned. This is defect~B\nbd 1,
		and correcting it changed the character of the smoke-scale results
		(\cref{sec:res-comparison}).
	\end{cautionbox}

	%=======================================================================
	\chapter{Data and exploratory analysis}
	\label{ch:data}
	%=======================================================================

	In a simulation-driven study the ``data'' is the environment: the distributions
	from which episodes are drawn, the signals recorded, and the sampling design.
	This chapter documents all three, together with the exploratory analysis that
	led to two corrections of the experimental design.

	%-----------------------------------------------------------------------
	\section{Data-generating process}
	%-----------------------------------------------------------------------

	Each episode samples, independently: a path kind, a centre
	$\vec c\sim\mathcal U([-1.5,1.5]^2\times[1.0,2.5])$, a radius
	$R\sim\mathcal U(0.5,2.0)$, a phase $\phi_0\sim\mathcal U(0,2\pi)$, a nominal
	speed, a wind magnitude and azimuth, and a vector of plant scale factors
	$(\lambda_m,\lambda_D,\lambda_\tau,\lambda_{T},\lambda_{K_w},\lambda_J)$ applied
	to mass, drag, rotor lag, maximum thrust, rate gains and inertia. The initial
	condition is offset from the reference by
	$\vec e_p(0)\sim\mathcal U([-0.5,0.5]^3)$.

	The environment is batched: $n$ vehicles are simulated in parallel as a single
	array program, each with its own episode parameters. This is what makes
	$\num{512}$-vehicle PPO rollouts and $\num{256}$-vehicle evaluations tractable.

	\begin{table}[htbp]
		\centering
		\caption{Disturbance suites. Multiplicative entries scale the nominal parameter;
			wind and speed are absolute. The moderated columns are produced by the
			contraction \cref{eq:moderate} and are used for all evaluation reported in
			\cref{ch:results}; the raw columns remain the training distribution. Bands
			are those printed by \texttt{study\_moderate.py} at run time.}
		\label{tab:dist}
		\small
		\begin{tabularx}{\linewidth}{@{}l X r r r r@{}}
			\toprule
			& & \multicolumn{2}{c}{\textbf{In distribution}} & \multicolumn{2}{c}{\textbf{Out of distribution}} \\
			\cmidrule(lr){3-4}\cmidrule(lr){5-6}
			\textbf{Channel} & \textbf{Physical meaning} & raw & moderated & raw & moderated \\
			\midrule
			wind [\si{\metre\per\second}] & Mean wind magnitude, \cref{eq:wind} & 0--15 & \textbf{0--2} & 18--25 & \textbf{2.5--4} \\
			speed [\si{\metre\per\second}] & Nominal reference speed & 0.5--4.0 & 0.5--4.0 & 5.0--6.5 & 5.0--6.5 \\
			$\lambda_m$ & Mass & 0.80--1.20 & 0.970--1.030 & 0.65--1.35 & 0.948--1.052 \\
			$\lambda_D$ & Translational drag & 0.50--2.00 & 0.925--1.150 & 2.20--3.00 & 1.180--1.300 \\
			$\lambda_\tau$ & Rotor time constant & 0.50--3.00 & 0.925--1.300 & 3.50--5.00 & 1.375--1.600 \\
			$\lambda_{T}$ & Maximum thrust & 0.85--1.15 & 0.978--1.022 & 0.75--1.30 & 0.963--1.045 \\
			$\lambda_{K_w}$ & Inner-loop rate gains & 0.60--1.40 & 0.940--1.060 & 0.40--1.80 & 0.910--1.120 \\
			$\lambda_J$ & Inertia tensor & 0.70--1.40 & 0.955--1.060 & 0.50--1.80 & 0.925--1.120 \\
			\bottomrule
		\end{tabularx}
	\end{table}

	%-----------------------------------------------------------------------
	\section{Exploratory analysis I: the reference must be flyable}
	%-----------------------------------------------------------------------

	The first exploratory step was to compute the acceleration each reference
	family actually demands and compare it with \cref{eq:latbudget}. This revealed
	that the superellipse family, introduced precisely to expose the value of a
	prediction horizon at the corners, was demanding
	\SI{40.85}{\metre\per\square\second} --- \num{3.06} times the available
	\vAlat\,\si{\metre\per\square\second}. \Cref{eq:omegacap} was introduced in
	response, and a unit test now asserts that every family stays inside
	$\alpha\,a_\mathrm{lat}^{\max}=\vAlatCap\,\si{\metre\per\square\second}$.
	\Cref{fig:feasibility} is the resulting audit.

	%-----------------------------------------------------------------------
	\section{Exploratory analysis II: the disturbance suites were saturating}
	\label{sec:moderation}
	%-----------------------------------------------------------------------

	The second exploratory step examined the evaluation suites of
	\cref{tab:dist}. The out-of-distribution suite reaches
	\SI{25}{\metre\per\second} of wind against an airframe with a
	\vAlat\,\si{\metre\per\square\second} lateral budget. The symptoms were
	unambiguous: actuator saturation of \SIrange{41}{78}{\percent}, every controller
	landing within \SI{0.3}{\metre} of every other, maximum error pinned at
	\SI{2.95}{\metre} --- the termination bound of \cref{def:crash} --- and the
	\texttt{circle}, \texttt{fig8} and \texttt{square} columns agreeing to within
	\SI{0.03}{\metre}. When the path no longer matters, the experiment is measuring
	the wind.

	A contraction was therefore applied to every multiplicative band, toward its
	nominal value $1$:
	\begin{equation}
		(\ell,h)\ \longmapsto\ \big(1+(\ell-1)\varphi,\ 1+(h-1)\varphi\big),
		\qquad \varphi=0.15,
		\label{eq:moderate}
	\end{equation}
	with the wind band replaced by an absolute mild one. Contracting rather than
	replacing preserves the \emph{shape} of the suite --- which channel is wide,
	which is narrow, which is asymmetric --- so the moderated suite is a strict
	subset of the original and results under it remain statements about the same
	disturbance structure. A guard asserts containment and narrowing channel by
	channel and raises naming the offending key, so that a typo cannot silently
	enter an evaluation.

	\begin{remark}[Honest limitation]
		The moderated suites are disjoint on drag, rotor lag and wind only. They
		\emph{overlap} on mass, thrust, rate gains and inertia, because the raw OOD
		bands on those channels were already supersets of the raw in-distribution
		bands. The claim of full disjointness that appears in the original
		specification is therefore too strong; the notebooks print the per-channel
		disjointness table at run time so the claim cannot drift again.
	\end{remark}

	\realfig{58mm}{\figSaturationGuard}
	{The saturation guard. Fraction of samples with the collective on its box, by
		suite and controller. Any cell above roughly \SI{40}{\percent} is reporting the
		actuator rather than the controller, and is flagged in the comparison of
		\cref{sec:res-comparison} rather than quietly averaged in.}
	{fig:satguard}

	%-----------------------------------------------------------------------
	\section{Recorded signals and derived quantities}
	%-----------------------------------------------------------------------

	\begin{table}[htbp]
		\centering
		\caption{Signals recorded per control step, per vehicle.}
		\label{tab:signals}
		\small
		\begin{tabularx}{\linewidth}{@{}l c X@{}}
			\toprule
			\textbf{Signal} & \textbf{Dim.} & \textbf{Role} \\
			\midrule
			$\plantstate$ & 23 & Full plant state (ground truth) \\
			$\vec o$ & 40 (+6) & Policy observation, \cref{eq:obs} \\
			$\err$ & 9 & Tracking error, \cref{eq:err} \\
			$\inp$ & 4 & CTBR command applied \\
			PWM & 4 & Normalised motor commands from \cref{eq:mixer} \\
			$\dres$ & 6 & Model residual, \cref{eq:dres} \\
			$\wrench$ & 6 & External wrench, \crefrange{eq:Fext}{eq:tauext} \\
			$\vec v_w$ & 3 & Wind realisation, \cref{eq:wind} \\
			$r$, crash & 2 & Reward and termination flag \\
			\bottomrule
		\end{tabularx}
	\end{table}

	\paragraph{The predictor window.} The residual predictor consumes a rolling
	window of $H=\vH$ frames at stride $1$, spanning
	$\vH\times\SI{0.02}{\second}=\SI{1.28}{\second}$, matching the window length of
	the source work. Each frame is $\vFrame$-dimensional:
	\begin{equation}
		\vec x_t=\Big[\underbrace{\vec p-\vec p_\refsym}_{3},\ \underbrace{\vecop(\mat R)}_{9},\
		\underbrace{\vec v}_{3},\ \underbrace{\bodyrate}_{3}\ \Big\|\
		\underbrace{\inp_{t-1}-\inp_\refsym}_{4}\ \Big\|\ \underbrace{\mathrm{PWM}_t}_{4}\Big].
		\label{eq:rdpframe}
	\end{equation}

	\begin{keybox}
		\keyhead{Why the motor commands are in the frame}
		With an integrating inner rate loop, a \emph{constant} external moment is driven
		out of the rate error in steady state --- measured at
		\SI{5.5e-4}{\radian\per\second} after \SI{5}{\second} --- while the mixer output
		holds a spread of $\num{0.16}$. Formally, as $\bodyrate\to\vec 0$ and
		$\dot{\bodyrate}\to\vec 0$, \cref{eq:tauext} reduces to
		$\vec\tau_\mathrm{ext}\to-\mat M_\tau^\mathrm{nom}K_T\vec\Omega^2$, which is built
		\emph{entirely} from the motor commands. Thrust and body-rate setpoints alone
		therefore cannot resolve a standing moment; the actuator commands can. This is
		the single most important design decision in the predictor, and it is also the
		reason the PX4 topic \texttt{actuator\_motors} must be published before any
		off-centre-payload dataset is recorded.
	\end{keybox}

	\paragraph{Prediction smoothing.} The predictor's output is averaged over a
	rolling buffer of $\vSmooth$ frames before it reaches the controller, following
	the source work. The buffer is a fixed-shape ring, not a growing list: a list
	that fills from one to $\vSmooth$ entries changes the shape of the reduction on
	every step, and an array compiler that specialises on shape then recompiles once
	per step. Measured, that cost \SI{68}{\milli\second} per step while the buffer
	filled against \SI{0.29}{\milli\second} once saturated, and since the latency
	probe runs only $30$ steps from a fresh episode the entire probe landed inside
	the warm-up --- reporting the adaptive controller at \SI{45}{\milli\second}
	median instead of \SI{4}{\milli\second}. The ring buffer is bit-identical to the
	growing list and is pinned by a test.

	%-----------------------------------------------------------------------
	\section{Two corrections to the experimental design}
	%-----------------------------------------------------------------------

	\begin{description}[leftmargin=*,style=nextline]
		\item[Saturation contaminating the training distribution.] Wrench randomisation
		at $0.32\,mg=\SI{6.478}{\newton}$ produced \SIrange{22}{40}{\percent}
		collective saturation for the whole of training. A policy pinned against its
		input box visits a state distribution that is not representative of the
		deployed one. The magnitude was reduced to $0.10\,mg=\SI{2.024}{\newton}$ and
		the saturation fraction became a gate rather than a diagnostic
		(\cref{tab:res-trainhealth}).
		\item[A displacement scale that had to be matched to the terminal set.] The
		\LLTC{} candidate displacement is now \SI{0.35}{\metre} against a measured
		terminal-set reach of \SI{0.5649}{\metre}, a ratio of \num{0.62}, so candidates
		lie inside the region where the local quadratic model holds. The earlier
		configuration had this ratio above one and was documented as ``matched'',
		which it was not.
	\end{description}

	%=======================================================================
	\chapter{Methodology and modelling}
	\label{ch:method}
	%=======================================================================

	The methodology is a sequence of four increments. Each adds exactly one
	mechanism, keeps everything else fixed, and is evaluated against the increment
	before it.

	\begin{figure}[htbp]\centering
		\resizebox{0.98\linewidth}{!}{%
\begin{tikzpicture}[
			box/.style={draw=RuleGrey, fill=PanelGrey, rounded corners=2.5pt,
				line width=0.7pt, align=center, text width=26mm,
				minimum height=12mm, font=\scriptsize},
			lrn/.style={box, fill=AccentBlue!12, draw=AccentBlue, line width=1.1pt},
			trn/.style={box, fill=white, draw=TextGrey, text=TextGrey,
				dash pattern=on 2.4pt off 1.8pt},
			arr/.style={-{Stealth[length=2.4mm]}, line width=0.85pt, color=NavyDeep},
			darr/.style={-{Stealth[length=2.4mm]}, line width=0.7pt, color=TextGrey,
				dash pattern=on 2.4pt off 1.8pt},
			lb/.style={font=\tiny, color=TextGrey, inner sep=1.6pt, fill=white}]

			%---- tier 1: the online loop, 50 Hz, on the vehicle ---------------
			\node[box]                       (err)   at (0,0)    {error coordinates\\$\err\in\R^{16}$};
			\node[lrn, right=9mm of err]     (cm)                {\textbf{cost map} $\theta$\\$(\mat S_k,\vec c_k)$};
			\node[box, right=9mm of cm]      (mpc)               {differentiable iLQR\\\cref{alg:ilqr}};
			\node[box, right=9mm of mpc]     (inner)             {PX4 rate loop\\$+$ mixer};
			\node[box, right=9mm of inner]   (plant)             {plant\\$\plantstate\in\R^{23}$};

			%---- tier 2: the residual path -----------------------------------
			\node[box, below=17mm of mpc, text width=22mm] (bd)  {$\mat B_d$ bridge\\\cref{eq:convert}};
			\node[lrn, below=17mm of inner]                (rdp) {\textbf{RDP} $\psi$\\$\vH\times\vFrame$ window};

			%---- tier 3: training only ---------------------------------------
			\node[trn, below=36mm of err] (obs) {observation $\vec o\in\R^{47}$};
			\node[trn, right=9mm of obs]  (ppo) {critic $V_\phi$, PPO\\\cref{eq:ppo}};
			\begin{scope}[on background layer]
				\node[draw=TextGrey, dash pattern=on 2.4pt off 1.8pt, rounded corners=3pt,
				fill=black!3, line width=0.6pt, inner sep=3mm, fit=(obs)(ppo)] (tb) {};
			\end{scope}
			\node[anchor=south west, font=\tiny\itshape, color=TextGrey]
			at ([xshift=6mm,yshift=1.6mm]tb.north west)
			{training only --- never on the vehicle};

			%---- online flow -------------------------------------------------
			\draw[arr] (err)   -- (cm);
			\draw[arr] (cm)    -- (mpc);
			\draw[arr] (mpc)   -- node[lb, above=0.5mm] {CTBR $\inp$} (inner);
			\draw[arr] (inner) -- (plant);

			%---- state feedback, closed above the top row --------------------
			\draw[arr] (plant.north) -- ++(0,9mm)
			-| node[lb, pos=0.22, above] {$\ctrlstate$ and the $\vec x_\refsym$ preview} (err.north);

			%---- residual loop: history down, wrench left, residual up -------
			\draw[arr] (plant.south) |- node[lb, pos=0.55, above] {state, actuator history} (rdp.east);
			\draw[arr] (rdp) -- node[lb, above] {$\hat{\wrench}$} (bd);
			\draw[arr] (bd)  -- node[lb, right, pos=0.5] {$\dres$} (mpc);

			%---- training flow -----------------------------------------------
			\draw[darr] (err.south) -- (obs.north);
			\draw[darr] (obs) -- (ppo);
			\draw[darr] (ppo.north) -- node[lb, right, pos=0.55] {$\nabla_\theta$} (cm.south);

			%---- legend -------------------------------------------------------
			\node[anchor=north west, font=\scriptsize, color=TextGrey,
			inner sep=0pt] at ([yshift=-7mm]tb.south west)
			{\tikz[baseline=-0.5ex]\node[lrn, text width=0pt, minimum height=2.8mm,
				minimum width=5mm, inner sep=0pt]{};\ learned \qquad
				\tikz[baseline=-0.5ex]\node[box, text width=0pt, minimum height=2.8mm,
					minimum width=5mm, inner sep=0pt]{};\ fixed structure \qquad
				\tikz[baseline=-0.5ex]\node[trn, text width=0pt, minimum height=2.8mm,
					minimum width=5mm, inner sep=0pt]{};\ training only};
		\end{tikzpicture}}
		\caption[Architecture of the adaptive actor--critic MPC]{Architecture of the
			adaptive actor--critic MPC. Exactly two blocks are learned. The cost map
			$\theta$ supplies $(\mat S_k,\vec c_k)$ to the \emph{objective} of the
			optimal control problem \cref{eq:ocp}; the residual predictor $\psi$
			supplies $\dres$ to its \emph{dynamics}. Everything else is fixed
			structure. That the two enter at different places is the subject of
			\cref{sec:decoupled} and is what lets either be trained, replaced or
			switched off without touching the other --- and it is what
			\cref{tab:variants} turns into an experiment. The $\mat B_d$ bridge is not
			decorative: $\psi$ emits a physical wrench in \si{\newton} and
			\si{\newton\metre}, the dynamics take a model residual in
			\si{\metre\per\square\second} and \si{\radian\per\square\second}, and the
			two differ by $m^{-1}$ and $\mat J^{-1}$ (\cref{prop:convert}). Supplying
			one where the other is expected is dimensionally wrong and numerically
			silent, which is why it survived into an earlier revision. The dashed
			block exists only while training; the vehicle carries the solid path
			alone, at \SI{50}{\hertz}.}
		\label{fig:architecture}
	\end{figure}

	%-----------------------------------------------------------------------
	\section{Increment 0: baselines and the cost of exactness}
	\label{sec:baselines}
	%-----------------------------------------------------------------------

	Three baselines are carried through the whole study. All are scored on the same
	$\mat Q$, $\mat R$, the same CTBR action space, the same actuator limits and the
	same reference preview, so that any difference is attributable to the controller
	and not to its information.

	\begin{description}[leftmargin=*,style=nextline]
		\item[LQR with reference feed-forward.] One offline Riccati solve, a $9\times4$
		gain applied online. The feed-forward is exact rather than linearised: the
		thrust required to realise $\vec a_\refsym$ is
		$T=m\|\vec a_\refsym+g\vec e_3\|$ and, because of the affine-then-quadratic map
		\cref{eq:thrustmap}, the collective is
		$c=\big(\sqrt{T/T_{\max}}-r\big)/(1-r)$. The body-rate set-point is the
		reference's own angular velocity, from centred differences of $\quat_\refsym$.
		Without this term a pure state-feedback law lags a moving reference for the
		same reason a frozen-reference MPC does, and the comparison would measure the
		missing term rather than the controller. A \texttt{no-feed-forward} row is
		retained precisely to put that difference on the record.
		\item[NMPC with horizon $N$.] \Cref{eq:ocp} with
		$\mat S_k=\blkdiag(\mat Q,\mat R)$ fixed, $\vec c_k=\vec 0$, and
		$\mat P_\mathrm{term}$ the Riccati matrix of the hover linearisation. Swept
		over $N\in\{1,2,3,5,10,20\}$.
		\item[PID.] A cascaded hand-tuned controller. \textbf{It was removed from the
			main comparison} because it shares no design object with anything else --- a
		different structure, nine hand gains, a different tuning budget --- so its
		column measured how well someone tuned a cascade rather than what a controller
		class can do. It is retained in the ROS\,2 experiments of
		\cref{sec:ros-experiments}, where it is the realistic incumbent.
	\end{description}

	\begin{keybox}
		\keyhead{The measurement that motivates everything after this}
		Increment~0 sweeps $\mat Q$ and $\mat R$ and reports the spread in closed-loop
		tracking error between the best and worst admissible choice
		(\cref{tab:res-qr}). If that spread is large, then a substantial part of what
		is usually called ``controller performance'' is really weight tuning, and
		automating the weights is worth doing.
	\end{keybox}

	%-----------------------------------------------------------------------
	\section{Increment 1: learned local time-varying terminal cost (\texorpdfstring{\LLTC}{LLTC})}
	\label{sec:lltc}
	%-----------------------------------------------------------------------

	\paragraph{What Increment~0 left unsolved.} The horizon sweep shows accuracy
	improving with $N$ and p95 latency growing with it, so the horizon is a
	straight trade of compute for optimality. Shortening it to $N{=}1$ settles the
	compute side by construction. The question this increment asks is the narrow
	one: can the \emph{terminal} cost alone carry what the removed stages were
	providing? It is the minimal form of the idea --- a single matrix at the end of
	a one-step problem --- and if it works, nothing more elaborate is needed.

	\subsection{Construction}

	The terminal cost is parametrised so that it is positive definite by
	construction, which keeps \cref{eq:ocp} a well-posed quadratic problem for any
	value of the parameters:
	\begin{equation}
		\mat P_\theta(\err_0)=\mat L_\theta(\err_0)\,\mat L_\theta(\err_0)\transp
		+\varepsilon\mat I,
		\qquad \varepsilon=10^{-3},
		\label{eq:lltc}
	\end{equation}
	with $\mat L_\theta$ lower triangular, its $45$ free entries emitted by a
	multilayer perceptron from the current error.

	\begin{keybox}
		\keyhead{This terminal cost is time-varying, and that is the whole point}
		Classical terminal-cost NMPC carries a \emph{constant} $\mat P$ --- one
		matrix, computed offline, usually the Riccati solution of the linearisation
		about a single operating point, and the same at every control step for the
		life of the controller. \Cref{eq:lltc} is not that. $\mat P_\theta$ is a
		\emph{function of the current error} $\err_0$, re-evaluated at every one of
		the \SI{50}{\hertz} control steps, so along a trajectory it traces out a
		\textbf{time-varying} terminal cost $\mat P_\theta(\err_0(t))$: the
		controller carries a different quadratic tail at every instant, chosen by
		the network from where the vehicle actually is.

		Two consequences follow, and they pull in opposite directions. In its
		favour, a constant $\mat P$ can only be correct near the point it was
		linearised about, whereas a state-dependent one can in principle be correct
		everywhere it has seen data --- which is what makes it worth a network at
		all. Against it, the standard terminal-cost argument (a control-invariant
		terminal set and a $\mat P$ that is a local Lyapunov function on it) is an
		argument about a \emph{fixed} $\mat P$. Letting it vary with the state
		forfeits that argument unless the variation is itself constrained, and
		nothing in \cref{eq:lltc} constrains it beyond positive definiteness.
		\Cref{sec:res-lltc-analysis} reports what that costs in practice: the fit
		is near-exact and the closed loop is not, because ``everywhere it has seen
		data'' turns out to be a much smaller region than the closed loop visits.
	\end{keybox}

	\subsection{Training objective}

	Candidate states $\err^{(i)}$ are sampled around the reference; for each, a
	long-horizon solve provides the realised cost-to-go
	$V_1^{(i)}=J^{(i)}-\ell_0^{(i)}$, and $\theta$ minimises
	\begin{equation}
		\mathcal L_\mathrm{LLTC}(\theta)=\frac1n\sum_{i=1}^{n}
		\Big(\tfrac12\,\err^{(i)\top}\mat P_\theta(\err^{(i)})\,\err^{(i)}-V_1^{(i)}\Big)^{2},
		\label{eq:lltcloss}
	\end{equation}
	restricted to candidates whose cost-to-go lies inside the terminal set.

	\subsection{Rationale and expected failure modes}

	If \cref{eq:lltcloss} is minimised well, then by the standard argument a
	one-step problem with terminal cost equal to the true cost-to-go reproduces the
	infinite-horizon optimum~\cite{lowrey2019polo}. Two failure modes are anticipated
	and both are tested for:

	\begin{enumerate}[leftmargin=*]
		\item \textbf{Insufficient acceptance or an unidentified regression.} A quality
		gate reports the coefficient of determination $R^2$ of \cref{eq:lltcloss} and
		flags $R^2\le0$ --- i.e.\ worse than predicting the mean of $V_1$ --- as a
		negative result rather than a tuning detail.
		\item \textbf{Locality.} $\mat P_\theta$ is fitted where the data is. Out of
		distribution it carries no positive-definiteness or dissipativity guarantee
		and no reason to remain a valid cost-to-go. An OOD sweep tests this directly.
	\end{enumerate}

	\begin{remark}
		\LLTC{} leaves the stage cost --- $N_e+N_u=20$ hand-chosen weights --- entirely
		untouched. The tuning criticism of \cref{sec:baselines} therefore still applies
		in full, which is what motivates Increment~2.
	\end{remark}

	%-----------------------------------------------------------------------
	\section{Increment 2: learned cost map (AC-MPC)}
	\label{sec:acmpc}
	%-----------------------------------------------------------------------

	\paragraph{What Increment~1 left unsolved.} A terminal cost acts at one point of
	the horizon and is fitted offline against a cost-to-go that was itself produced
	by the hand-chosen weights. It therefore inherits those weights rather than
	replacing them, and leaves untouched the third liability of
	\cref{sec:argument}: that a human still picks $\mat Q$ and $\mat R$, and that
	choice moves closed-loop accuracy by a factor \num{2.35} over the admissible
	grid (\cref{tab:res-qr}), and \vQRratio{} over the wider grid of
	\cref{tab:res-cost-spread}. Learning the cost at \emph{every} stage subsumes the
	terminal case and removes the hand-chosen weights entirely --- the
	\texttt{tuned} column of the ledger goes to zero, which is the point.

	\subsection{Construction}

	A trunk network maps $\vec o$ to the full stage cost of \cref{eq:ocp}. Three
	parametrisations are compared, all positive semi-definite by construction. Both
	the quadratic term $\mat S_k$ \emph{and} the linear term $\vec c_k$ are emitted:
	an earlier version emitted only $\mat S_k$ and held $\vec c_k=\vec 0$, which
	discards exactly the part of \cref{eq:groszanon} that absorbs a \emph{bias} in
	the model, and is defect~B\nbd 2.

	\begin{table}[htbp]
		\centering
		\caption{Cost-matrix parametrisations. $Q_\mathrm{LO}$, $Q_\mathrm{HI}$ bound
			the admissible spectrum so the solve stays conditioned. Free parameters per
			stage include the $13$ entries of $\vec c_k$.}
		\label{tab:reps}
		\small
		\begin{tabularx}{\linewidth}{@{}l c X@{}}
			\toprule
			\textbf{\texttt{rep}} & \textbf{Free params / stage} & \textbf{Construction} \\
			\midrule
			\texttt{diag} & 26 & $\mat S=\diag\!\big(Q_\mathrm{LO}+(Q_\mathrm{HI}-Q_\mathrm{LO})\,\sigma(\vec z)\big)$ \\
			\texttt{chol} & 104 & $\mat S=\mat A\mat A\transp+Q_\mathrm{LO}\mat I$, $\mat A$ lower triangular, positive diagonal \\
			\texttt{full} & 182 & $\mat S=\mat A\mat A\transp+Q_\mathrm{LO}\mat I$, $\mat A$ unconstrained \\
			\bottomrule
		\end{tabularx}
	\end{table}

	\subsection{Why an initialisation argument, not a capacity argument}

	Romero et al.~\cite{romero2024acmpc} report that the richer forms fail to learn
	within the diagonal's sample budget and attribute it to the enlarged search
	space. This work advances a sharper hypothesis, which is testable and makes a
	different prediction. With the head's weights scaled small and zero bias,
	$\mat A\approx\mat 0$, so $\mat S\approx Q_\mathrm{LO}\mat I$: the quadratic term
	starts \emph{effectively absent} and the solve is dominated by the linear term.
	The diagonal parametrisation, by contrast, starts at the sigmoid mid-range. The
	measured initial spectra make the asymmetry explicit --- \texttt{diag} begins
	with $\mat S$ diagonal at \num{3.162}, \texttt{chol} at \num{0.4905} and
	\texttt{full} at \num{0.01} --- so the three do not start from comparable
	places.

	\textbf{Prediction.} The richer forms should require \emph{longer}, not be
	incapable. \Cref{sec:res-acmpc} reports whether the reproduction supports the
	published claim or this one.

	\subsection{Training}

	Trained by PPO through \cref{alg:ilqr}. The design point studied is the one
	that is a plausible failure mode of differentiating through an optimiser
	rather than of reinforcement learning in general:

	\begin{enumerate}[leftmargin=*]
		\item \textbf{Exploration noise.} Additive Gaussian exploration on a CTBR command
		whose collective is box-constrained does not merely perturb the action: it
		drives the input onto the box, which corrupts the linearisation
		$\mat A_k,\mat B_k$ inside the solve. A differentiable MPC layer is therefore
		expected to be markedly more noise-sensitive than a model-free MLP with the
		same reward, and the exploration standard deviation is swept to test this.
	\end{enumerate}

	\begin{remark}[A trust region is required, not optional]
		Because the policy mean is the \emph{solution of an optimisation problem}, a
		single clipped gradient step can move $\mu_\theta$ far more than the change in
		$\theta$ suggests: a \SI{0.5}{\percent} change in $\mat S$ was measured to move
		the mean by \num{0.32} in normalised command units, taking the sampled KL to
		order $10$. PPO's ratio clip does not bound that, because it constrains the
		ratio and not the mean. A KL trip that reverts the update, together with an
		adaptive learning rate, is therefore part of the training loop rather than a
		refinement of it.
	\end{remark}

	%-----------------------------------------------------------------------
	\section{Increment 3: domain randomisation}
	\label{sec:dr}\label{sec:noise}
	%-----------------------------------------------------------------------

	\paragraph{What Increment~2 left unsolved.} A cost map trained on one plant is a
	cost map \emph{for} that plant. Increment~2 removes the hand-chosen weights but
	says nothing about the first liability of \cref{sec:argument} --- mismatch
	between the plant that is flown and the model that is planned with. Domain
	randomisation is the standard answer, and it comes with a structural cost that
	this increment is built to measure rather than assume: a policy that cannot
	\emph{observe} which member of the randomised family it is flying cannot
	specialise to it. It can only find one compromise acceptable everywhere.
	Generality is bought with specificity, and the currency is the conservatism
	premium of question~2 below.

	Training is repeated over the parameter bands of \cref{tab:dist}, and four
	questions are asked, each with its own experiment.

	\begin{enumerate}[leftmargin=*]
		\item \textbf{Does it buy robustness?} Measured as the inter-quartile range of
		per-episode RMSE across the randomised family, not as the mean.
		\item \textbf{What does it cost?} Every controller is evaluated on the
		\emph{nominal} plant. The gap between a DR policy and a nominal-only policy on
		the nominal plant is the conservatism premium, charged on every flight whether
		or not a disturbance arrives.
		\item \textbf{Which variance moves?} Two different variances exist ---
		seed-to-seed and episode-to-episode --- and they are expected to move in
		opposite directions. Reporting only one is how a robustness claim is
		accidentally manufactured.
		\item \textbf{Does sequential randomisation forget?} Tasks are presented in
		sequence and backward transfer is measured, together with two plasticity
		diagnostics: the fraction of dead trunk units and the fraction of saturated
		cost-map outputs.
	\end{enumerate}

	\begin{keybox}
		\keyhead{The structural limitation that motivates Increment 4}
		A DR policy cannot tell which plant it is on, so it cannot specialise its
		response --- it can only find one compromise acceptable everywhere. That is a
		ceiling on what DR can achieve, whatever the conservatism premium turns out to
		be, and it does not depend on the sign of that premium. So: give the policy
		eyes.
	\end{keybox}

	%-----------------------------------------------------------------------
	\section{Increment 4: adaptive AC-MPC with a residual predictor}
	\label{sec:adaptive}
	%-----------------------------------------------------------------------

	\paragraph{What Increment~3 left unsolved.} If the conservatism premium is the
	price of \emph{not knowing} the plant, then the remedy is not more
	randomisation but an observation. Estimating the disturbance online should
	recover the specificity that Increment~3 averaged away, while keeping the
	generality it bought --- and it closes the chain of \cref{sec:argument}, because
	the residual predictor is the \emph{enhanced model} standing beside the cost map
	as the \emph{enhanced cost}. Whether it actually pays is the question
	\cref{sec:res-adaptive} answers, and the decomposition below is what makes the
	answer attributable rather than a single aggregate number.

	\subsection{Decoupled construction}
	\label{sec:decoupled}

	The adaptive scheme is built in two stages, which is what makes the result
	attributable.

	\begin{enumerate}[leftmargin=*]
		\item \textbf{Train an oracle.} A policy that receives the ground-truth
		disturbance, both in its observation and in its internal model. This
		establishes the ceiling: it is the performance available if estimation were
		perfect.
		\item \textbf{Replace the oracle with an estimator.} A recurrent network
		regresses the disturbance from onboard signals only. The gap between oracle
		and estimator is attributable to estimation; the gap between oracle and the
		disturbance-blind baseline is attributable to the information itself.
	\end{enumerate}

	\begin{keybox}
		\keyhead{A precondition that must be checked first}
		The decomposition is only meaningful if the oracle is \emph{trainable}. If the
		oracle policy spends training saturated against its input box, it never learns a
		disturbance-response manifold to specialise, and the oracle signal is correct but
		useless. \Cref{sec:res-adaptive} reports the gate, which is checked before any
		closed-loop number is read.
	\end{keybox}

	\subsection{The residual dynamics predictor}
	\label{sec:rdp-theory}

	Four temporal encoders are compared on the identical $\vH\times\vFrame$ window
	\cref{eq:rdpframe} and the identical $6$-dimensional target, so the comparison
	isolates the encoder. The recurrent encoders use two layers of width $64$,
	matching the source work.

	The loss is normalised per channel, which is not cosmetic: the target's first
	three components are in \si{\metre\per\square\second} and its last three in
	\si{\radian\per\second}, and an unnormalised squared error fits only the larger
	block.
	\begin{equation}
		\mathcal L_\mathrm{RDP}(\psi)=\frac{1}{6n}\sum_{i=1}^{n}\sum_{j=1}^{6}
		\left(\frac{\hat d_j^{(i)}(\psi)-d_j^{(i)}}{s_j}\right)^{\!2},
		\qquad s_j=\mathrm{std}_j(d).
		\label{eq:rdploss}
	\end{equation}

	\paragraph{Training objective: position hold.} The predictor's training set is
	flown on a \emph{hold} setpoint, not on the tracking references. On a hold, the
	disturbance is the only thing moving the vehicle, so the regression target is
	the residual rather than the residual plus the controller's own tracking
	transient. Evaluation stays on the tracking suites, so the predictor is always
	tested off its training distribution and the generalisation claim is tested
	rather than assumed. The same split is applied in the ROS\,2 workspace, where
	experiment E\nbd 0 generates the training set on a hold and every
	evaluation experiment flies the Lissajous reference.

	\subsection{Three ways to use the estimate}

	\begin{table}[htbp]
		\centering
		\caption{Variants of the adaptive scheme. ``Model'' means the estimate is
			written into $\dres$ inside the solver; ``observation'' means it is
			concatenated to $\vec o$.}
		\label{tab:variants}
		\small
		\begin{tabularx}{\linewidth}{@{}c c c X@{}}
			\toprule
			\textbf{Variant} & \textbf{To model} & \textbf{To observation} & \textbf{Interpretation} \\
			\midrule
			A & no  & yes & The cost map may react to the disturbance, the model stays nominal \\
			B & yes & no  & The model is corrected; the cost map remains disturbance-blind \\
			C & yes & yes & Both \\
			\bottomrule
		\end{tabularx}
	\end{table}

	\begin{cautionbox}
		\textbf{Interface correctness.} Writing into $\dres$ requires
		\cref{def:residual} units. Supplying \cref{def:wrench} instead --- forces in
		\si{\newton} and moments in \si{\newton\metre} --- is dimensionally wrong and,
		because it is wrong \emph{consistently} between training and evaluation,
		produces no error message. \Cref{prop:convert} gives the correct conversion,
		exact in both blocks; the original implementation discarded the moment block
		entirely.
	\end{cautionbox}

	\subsection{Disturbance scenarios}

	\begin{table}[htbp]
		\centering
		\caption{Physical disturbance scenarios, as configured in
			\texttt{adaptive\_core\_jax.py}. Levels are fractions of vehicle mass except
			for the slung load, where the level is the figure-eight period in seconds
			(shorter is harsher). Moderated levels are used for evaluation.}
		\label{tab:scenarios}
		\small
		\begin{tabularx}{\linewidth}{@{}l X l@{}}
			\toprule
			\textbf{Scenario} & \textbf{Physics} & \textbf{Raw levels} \\
			\midrule
			\texttt{central} & Payload at $\vec r_p=(0,0,-0.05)$: pure vertical force, no CG shift & $0,\ 0.075,\ 0.15,\ 0.25$ \\
			\texttt{asym} & Payload at an arm tip, $\vec r_p=(0.174,0.174,-0.05)$: coupled force \emph{and} standing moment & $0,\ 0.01,\ 0.07,\ 0.11$ \\
			\texttt{slung} & Spherical pendulum, $L=\vTether\,\si{\metre}$, tracked on a figure eight & $15,\ 10,\ 5,\ 3$ \si{\second} \\
			\bottomrule
		\end{tabularx}
	\end{table}

	The \texttt{asym} ladder is bounded by a static trim limit rather than chosen
	freely: with the arms and $f_{\max}$ of \cref{tab:vehicle}, the largest
	off-centre load the allocator can trim is $f_{\max}=\num{0.0812}$ of vehicle
	weight. The raw top level $0.11$ exceeds it and the moderated one does not, and
	the notebook prints that comparison before the sweep so an infeasible cell
	cannot be mistaken for a control failure.

	The slung load is integrated as a spherical pendulum coupled through the tether
	tension rather than as a stiff spring, because a spring cable requires a physics
	step well below its natural period and is the usual reason such simulations
	diverge. With $\hat{\vec n}$ the unit tether direction in world coordinates,
	\begin{equation}
		F_\mathrm{ten}=m_\ell\!\left(g\,\hat{\vec n}\cdot\vec e_3+L\|\dot{\hat{\vec n}}\|^2\right),
		\qquad
		\vec F^{\,b}_\mathrm{ext}=\mat R\transp\big(F_\mathrm{ten}\hat{\vec n}\big),
		\label{eq:tether}
	\end{equation}
	with $\hat{\vec n}$ propagated under the constraint $\|\hat{\vec n}\|=1$ and
	$\dot{\hat{\vec n}}\perp\hat{\vec n}$. At $L=\vTether\,\si{\metre}$ --- the arm
	length, so that the tether is physically mountable on this airframe --- the
	pendulum period is $2\pi\sqrt{L/g}=\vTetherT\,\si{\second}$.

	\begin{remark}[Two simplifications to be validated]
		\Cref{eq:tether} produces no moment, which is exact only if the attachment point
		coincides with the centre of gravity; and the pendulum is advanced by explicit
		Euler at $\Delta t_c$ while the plant uses RK4 at $\Delta t_c/10$. Explicit Euler
		on an oscillator injects energy monotonically, so tether energy drift over an
		episode is monitored, and a symplectic variant --- a one-line reordering --- is
		available if the drift exceeds a few percent.
	\end{remark}
	%=======================================================================
	\chapter{Implementation and experimentation}
	\label{ch:impl}
	%=======================================================================

	%-----------------------------------------------------------------------
	\section{Software architecture}
	%-----------------------------------------------------------------------

	The codebase is deliberately split into two halves that never depend on each
	other. The study half must import with no ROS on the path; the robotics half
	must build with neither JAX nor PyTorch. They meet only through a serialised
	model file, which keeps each side testable in isolation.

	\begin{table}[htbp]
		\centering
		\caption{Principal modules, with line counts at the revision this report
			describes.}
		\label{tab:modules}
		\small
		\begin{tabularx}{\linewidth}{@{}l r X@{}}
			\toprule
			\textbf{Module} & \textbf{Lines} & \textbf{Responsibility} \\
			\midrule
			\multicolumn{3}{@{}l}{\itshape Study (pure JAX/Python)} \\
			\texttt{x500\_core\_jax.py} & 2558 & Plant, error coordinates, iLQR layer, cost map, PPO, batched environment, baselines, metrics \\
			\texttt{adaptive\_core\_jax.py} & 759 & Payload/slung scenarios, GRU/LSTM/TCN/CNN predictors \\
			\texttt{study\_prelude.py} & 183 & Scale table, shared configuration, controller factory, reporting helpers \\
			\texttt{study\_moderate.py} & 248 & Moderated suites \cref{eq:moderate}, on-path initialisation, terminal-cost reconstruction \\
			\texttt{viz.py} & 296 & Three-dimensional flight clip renderer \\
			\texttt{export\_estimator.py} & --- & Framework-neutral \texttt{.npz} export with a parity assertion \\
			Notebooks 1--7 & 2696 & One increment per notebook, each writing versioned artefacts \\
			Test suite & 1202 & 86 tests, of which 26 pin audit findings \\
			\midrule
			\multicolumn{3}{@{}l}{\itshape ROS\,2 / PX4 workspace} \\
			\texttt{acmpc\_controller} & --- & CTBR bridge, $B_d$ bridge, glue and action-mapping self-tests \\
			\texttt{rdp\_estimator} & --- & Ring buffer and pure-NumPy inference, no JAX \\
			\texttt{reference\_generator} & --- & Lissajous and position-hold references, feasibility check \\
			\texttt{experiment\_manager} & --- & Sweep plans (134 runs) and config capture \\
			\texttt{disturbance\_manager} & --- & Gazebo payload and wind assets \\
			\texttt{state\_logger}, \texttt{visualization} & --- & Recording and live panels \\
			\bottomrule
		\end{tabularx}
	\end{table}

	\subsection{Why JAX}

	Three properties of the problem dictated the framework choice.

	\begin{enumerate}[leftmargin=*]
		\item \textbf{The environment must be differentiable and batched.}
		\Cref{eq:edyn} is differentiated inside the training loop, and $512$ vehicles
		are advanced as one array program.
		\item \textbf{The solver must be compiled.} \Cref{alg:ilqr} is a fixed-iteration
		loop with a line search; expressed with \texttt{lax.fori\_loop} and
		\texttt{lax.scan} it compiles to a single kernel launch and admits
		\texttt{stop\_gradient} on the non-differentiated prefix.
		\item \textbf{Environments must be hashable to avoid recompilation.} The static
		configuration is a frozen dataclass used as a JIT static argument, so two
		environments with the same configuration share one compiled program.
	\end{enumerate}

	Numerics are in \texttt{float64} throughout. This is unusual for RL and
	deliberate here: the backward Riccati recursion of \cref{alg:ilqr} inverts
	$\mat Q_{uu}$ at every stage, and the cost matrices span several orders of
	magnitude by construction (\cref{tab:reps}).

	\subsection{Numerical parity and export}

	The four predictor encoders are written out explicitly --- the recurrences as
	\texttt{lax.scan} over the exact gate equations with the same gate ordering and
	bias convention, the convolutions as \texttt{conv\_general\_dilated} with the
	reference padding --- rather than taken from a neural network library. This
	keeps a comparison of encoders a comparison of encoders, and lets a parity test
	transplant reference weights and assert that the outputs agree.

	The deployment path is a pure-NumPy re-implementation, and the export refuses to
	write a model whose NumPy forward pass disagrees with the JAX one: measured, all
	five exported encoders agree to $\approx10^{-14}$ against a $10^{-5}$ tolerance.
	A deployment model that has silently diverged from the trained one is the single
	most expensive bug available here, because nothing downstream would notice.

	Two subtleties were necessary for parity and are recorded because they are easy
	to reintroduce:

	\begin{itemize}[leftmargin=*]
		\item \textbf{Clamp gradients.} \texttt{torch.clamp} and \texttt{jnp.clip} agree
		on the value and disagree on the derivative at the boundary. This is not a
		corner case, because the iLQR line search parks the collective exactly on the
		box edge; the measured consequence is a $0.297$ absolute error in $\mat B$,
		enough to change the descent path. A custom JVP pins the convention.
		\item \textbf{Variance convention.} The biased/unbiased choice in standard
		deviations used for observation normalisation must match the reference, and the
		variance must be floored: several observation channels are near-constant within
		a batch --- a frozen preview under position hold, an integral that has not
		moved --- and an unfloored $1/\sqrt{\mathrm{var}}$ amplifies the first sample
		that does move by $\sim10^4$, which took the policy to \texttt{NaN} within
		three iterations.
	\end{itemize}

	%-----------------------------------------------------------------------
	\section{Experimental scale}
	%-----------------------------------------------------------------------

	A single environment variable selects between a wiring-check configuration and
	the configuration the study reports.

	\begin{table}[htbp]
		\centering
		\caption{Scale table, from \texttt{study\_prelude.py}. \texttt{smoke} exists to
			verify that the pipeline executes; it produces undertrained policies and its
			learned-controller numbers are not results.}
		\label{tab:scale}
		\small
		\begin{tabularx}{\linewidth}{@{}l r r r X@{}}
			\toprule
			\textbf{Key} & \textbf{smoke} & \textbf{medium} & \textbf{full} & \textbf{Meaning} \\
			\midrule
			\texttt{n\_env} & 16 & 128 & 512 & Parallel vehicles during training \\
			\texttt{T\_rollout} & 16 & 32 & 64 & PPO steps per iteration \\
			\texttt{minib} & 8 & 16 & 32 & Minibatches per epoch \\
			\texttt{epochs} & 4 & 6 & 10 & PPO epochs per iteration \\
			\texttt{iters} & 6 & 80 & 400 & PPO iterations \\
			\texttt{iters\_sweep} & 4 & 30 & 150 & Iterations per sweep point \\
			\texttt{hid} & 64 & 256 & 1024 & Trunk width \\
			\texttt{ilqr} & 5 & 10 & 25 & iLQR iterations per solve \\
			\texttt{n\_eval} & 8 & 64 & 256 & Vehicles per evaluation \\
			\texttt{T\_eval} & 250 & 600 & 1200 & Steps per evaluation \\
			\texttt{est\_epochs} & 2 & 15 & 60 & Predictor training epochs \\
			\texttt{est\_batch} & 128 & 1024 & 4096 & Predictor batch size \\
			\midrule
			\multicolumn{5}{@{}l}{\itshape Training work, $\mathtt{n\_env}\times\mathtt{T\_rollout}\times\mathtt{iters}\times\mathtt{ilqr}$} \\
			relative & $1\times$ & $427\times$ & $42\,667\times$ & \\
			\bottomrule
		\end{tabularx}
	\end{table}

	\begin{cautionbox}
		\textbf{The full-scale campaign has not been run.} Every quantitative cell in
		\cref{ch:results} is therefore left blank. The pipeline has been executed end
		to end at \texttt{smoke} scale, which verifies that it runs and that its gates
		fire, but \texttt{smoke} produces undertrained policies by construction and its
		learned-controller numbers are not results. Scaling from the measured
		\texttt{smoke} timing, a single full-scale training run is of order
		\SIrange{90}{190}{\hour} and Notebook~3 alone trains roughly twenty-five
		policies, so a full campaign is weeks of wall-clock on one workstation;
		\texttt{medium} is the sensible intermediate rung.
	\end{cautionbox}

	\paragraph{Where the wall-clock goes.} Every stage of every notebook is
	timed, and the per-stage totals are written to the run directory as
	\texttt{timing\_summary.csv}. For the
	domain-randomisation notebook at \texttt{medium} scale --- the most
	expensive of the seven, and the only one whose cost is dominated by repeated
	training rather than by a single run --- the total is \SI{2.03}{\hour},
	distributed as
	\SI{73.8}{\percent} policy training (three arms, \SIrange{15}{41}{\minute}
	each), \SI{20.0}{\percent} the four-stage sequential curriculum, and
	\SI{6.2}{\percent} all evaluation and axis rollouts together. The ratio is
	what matters rather than the absolute figures, which scale with the machine:
	\textbf{training dominates by an order of magnitude}, so the only effective
	economy is to train fewer arms, which is what the reuse below does.


	\begin{remark}[Where the domain-randomisation cost comes from]
		The seed-variance study retrains arms from scratch, and its first seed
		duplicates a policy already trained by the headline stage. Reusing that policy
		avoids retraining three arms. The trade is that the axis sweep's binning
		changes from a delta at the bin centre to a random sample within the bin, which
		is stated explicitly rather than absorbed silently.
	\end{remark}

	%-----------------------------------------------------------------------
	\section{Evaluation protocol}
	%-----------------------------------------------------------------------

	\subsection{Conditions}

	Three conditions, all using the moderated suites of \cref{tab:dist}:

	\begin{description}[nosep,leftmargin=*,style=nextline]
		\item[S1 --- nominal.] No wind, nominal parameters. Measures raw tracking.
		\item[S2 --- in distribution.] The training suite, moderated. Measures
		robustness inside the family the policy saw.
		\item[S3 --- out of distribution.] Disjoint on drag, rotor lag and wind;
		overlapping on the remaining channels, as \cref{sec:moderation} records.
	\end{description}

	Each condition is evaluated on \texttt{circle}, \texttt{fig8} and
	\texttt{square}, over fixed evaluation seeds recorded in a versioned artefact so
	that a rerun is comparable.

	\subsection{Metrics, and why each was chosen}

	\begin{table}[htbp]
		\centering
		\caption{Evaluation metrics. All are computed per vehicle and reported as the
			median with the inter-quartile range across vehicles, because the distributions
			are right-skewed by occasional terminations and a mean is therefore
			uninformative.}
		\label{tab:metrics}
		\small
		\begin{tabularx}{\linewidth}{@{}l l X@{}}
			\toprule
			\textbf{Metric} & \textbf{Definition} & \textbf{Why} \\
			\midrule
			RMSE & $\sqrt{\frac1T\sum_t\|\vec p_t-\vec p_{\refsym,t}\|^2}$ & The task \\
			IQR(RMSE) & Inter-quartile range across vehicles & Robustness is a statement about spread, not mean \\
			$e_{\max}$ & $\max_t\|\vec p_t-\vec p_{\refsym,t}\|$ & Worst case; saturates at the bound of \cref{def:crash} \\
			tilt & $\max_t\angle(\vec z_b,\vec e_3)$ & Attitude excursion; a proxy for aggressiveness \\
			effort & $\sum_t\|\inp_t-\inp_\refsym\|^2$ & Energy and actuator wear \\
			smooth & $\sum_t\|\inp_t-\inp_{t-1}\|^2$ & Jerk; what an operator perceives as quality \\
			sat & Fraction of samples with the collective on the box & \textbf{A guard}: any row with high saturation is reporting the actuator, not the controller \\
			crash & Fraction of episodes terminating per \cref{def:crash} & Feasibility \\
			$t_\mathrm{solve}$ & Median and p95 \emph{single-vehicle} solve latency & Admissibility against \SI{20}{\milli\second} \\
			\midrule
			$R^2_j$ & $1-\mathrm{SS}_{\mathrm{res},j}/\mathrm{SS}_{\mathrm{tot},j}$ per channel & Predictor quality; break-even at $0$ \\
			NRMSE & Channel-normalised RMSE of \cref{eq:rdploss} & Predictor error in comparable units \\
			\bottomrule
		\end{tabularx}
	\end{table}

	\begin{cautionbox}
		\textbf{One latency, not two.} \texttt{ms\_per\_step} divided by the batch size
		is \emph{throughput} over many vehicles and has nothing to do with what one
		aircraft experiences. Only a batch-of-one measurement may be compared against
		the control period. Mixing the two produced an apparent $200\times$ speed-up for
		the learned controllers in an earlier revision; see defect~B\nbd 1 in
		\cref{app:audit}.
	\end{cautionbox}

	\subsection{Statistical treatment}

	Two distinct variances are reported separately and never pooled:

	\begin{itemize}[nosep,leftmargin=*]
		\item \textbf{Episode variance} --- spread across vehicles within one trained
		policy. This is what robustness claims are about.
		\item \textbf{Seed variance} --- spread across independently trained policies.
		This is what ``method A beats method B'' claims are about, and a difference
		smaller than the seed spread is not a result.
	\end{itemize}

	Comparisons between trained methods are therefore made across $\ge3$ training
	seeds, and any claimed improvement is reported alongside the measured seed
	spread (\cref{tab:res-seeds}) so the reader can judge it.

	%-----------------------------------------------------------------------
	\section{Flight visualisation}
	%-----------------------------------------------------------------------

	Clips are rendered for each controller on a common path under a common
	condition. Three properties were necessary to make them comparable and each cost
	a bug to discover:

	\begin{enumerate}[leftmargin=*]
		\item \textbf{The vehicle starts on the reference.} The environment seeds each
		episode with a $\approx\SI{0.5}{\metre}$ offset, so an unmodified clip opens
		with a recovery transient rather than tracking. A \emph{fixed pilot} flies the
		vehicle onto the path and the controller under test takes over at frame zero.
		Letting each controller settle itself is circular: a controller that cannot
		hold the path drives the error \emph{up} during the warm-up, so it can never be
		made to start on the path. The clips are asserted to open in an identical state
		to $10^{-6}$.
		\item \textbf{The episode outlasts the clip.} Recording length and episode length
		are coupled by construction.
		\item \textbf{The whole reference is drawn.} The clip must span at least one full
		path period, and the dashed reference must be the complete closed curve, not
		the fragment that was flown. Since $\omega$ depends on the sampled radius
		through \cref{eq:omegacap}, the clip length is derived from
		$2\pi/\omega$ rather than fixed in steps.
	\end{enumerate}

	\realfigtwo{58mm}{\figContactSheet}{\figGroundTracks}
	{Flight clips, one per controller, on a common path and condition, all opening in
		an identical on-path state. Drawn from the run manifest rather than by
		enumerating the output directory.}
	{Contact sheet: one frame from each clip.}
	{Ground tracks against the complete reference curve.}
	{fig:contact}

	%-----------------------------------------------------------------------
	\section{Sim-to-real: the PX4/Gazebo workspace}
	\label{sec:ros}
	%-----------------------------------------------------------------------

	\subsection{Cascade and interfaces}

	\begin{center}
		\begin{tikzpicture}[
			node distance=5mm and 7mm,
			box/.style={draw=RuleGrey, fill=PanelGrey, rounded corners=2pt, line width=0.7pt,
				align=center, text width=30mm, minimum height=11mm, font=\footnotesize},
			arr/.style={-{Stealth[length=2.2mm]}, line width=0.7pt, color=NavyDeep}]
			\node[box] (sp) {setpoint\\\scriptsize\texttt{reference\_generator}};
			\node[box, right=of sp] (outer) {\textbf{outer loop}\\\scriptsize 50\,Hz, offboard\\\scriptsize PID / NMPC / AC-MPC};
			\node[box, right=of outer] (inner) {\textbf{PX4 inner loop}\\\scriptsize $\approx$1\,kHz, onboard\\\scriptsize rate PID + allocator};
			\node[box, right=of inner] (gz) {\textbf{Gazebo}\\\scriptsize 1\,kHz physics\\\scriptsize $F=K_T\Omega^2$};
			\draw[arr] (sp) -- (outer);
			\draw[arr] (outer) -- node[above, font=\scriptsize, color=TextGrey] {CTBR} (inner);
			\draw[arr] (inner) -- node[above, font=\scriptsize, color=TextGrey] {motors} (gz);
			\draw[arr] (gz.south) -- ++(0,-6mm) -| node[pos=0.25, below, font=\scriptsize, color=TextGrey]
			{odometry, IMU, actuator commands} (outer.south);
		\end{tikzpicture}
	\end{center}

	The outer loop never commands motors: it commands collective thrust and body
	rates, and PX4 owns everything below. That modularity is what allows the
	learned controller to replace one block without touching anything else.

	\subsection{Frame conversions and the tests that protect them}

	\begin{table}[htbp]
		\centering
		\caption{Frame conventions at the PX4 boundary.}
		\label{tab:frames}
		\small
		\begin{tabularx}{\linewidth}{@{}l l X@{}}
			\toprule
			\textbf{Quantity} & \textbf{Frame} & \textbf{Handling} \\
			\midrule
			\texttt{vehicle\_odometry.position/velocity} & NED & converted on entry \\
			\texttt{vehicle\_angular\_velocity.xyz} & FRD & converted on entry \\
			\texttt{vehicle\_odometry.q} & body FRD $\to$ world NED & $\quat_\mathrm{ENU/FLU}=\quat_{\mathrm{NED}\to\mathrm{ENU}}\otimes\quat_\mathrm{PX4}\otimes\quat_{\mathrm{FRD}\to\mathrm{FLU}}$ \\
			everything inside the package & ENU/FLU & --- \\
			\texttt{VehicleRatesSetpoint} & FRD & $\mathrm{roll}=\omega_x$, $\mathrm{pitch}=-\omega_y$, $\mathrm{yaw}=-\omega_z$ \\
			\texttt{thrust\_body} & FRD, down-positive & $[0,0,-c]$ \\
			\bottomrule
		\end{tabularx}
	\end{table}

	The mapping is asserted numerically rather than trusted: the composed rotation
	agrees with the direct construction to $10^{-15}$ over $500$ random attitudes,
	the tilt is independent of heading to $10^{-17}$, and the CTBR exit mapping
	passes $9/9$ checks. A separate consistency checker asserts that every physical
	constant shared between the study and the workspace agrees, and that every
	exported predictor carries the expected window length and channel ordering.

	\subsection{The thrust curve is not linear, and calibration matters}

	Two nonlinearities compose between the commanded collective $c$ and the force
	produced: the PX4 allocator applies $c=(1-f)y+fy^2$ with $f$ a model factor, and
	the simulated motor gives $F=F_{\max}y^2$ on top of the idle floor of
	\cref{eq:thrustmap}. On this airframe $f=0$, so $c=y$ and the hover command is
	$u_\mathrm{hover}=\vUhover$.

	\begin{cautionbox}
		\textbf{A dataset-wide bias caused by one constant.} An earlier version used an
		uncalibrated $T_{\max}$ in the wrench ground truth while the thrust model used a
		calibrated one. Every recorded $\vec F_\mathrm{ext}$ was therefore biased in the
		same direction across the entire training set --- a systematic error invisible
		in any per-sample check. All modules now derive the constant from one place and
		a test asserts agreement.
	\end{cautionbox}

	\subsection{Experiments for the residual predictor}
	\label{sec:ros-experiments}

	\begin{table}[htbp]
		\centering
		\caption{ROS\,2 / Gazebo experiment matrix, as planned by
			\texttt{experiment\_manager}. E\nbd 0 generates the predictor's
			training set on a \emph{position hold}; every evaluation experiment flies the
			tracking reference, so the predictor is always tested off its training
			distribution. All 134 planned runs record their reference mode explicitly.}
		\label{tab:rosexp}
		\small
		\begin{tabularx}{\linewidth}{@{}l l X r@{}}
			\toprule
			\textbf{Family} & \textbf{Reference} & \textbf{Purpose} & \textbf{Runs} \\
			\midrule
			E-0 & hold & RDP training set: S0,S1,S2,S3,S5 $\times$ 3 seeds & 15 \\
			E-A & tracking & Scenario matrix S0--S6 $\times$ C0--C3 $\times$ 3 seeds, plus PID and NMPC incumbents on S0/S1 & 96 \\
			E-B & tracking & Frequency response and chirp from S4 & 7 \\
			E-C & tracking & MPC horizon interaction, $N\in\{10,20,30,40\}$ & 12 \\
			E-D & hold (train) / tracking (score) & History length $H\in\{16,32,\vH,128\}$ & 4 \\
			\bottomrule
		\end{tabularx}
	\end{table}

	For a payload of mass fraction $f$ at body position $\vec r_p$, in steady hover,
	\begin{equation}
		F_z^\mathrm{ext}\approx-f\,m\,g,
		\qquad
		\vec\tau^\mathrm{ext}\approx\vec r_p\times\big(-f\,m\,g\,\vec e_3^{\,b}\big),
		\label{eq:analytic}
	\end{equation}
	a physical prediction with no free parameters, drawn on the live plot as a
	correctness check.

	\paragraph{Protocol.} Level $0$ is flown first for all controllers. If the
	measured wrench is not near zero at level $0$ --- acceptance
	$|F_z|<\SI{0.3}{\newton}$ and $\|\vec\tau\|<\SI{0.05}{\newton\metre}$ --- the
	ground-truth computation is wrong and nothing downstream is meaningful, so the
	campaign stops. Levels are then swept in ascending order with three repeats per
	cell at different simulator seeds, reporting the median and the spread.
	Publication of the motor-command topic is checked before every recording, for
	the reason given in \cref{sec:rdp-theory}.

	\realfig{58mm}{\figWrenchSeries}
	{Wrench ground truth against the predictor over a disturbed episode, one panel
		per channel. The analytic steady-state value of \cref{eq:analytic} is the
		horizontal guide; it has no free parameters, which makes it the most
		convincing available correctness check on the ground-truth computation.}
	{fig:rospanel}

	%=======================================================================
	\chapter{Results and evaluation}
	\label{ch:results}
	%=======================================================================

	\begin{cautionbox}
		\textbf{Reading convention, and the scale these numbers come from.} Every
		quantitative cell in this chapter is a measured value, written into this
		document from the artefact CSVs by \texttt{study/fill\_report.py}; none is
		estimated and none is left blank. They were produced at \texttt{medium}
		scale (\cref{tab:scale}), not at \texttt{full}. That matters for how far
		they generalise, and it is stated once here rather than attached to every
		number: \texttt{medium} trains for $80$ PPO iterations against $400$, with a
		trunk of $256$ against $1024$. Differences of the same order as the seed
		spread of \cref{tab:res-seeds} are therefore not resolved, and this chapter
		says so where it matters instead of reporting them as effects.

		Numbers quoted in the running text of earlier chapters are constants of the
		platform (\cref{tab:vehicle}), properties of the reference geometry
		(\cref{sec:feasibility}), or engineering diagnostics such as parity
		tolerances --- none is a controller result.

		Figures resolve through the \texttt{FIGURE FILES} block at the head of
		this document. It searches \texttt{report/figures/} first and
		\texttt{artifacts/figures/} second, so re-running the study updates every
		figure in place, and a plot of your own dropped into the first of those
		directories overrides the generated one without any edit to this
		document.
	\end{cautionbox}

	%-----------------------------------------------------------------------
	\section{What the pipeline execution established}
	\label{sec:res-pipeline}
	%-----------------------------------------------------------------------

	Before any controller claim, the apparatus itself has to be trustworthy. The
	following are properties of the code and the platform, not of any trained
	policy, and they hold at every scale.

	\begin{table}[htbp]
		\centering
		\caption{Verification status of the apparatus. These are the checks that must
			pass before a result is readable; none of them is itself a result.}
		\label{tab:res-verification}
		\small
		\begin{tabularx}{\linewidth}{@{}X l@{}}
			\toprule
			\textbf{Check} & \textbf{Status} \\
			\midrule
			Physics self-test: 20 derived quantities against the SDF (T-1\,\dots\,T-5, T-8) & pass \\
			Hover linearisation asserted at import (C-1, C-2) & pass \\
			Residual and wrench both zero on the nominal plant (T-8) & pass \\
			Test suite, of which 26 tests pin audit findings & 86 pass \\
			Study $\leftrightarrow$ workspace constant agreement (\texttt{check\_consistency}) & pass \\
			Frame, CTBR and $B_d$ bridge self-tests (\texttt{check\_glue}) & pass \\
			Exported predictor parity, JAX against NumPy, tolerance $10^{-5}$ & $\approx10^{-14}$ \\
			Notebooks 1--7 executed end to end in one chain & pass \\
			\bottomrule
		\end{tabularx}
	\end{table}

	\begin{keybox}
		\keyhead{The gates fire, which is the point of having them}
		Executing the chain at \texttt{smoke} scale caused every quality gate in the
		study to trip: the training-health gate on collective saturation, the
		terminal-cost fit gate on $R^2\le0$, the adaptive precondition gate, and the
		undertrained-policy detector, which reports that the split between success and
		failure falls exactly along the training axis and names the scale table
		responsible. A pipeline whose gates stay silent on undertrained policies would
		be the more worrying outcome.
	\end{keybox}

	%-----------------------------------------------------------------------
	\section{Increment 0: baselines, horizon and cost sensitivity}
	\label{sec:res-baselines}
	%-----------------------------------------------------------------------

	\subsection{Classical anchors}

	\begin{table}[htbp]
		\centering
		\caption{Classical baselines under nominal conditions, median over $n_\mathrm{eval}$
			vehicles. The two LQR rows differ only in the reference feed-forward of
			\cref{sec:baselines}. Source: \texttt{artifacts/lltc/nb1\_classical.csv}.}
		\label{tab:res-classic}
		\small
		\begin{tabularx}{\linewidth}{@{}l X r r r r r r@{}}
			\toprule
			\textbf{Path} & \textbf{Controller} & \textbf{RMSE} & \textbf{IQR} & \textbf{$e_{\max}$} & \textbf{effort} & \textbf{smooth} & \textbf{sat} \\
			\midrule
			hover  & LQR                 & 0.0281 & 0.0099 & 0.2246 & 0.5311 & $3.113\times10^{-9}$ & 0.000 \\
			hover  & LQR (no feed-fwd)   & 0.0281 & 0.0099 & 0.2246 & 0.5311 & $3.113\times10^{-9}$ & 0.000 \\
			circle & LQR                 & 0.0626 & 0.0116 & 0.2666 & 0.5393 & $1.659\times10^{-6}$ & 0.000 \\
			circle & LQR (no feed-fwd)   & 0.0548 & 0.0146 & 0.2584 & 0.5398 & $1.832\times10^{-6}$ & 0.000 \\
			fig8   & LQR                 & 0.0945 & 0.0122 & 3.0172 & 0.5494 & $9.819\times10^{-5}$ & 0.002 \\
			fig8   & LQR (no feed-fwd)   & 0.0964 & 0.0151 & 2.6243 & 0.5519 & $8.498\times10^{-5}$ & 0.002 \\
			square & LQR                 & 0.1250 & 0.0940 & 0.6880 & 0.5670 & 0.004 & 0.000 \\
			square & LQR (no feed-fwd)   & 0.0599 & 0.0154 & 0.2840 & 0.5545 & 0.002 & 0.000 \\
			\bottomrule
		\end{tabularx}
	\end{table}

	\realfig{56mm}{\figLqrFeedforward}
	{The reference feed-forward in isolation: identical gain, identical model,
		with and without the term. The gap is what a pure state-feedback law loses by
		lagging a moving reference.}
	{fig:lqrff}

	\subsection{Horizon: accuracy against computation}
	\label{sec:res-horizon}

	\begin{table}[htbp]
		\centering
		\caption{Horizon sweep. The \emph{frozen} columns repeat the sweep with the
			reference held at $\ctrlstate_\refsym(t)$ over the horizon. Latency is
			single-vehicle and must be read against the \SI{20}{\milli\second} period.
			Source: \texttt{artifacts/common/nb1\_horizon.csv}.}
		\label{tab:res-horizon}
		\small
		\begin{tabularx}{\linewidth}{@{}r r r r r r r X@{}}
			\toprule
			& \multicolumn{3}{c}{\textbf{RMSE, preview}} & \multicolumn{3}{c}{\textbf{RMSE, frozen}} & \\
			\cmidrule(lr){2-4}\cmidrule(lr){5-7}
			$N$ & circle & fig8 & square & circle & fig8 & square & $t_\mathrm{solve}$ [\si{\milli\second}] \\
			\midrule
			1  & 0.0658 & 0.0676 & 0.0701 & 0.0758 & 0.0784 & 0.0889 & 12.20 \\
			2  & 0.0664 & 0.0664 & 0.0660 & 0.0990 & 0.0952 & 0.1045 & 28.58 \\
			3  & 0.0656 & 0.0647 & 0.0660 & 0.1218 & 0.1165 & 0.1234 & 39.78 \\
			5  & 0.0649 & 0.0638 & 0.0666 & 0.1681 & 0.1637 & 0.1700 & 61.88 \\
			10 & 0.0645 & 0.0626 & 0.0664 & 0.2787 & 0.2713 & 0.2761 & 121.54 \\
			20 & 0.0644 & 0.0625 & 0.0655 & 0.4606 & 0.4326 & 0.4538 & 257.61 \\
			\bottomrule
		\end{tabularx}
	\end{table}

	\realfigtwo{52mm}{\figHorizon}{\figCorner}
	{Effect of the prediction horizon.}
	{Tracking error and solve latency against $N$, preview versus frozen reference.}
	{Trajectory detail at a superellipse corner, showing where anticipation is worth
		its computation.}
	{fig:horizon}

	\begin{keybox}
		\keyhead{Expected reading}
		With the preview supplied, accuracy should improve with $N$ and saturate; with
		the reference genuinely frozen it should \emph{degrade} with $N$, because the
		solver plans to stop on the current waypoint. An assertion in the pipeline fails
		if the two families of columns ever coincide --- which is how the original
		implementation's silent failure was detected. A second reading to take from
		\cref{tab:res-horizon} is where p95 latency crosses \SI{20}{\milli\second},
		since every $N$ beyond that point is inadmissible whatever its accuracy.
	\end{keybox}

	\subsection{Noise sensitivity}
	\label{sec:res-noise}

	\begin{table}[htbp]
		\centering
		\caption[Baselines under measurement noise]{Degradation of the hand-built baselines under measurement noise.
			Source: \texttt{artifacts/common/nb1\_noise.csv}.}
		\label{tab:res-noise}
		\small
		\begin{tabularx}{\linewidth}{@{}l X r r r r@{}}
			\toprule
			\textbf{Noise} & \textbf{Controller} & \textbf{RMSE} & \textbf{$e_{\max}$} & \textbf{smooth} & \textbf{sat} \\
			\midrule
			off  & LQR & 0.0678 & 0.256 & $1.177\times10^{-6}$ & 0.000 \\
			     & NMPC $N{=}1$ & 0.0692 & 0.304 & $3.820\times10^{-6}$ & 0.000 \\
			     & NMPC $N{=}10$ & 0.0679 & 0.265 & $1.722\times10^{-5}$ & 0.000 \\
			\addlinespace[2pt]
			low  & LQR & 0.0679 & 0.254 & 0.009 & 0.000 \\
			     & NMPC $N{=}1$ & 0.0695 & 0.304 & 0.018 & 0.000 \\
			     & NMPC $N{=}10$ & 0.0681 & 0.262 & 0.011 & 0.000 \\
			\addlinespace[2pt]
			high & LQR & 0.0704 & 0.254 & 0.127 & 0.056 \\
			     & NMPC $N{=}1$ & 0.0716 & 0.302 & 0.201 & 0.002 \\
			     & NMPC $N{=}10$ & 0.0690 & 0.256 & 0.170 & 0.001 \\
			\midrule
			\multicolumn{2}{@{}l}{\textbf{Degradation high/off}} & 1.04 / 1.03 / 1.02 & --- & --- & --- \\
			\bottomrule
		\end{tabularx}
	\end{table}

	\realfig{52mm}{\figNoise}
	{Baseline degradation against measurement-noise level. Smoothness is the metric
		to watch: a predictive controller differentiating a noisy state produces jerk
		before it produces tracking error.}
	{fig:noise}

	\subsection{Sensitivity to the cost weights}
	\label{sec:res-cost}

	\begin{table}[htbp]
		\centering
		\caption{$\mat Q$/$\mat R$ sweep, $N=1$, nominal conditions. The spread between
			the best and worst admissible choice is the quantity that motivates learning the
			cost. Source: \texttt{artifacts/common/nb1\_qr\_spread.csv}.}
		\label{tab:res-qr}
		\small
		\begin{tabularx}{\linewidth}{@{}X r r r@{}}
			\toprule
			\textbf{Weighting} & $Q_\mathrm{pos}$ & $R_\mathrm{rate}$ & \textbf{RMSE [\si{\metre}]} \\
			\midrule
			Best admissible & 10 & 0.05 & 0.0858 \\
			Median admissible & --- & --- & 0.1250 \\
			Worst admissible & 0.5 & 20 & 0.2014 \\
			\midrule
			\multicolumn{3}{@{}l}{\textbf{Best-to-worst ratio}} & \textbf{2.35$\times$} \\
			\bottomrule
		\end{tabularx}
	\end{table}

	\realfig{56mm}{\figQrHeatmap}
	{Heat map of closed-loop RMSE over the $(\mat Q,\mat R)$ grid, with the
		hand-chosen operating point marked. The sensitivity of a nominally exact NMPC to
		weights that no part of its theory determines.}
	{fig:qrsweep}

	\begin{remark}[The operating point is chosen on one controller and scored on another]
		The hand-tuned weights were selected on an \textbf{LQR} sweep, while
		\cref{tab:res-qr} scores \textbf{NMPC $N{=}1$}. The two disagree slightly about
		the optimum, and that disagreement is itself part of the answer: even ``the
		right weights'' is controller dependent, which is exactly what learning the
		cost is meant to remove.
	\end{remark}

	\subsection{Sensitivity to model mismatch, and to the direction of it}
	\label{sec:res-sensitivity}

	\Cref{sec:argument} claims the classical answer pays three times. The cost-weight
	sweep above measures the third. This subsection measures the first: what happens
	when the plant is not the one the controller plans with. Mass and inertia are
	swept on \texttt{fig8}, which is the path that is \emph{partly} at the
	feasibility cap (\SI{6.2}{\percent} of radii, against the superellipse's
	\SI{100}{\percent} and the circle's none), so the sweep stresses the controller
	without the collective already being pinned against its box.

	\begin{table}[htbp]
		\centering
		\caption[Where the hand-built controllers give way]{Where each hand-built
			controller gives way under plant mismatch. \textbf{Break} is
			\emph{relative} --- the first factor at which RMSE exceeds three times that
			controller's \emph{own} nominal --- and must not be read alone: a controller
			that starts more accurate trips it sooner while still leading.
			\textbf{Crossover} is \emph{absolute}: the factor at which NMPC $N{=}1$
			stops beating LQR, computed separately on each side of nominal, because
			mismatch has two directions. Source:
			\texttt{artifacts/sensitivity/model\_hand\_breakpoints.csv}.}
		\label{tab:res-mismatch-hand}
		\small
		\begin{tabularx}{\linewidth}{@{}X r r r r@{}}
			\toprule
			& \multicolumn{2}{c}{\textbf{Nominal RMSE [\si{\metre}]}}
			& \multicolumn{2}{c}{\textbf{Break ($3\times$ own nominal)}} \\
			\cmidrule(lr){2-3}\cmidrule(lr){4-5}
			\textbf{Axis} & LQR & NMPC $N{=}1$ & LQR & NMPC $N{=}1$ \\
			\midrule
			$\lambda_m$ (mass) & 0.1016 & 0.0517 & 1.75$\times$ & 1.75$\times$ \\
			\midrule
			\multicolumn{5}{@{}X}{\footnotesize NMPC $N{=}1$ stops beating LQR at $\lambda_m=1.75$ above nominal; below nominal it still leads at the widest factor swept. The inertia axis did not reach either threshold inside the swept range and is therefore not tabulated.} \\
			\bottomrule
		\end{tabularx}
	\end{table}

	\realfig{58mm}{\figMismatchHand}
	{Hand-built controllers against a plant they did not plan with: mass and inertia
		scaled about nominal, on \texttt{fig8}. The vertical line is the nominal
		plant. The shape to read is not the height but the \emph{slope} either side of
		it, and whether the two curves cross.}
	{fig:mismatch-hand}

	\begin{keybox}
		\keyhead{Expected reading, and the trap in it}
		The online solve plans several steps against the model, so a wrong model
		compounds where a fixed gain merely reacts; NMPC $N{=}1$ should therefore
		degrade \emph{faster} than LQR. But it also starts from a lower error, so a
		threshold defined relative to each controller's own nominal flatters the worse
		controller automatically --- measured here, by roughly a factor of four. The
		two columns of \cref{tab:res-mismatch-hand} answer different questions and are
		expected to disagree; the absolute crossover is the one that says which
		controller to fly.
	\end{keybox}

	\realfig{62mm}{\figMismatchCost}
	{Cost-weight stress on a grid deliberately wider than the admissible region, for
		both hand-built controllers on all three paths. This is the quantity
		\cref{sec:acmpc} removes, measured where it is worst rather than where it is
		comfortable.}
	{fig:mismatch-cost}

	\begin{table}[htbp]
		\centering
		\caption[Cost-weight spread, wide grid]{Spread of closed-loop RMSE over the
			wide $(Q_\mathrm{pos}, R_\mathrm{rate})$ grid. The ratio is the fraction of
			apparent ``controller performance'' that is really the human's choice of
			weights. Source: \texttt{artifacts/sensitivity/cost\_hand\_spread.csv}.}
		\label{tab:res-cost-spread}
		\small
		\begin{tabularx}{\linewidth}{@{}l X r r r@{}}
			\toprule
			\textbf{Path} & \textbf{Controller} & \textbf{Best} & \textbf{Worst}
			& \textbf{Ratio} \\
			\midrule
			circle & LQR / NMPC $N{=}1$ & 0.0622 / 0.1054 & 0.3700 / 0.1211 & 5.9 / 1.1$\times$ \\
			fig8   & LQR / NMPC $N{=}1$ & 0.0621 / 0.1048 & 0.3570 / 0.1337 & 5.7 / 1.3$\times$ \\
			square & LQR / NMPC $N{=}1$ & 0.1385 / 0.1094 & 0.4576 / 0.1350 & 3.3 / 1.2$\times$ \\
			\bottomrule
		\end{tabularx}
	\end{table}

	The same two axes are swept for the learned controllers in
	\cref{sec:res-acmpc-mismatch}, on the same path and the same grid, so the two
	families are directly comparable.

	%-----------------------------------------------------------------------
	\section{Increment 1: learned terminal cost}
	\label{sec:res-lltc}
	%-----------------------------------------------------------------------

	\subsection{Fit quality}

	\begin{table}[htbp]
		\centering
		\caption{Terminal-cost regression diagnostics. $R^2\le0$ means the fit is worse
			than predicting the mean of $V_1$; the gate treats that as a negative result
			about the regression, not as a tuning detail. Source:
			\texttt{artifacts/lltc/nb2\_fit.csv}.}
		\label{tab:res-lltcfit}
		\small
		\begin{tabularx}{\linewidth}{@{}X r r r r r@{}}
			\toprule
			\textbf{Quantity} & \textbf{$R^2$} & \textbf{NRMSE} & \textbf{acceptance} & $n$ & \textbf{$X_T$ reach [\si{\metre}]} \\
			\midrule
			Regression on $V_1$ & 0.999998 & 0.0014 & 100\% & 256 & 0.5649 \\
			\midrule
			\multicolumn{6}{@{}l}{\itshape Candidate displacement \SI{0.35}{\metre}; ratio to reach \num{0.62}; gate passed} \\
			\bottomrule
		\end{tabularx}
	\end{table}

	\realfigtwo{52mm}{\figLltcFit}{\figLltcSpectrum}
	{Behaviour of the learned terminal cost.}
	{Predicted terminal cost against realised cost-to-go; the diagonal is a perfect
		fit and the gate is at $R^2=0$.}
	{Acceptance rate and terminal-set reach against the terminal weighting --- the
		mechanism behind the fit quality.}
	{fig:lltcfit}

	\subsection{Horizon equivalence}

	\begin{table}[htbp]
		\centering
		\caption[Does \LLTC{} at $N=1$ reach \NMPC{} at $N=1$?]{The claim under test: does \LLTC{} at $N=1$ reach \NMPC{} at $N=1$ or
			better? A ratio below $1.0$ supports it. Source:
			\texttt{artifacts/lltc/nb2\_horizon\_equivalence.csv}.}
		\label{tab:res-lltc}
		\small
		\begin{tabularx}{\linewidth}{@{}X r r r@{}}
			\toprule
			\textbf{Controller} & \textbf{circle} & \textbf{fig8} & \textbf{square} \\
			\midrule
			NMPC $N{=}1$ & 0.0712 & 0.0683 & 0.0723 \\
			NMPC $N{=}3$ & 0.0709 & 0.0639 & 0.0702 \\
			NMPC $N{=}10$ & 0.0697 & 0.0620 & 0.0704 \\
			\LLTC{} $N{=}1$ & 1.2615 & 1.2551 & 1.2934 \\
			\midrule
			\textbf{Ratio \LLTC/NMPC $N{=}1$} & 17.7$\times$ & 18.4$\times$ & 17.9$\times$ \\
			\bottomrule
		\end{tabularx}
	\end{table}

	\realfig{54mm}{\figLltcLocality}
	{In-distribution versus out-of-distribution performance of the learned terminal
		cost. $\mat P_\theta$ is a \emph{local} object: it is fitted where the data is,
		and away from that region it carries no positive-definiteness or dissipativity
		guarantee and no reason to remain a valid cost-to-go.}
	{fig:lltcood}

	\subsection{Analysis}
	\label{sec:res-lltc-analysis}

	\textbf{The claim did not hold, and it failed for a reason that is not fit
	quality.} Both halves of that sentence matter. The regression of
	\cref{tab:res-lltcfit} is as good as a regression gets: $R^2=\num{0.999998}$,
	NRMSE \num{0.0014}, every one of the $256$ candidates accepted, and a
	displacement-to-reach ratio of \num{0.62} that puts every candidate well
	inside the region where the local quadratic model is supposed to hold. The
	gate did not fire. $\mat P_\theta$ predicts the realised cost-to-go almost
	exactly on the distribution it was fitted on.

	And the closed loop is \emph{eighteen times worse}. \Cref{tab:res-lltc} puts
	\LLTC{} at $N{=}1$ at \SIrange{1.2615}{1.2934}{\metre} against NMPC $N{=}1$ at
	\SIrange{0.0683}{0.0723}{\metre} on the same three paths --- a ratio of
	\numrange{17.7}{18.4}, consistent across paths and therefore not an artefact
	of one geometry. Note also what \cref{tab:res-lltc} says about the premise:
	NMPC improves from \num{0.0712} at $N{=}1$ to \num{0.0697} at $N{=}10$ on
	\texttt{circle}, a gain of \SI{2}{\percent}. The horizon was worth very
	little here to begin with, so even a terminal cost that worked perfectly had
	almost nothing to recover.

	The failure is \textbf{locality}, and it is structural rather than a matter
	of budget. $\mat P_\theta$ is fitted where the candidates are --- around the
	reference, inside the terminal set. A closed loop does not stay there: it is
	driven out by disturbance and by its own transients, and outside the fitted
	region $\mat P_\theta$ carries no positive-definiteness argument, no
	dissipativity argument, and no reason whatever to remain a valid cost-to-go.
	\Cref{fig:lltcood} shows the in- against out-of-distribution behaviour
	directly. The second, quieter problem is that $V_1$ was itself produced by a
	long-horizon solve driven by the hand-chosen $\mat Q$ and $\mat R$: the
	terminal cost therefore \emph{inherits} those weights rather than replacing
	them, so even a perfect fit leaves the third liability of
	\cref{sec:argument} exactly where it was.

	Both observations point the same way, and that is why Increment~2 learns the
	cost at every stage rather than only at the tail. A cost that acts at one
	point of the horizon can only be as good as the region it was fitted in; a
	cost that acts at every stage is trained by the closed loop it will be used
	in, which is a different object.

	%-----------------------------------------------------------------------
	\section{Increment 2: learned cost map}
	\label{sec:res-acmpc}
	%-----------------------------------------------------------------------

	\begin{table}[htbp]
		\centering
		\caption{Cost-matrix parametrisation, final reward and closed-loop RMSE after
			the common budget. The published claim is that only \texttt{diag} learns;
			this table decides it on these tasks. Differences smaller than the seed
			spread of \cref{tab:res-seeds} are not results. Source:
			\texttt{artifacts/acmpc/nb3\_representation.csv}.}
		\label{tab:res-rep}
		\small
		\begin{tabularx}{\linewidth}{@{}X r r r r@{}}
			\toprule
			\textbf{Task} & \texttt{diag} & \texttt{chol} & \texttt{full} & \textbf{Winner} \\
			\midrule
			circle & 0.1042 & 0.0872 & 0.1012 & \texttt{chol} \\
			fig8 & 0.1113 & 0.0863 & 0.0954 & \texttt{chol} \\
			stabilize & 0.0723 & 0.0779 & 0.0864 & \texttt{diag} \\
			\midrule
			\textbf{Margin over seed spread} & 0.0170 & 0.0251 & 0.0141 & --- \\
			\bottomrule
		\end{tabularx}
	\end{table}

	\begin{table}[htbp]
		\centering
		\caption[Exploration-noise sweep]{Exploration-noise sweep for the
			differentiable layer, \ACMPC{} against a model-free MLP trained on the same
			reward with the same action space. \textbf{sat} is the fraction of training
			steps with the collective on its box. The seed spread of
			\cref{tab:res-seeds} is the scale against which differences must be judged.
			Source: \texttt{artifacts/acmpc/nb3\_exploration.csv}.}
		\label{tab:res-acmpc-sweeps}
		\small
		\begin{tabularx}{\linewidth}{@{}X r r r r@{}}
			\toprule
			& \multicolumn{2}{c}{\textbf{RMSE [\si{\metre}]}} & & \\
			\cmidrule(lr){2-3}
			\textbf{Exploration $\sigma$} & \ACMPC & MLP & \textbf{ratio} & \textbf{sat, \ACMPC} \\
			\midrule
			0.05 & 0.1042 & 1.2042 & 11.6$\times$ & 0.13\% \\
			0.15 & 0.0939 & 0.9946 & 10.6$\times$ & 4.11\% \\
			0.30 & 0.1346 & 1.1386 & 8.5$\times$ & 21.51\% \\
			\midrule
			\textbf{Degradation, $\sigma=0.15\to0.30$} & \textbf{$+43\,\%$} & $+14\,\%$ & --- & \textbf{5.2$\times$} \\
			\bottomrule
		\end{tabularx}
	\end{table}


	\subsection{Mismatch, with no weights left to blame}
	\label{sec:res-acmpc-mismatch}

	The two axes of \cref{sec:res-sensitivity} are now swept for the learned arms,
	on the same path, the same grid and the same seed, so the two families are read
	off one picture. The cost-weight axis has no counterpart here, which is the
	point of this increment: there is no hand-chosen stage weight left to stress.
	What remains is the axis the cost map cannot touch --- inside the layer, the
	dynamics are still the nominal model, so a learned cost can reshape \emph{what}
	the solver wants but not \emph{what it believes the vehicle will do}.

	\begin{table}[htbp]
		\centering
		\caption[Where the learned controllers give way]{Where each learned controller
			gives way under plant mismatch, on \texttt{fig8}. Columns are read exactly as
			in \cref{tab:res-mismatch-hand}: \textbf{Break} is relative to that
			controller's own nominal and therefore flatters whichever arm starts worse,
			\textbf{Crossover} is the absolute factor at which the adaptive arm stops
			beating \ACMPC{}, computed separately below and above nominal. Source:
			\texttt{artifacts/sensitivity/model\_learned\_breakpoints.csv}.}
		\label{tab:res-mismatch-learned}
		\small
		\begin{tabularx}{\linewidth}{@{}X r r@{}}
			\toprule
			& \multicolumn{2}{c}{\textbf{Break ($3\times$ own nominal)}} \\
			\cmidrule(lr){2-3}
			\textbf{Axis} & \ACMPC{} $N{=}1$ & Adaptive $N{=}1$ \\
			\midrule
			$\lambda_m$ (mass) & 0.55$\times$ & 0.55$\times$ \\
			\midrule
			\multicolumn{3}{@{}X}{\footnotesize Both arms break at the same factor, and on the \emph{light} side: at $\lambda_m=0.55$ the vehicle is \SI{45}{\percent} under nominal mass and the hover collective leaves its trim. Neither arm crossed the other inside the swept range, so no crossover is reported. The inertia axis did not reach the threshold and is not tabulated.} \\
			\bottomrule
		\end{tabularx}
	\end{table}

	\realfig{58mm}{\figMismatchLearned}
	{The learned arms against a plant they did not plan with, on the same mass and
		inertia grid as \cref{fig:mismatch-hand}. The adaptive arm carries its residual
		predictor, so a parameter error is exactly the condition the estimator is
		supposed to notice: this is the sharpest available test of whether estimating
		beats averaging.}
	{fig:mismatch-learned}

	\begin{keybox}
		\keyhead{What this figure can and cannot settle}
		A scaled mass or inertia is a \emph{structured} error --- constant in the body
		frame and, for mass, exactly the standing force the predictor of
		\cref{sec:adaptive} is trained to reconstruct. If the adaptive arm does not
		separate from \ACMPC{} \emph{here}, the estimator is not adding information
		anywhere, and \cref{sec:res-adaptive} should be read with that in mind. If it
		does, the separation still has to clear the seed spread of
		\cref{tab:res-seeds} before it counts. Note also that the mass axis moves the
		hover thrust and therefore the distance to the collective box: the saturation
		column of the source CSV must be read beside the RMSE, or a cell that has run
		out of actuator will be reported as a cell that has run out of controller.
	\end{keybox}

	\subsection{Analysis}

	\textbf{The weights are gone, and they were not free.} Every hand-chosen
	stage weight was removed: the \texttt{tuned} column of \cref{tab:ledger}
	reads \textbf{0} for \ACMPC{} against \num{13} for every hand-built
	controller, which is the claim this increment was built to make. The price
	is in the same table. On S1 \ACMPC{} $N{=}1$ tracks at \SI{0.1040}{\metre}
	against \SI{0.0739}{\metre} for the tuned NMPC $N{=}1$ --- \SI{41}{\percent}
	worse --- and \SI{0.0945}{\metre} for LQR, \SI{10}{\percent} worse. That is
	a \textbf{trade}, and calling it a win would misrepresent it: thirteen
	constants that a human chose were exchanged for roughly forty per cent of
	accuracy against the controller those constants were tuned for. Whether the
	trade is worth taking depends on how much the tuning generalises, which is
	what \cref{sec:res-sensitivity} measures and what \cref{tab:res-cost-spread}
	prices at up to \num{5.9}$\times$ for LQR across the wide grid.

	\textbf{The published representation finding did not reproduce.} The claim
	under test was that only the diagonal parametrisation learns.
	\Cref{tab:res-rep} has \texttt{chol} winning on \texttt{circle}
	(\num{0.0872} against \num{0.1042}) and on \texttt{fig8} (\num{0.0863}
	against \num{0.1113}), with \texttt{diag} winning only on
	\texttt{stabilize}. These margins are \numrange{0.014}{0.025}, between two
	and four times the seed spread of \SI{0.0064}{\metre}
	(\cref{tab:res-seeds}), so they are resolved rather than noise. The result
	is consistent with the initialisation argument of \cref{sec:acmpc}: the
	richer forms are not incapable, they start further from a useful cost and
	need longer to get there. At the \texttt{medium} budget used here they have
	evidently had long enough. What cannot be concluded from this campaign is
	the converse --- that \texttt{diag} would not catch up with more iterations
	--- because no budget sweep was run to test it.

	\textbf{The saturation mechanism is supported, in the relative sense the
	hypothesis actually makes.} \Cref{tab:res-acmpc-sweeps} is unambiguous on
	absolute performance: \ACMPC{} beats the model-free MLP by between
	\num{8.5} and \num{11.6}$\times$ at every noise level, and never loses to
	it. But the mechanism claim was not about who wins, it was that exploration
	noise driving the collective onto its box corrupts the linearisation
	$\mat A_k,\mat B_k$ \emph{inside} the solve, and therefore hurts a
	differentiable layer in a way it cannot hurt an MLP. Both columns the claim
	needs move together: raising $\sigma$ from \num{0.15} to \num{0.30} takes
	\ACMPC{} saturation from \SI{4.11}{\percent} to \SI{21.51}{\percent}, a
	factor of \num{5.2}, and its RMSE degrades by \SI{43}{\percent} against the
	MLP's \SI{14}{\percent} over the same step. The differentiable layer
	degrades three times faster, and it does so exactly where the saturation
	column says the box is binding. Reporting the reward alone would have
	missed this: \ACMPC{}'s reward is better than the MLP's by two orders of
	magnitude at every $\sigma$, which on its own would have read as
	insensitivity.

	%-----------------------------------------------------------------------
	\section{Increment 3: domain randomisation}
	\label{sec:res-dr}
	%-----------------------------------------------------------------------

	\begin{table}[htbp]
		\centering
		\caption{Robustness and its price. The conservatism premium is evaluated on the
			\emph{nominal} plant and is \textbf{signed}: a positive value is the insurance
			premium the section predicts, a negative one means DR tracked the nominal
			plant better. Source: \texttt{artifacts/domrand/T6\_robustness.csv}.}
		\label{tab:res-dr}
		\small
		\begin{tabularx}{\linewidth}{@{}X r r r r@{}}
			\toprule
			& \multicolumn{1}{c}{\textbf{nominal}}
			& \multicolumn{3}{c}{\textbf{S2, moderated}} \\
			\cmidrule(lr){2-2}\cmidrule(lr){3-5}
			\textbf{Policy} & RMSE & RMSE & IQR & $e_{\max}$ \\
			\midrule
			Nominal-only & 0.1025 & 0.1856 & 0.0406 & 0.5381 \\
			DR, all axes & 0.1002 & 0.1554 & 0.0497 & 0.4497 \\
			DR $+$ measurement noise & 0.1118 & 0.1555 & 0.0459 & 0.4232 \\
			NMPC $N{=}1$ (no learning) & 0.0980 & 0.1521 & 0.0386 & 0.3690 \\
			\midrule
			\multicolumn{5}{@{}X}{\footnotesize Seed-averaged over the two seeds of \cref{tab:res-seeds}: conservatism premium $+0.70\,\%$ on the nominal plant, robustness gain $+6.90\,\%$ on S2. Both are signed; both are an order of magnitude below the seed spread, and \cref{sec:res-dr-analysis} reads them accordingly.} \\
			\bottomrule
		\end{tabularx}
	\end{table}

	\begin{table}[htbp]
		\centering
		\caption[Seed variance against vehicle variance]{The two variances, reported separately and never pooled. The seed
			spread is the resolution of every comparison in this chapter. Sources:
			\texttt{artifacts/domrand/seed\_study.csv}, \texttt{artifacts/acmpc/nb3\_seeds.csv}.}
		\label{tab:res-seeds}
		\small
		\begin{tabularx}{\linewidth}{@{}X r r@{}}
			\toprule
			\textbf{Policy} & \textbf{Seed variance} & \textbf{Episode IQR} \\
			\midrule
			Nominal-only & 0.0004 & 0.0314 \\
			DR, all axes & 0.0064 & 0.0330 \\
			\ACMPC{} (three seeds) & 0.0061 & --- \\
			\midrule
			\textbf{Seed spread used as the comparison floor} & \textbf{0.0064\,\si{\metre}} & --- \\
			\bottomrule
		\end{tabularx}
	\end{table}

	\realfigtwo{52mm}{\figAxisCompetence}{\figVariances}
	{Domain randomisation.}
	{Per-axis competence: RMSE against mass, rotor lag, inertia and thrust scale.
		The abscissa is the \emph{realised} mean factor within each bin, not the bin
		centre, because bins have random membership.}
	{Seed variance against episode variance --- the two move in opposite directions,
		and reporting only one manufactures a robustness claim.}
	{fig:dr}

	\begin{table}[htbp]
		\centering
		\caption{Sequential randomisation: backward transfer and plasticity. Positive
			backward transfer means the stage improved after later training; negative
			means it was forgotten. Source:
			\texttt{artifacts/domrand/\{curriculum,backward\_transfer\}.csv}.}
		\label{tab:res-forget}
		\small
		\begin{tabularx}{\linewidth}{@{}X r r r@{}}
			\toprule
			\textbf{Stage} & \textbf{BWT [\si{\metre}]} & \textbf{Dead units} & \textbf{Saturated outputs} \\
			\midrule
			nominal & $-0.0070$ & 0.11\,\% & 0.00\,\% \\
			mass & $-0.0037$ & 0.14\,\% & 0.00\,\% \\
			lag & $-0.0072$ & 0.14\,\% & 0.00\,\% \\
			all & --- & 0.12\,\% & 0.00\,\% \\
			\bottomrule
		\end{tabularx}
	\end{table}

	\realfig{52mm}{\figCurriculum}
	{Sequential curriculum: performance on every earlier stage after training
		through each later one. The diagonal is performance when the stage was
		current; the columns below it are what survived.}
	{fig:curriculum}

	\subsection{Analysis}
	\label{sec:res-dr-analysis}

	\textbf{The conservatism premium is positive in sign and unresolved in
	magnitude, and the second half of that matters more than the first.} The
	seed-averaged premium is $+\SI{0.70}{\percent}$ --- \SI{0.0007}{\metre} on a
	nominal RMSE of about \SI{0.10}{\metre}. The seed spread of
	\cref{tab:res-seeds} is \SI{0.0064}{\metre}, nine times larger. This
	campaign therefore \emph{cannot measure} the conservatism premium, and the
	single-seed rows of \cref{tab:res-dr} make that concrete by disagreeing with
	the average about its sign: there DR-all tracks the nominal plant at
	\SI{0.1002}{\metre} against nominal-only at \SI{0.1025}{\metre}, a premium
	of $-\SI{2.2}{\percent}$. Two seeds are not enough to tell $+0.7$ from
	$-2.2$. Separating an effect of \SI{0.0007}{\metre} from a spread of
	\SI{0.0064}{\metre} needs on the order of $(0.0064/0.0007)^2\approx80$ times
	the seeds used here, so roughly a hundred --- which this campaign does not
	have and was not designed to have. The honest statement is that the premium
	is \textbf{indistinguishable from zero at this seed count}, and two readings
	survive it: randomisation acting as a regulariser at this iteration budget,
	or the nominal-only control group being undertrained. Nothing here
	distinguishes them.

	The robustness gain stands on firmer ground. On S2 the seed-averaged gain is
	$+\SI{6.90}{\percent}$, \SI{0.0116}{\metre}, about \num{1.8} times the seed
	spread, and the single-seed rows agree in sign and exceed it
	(\SI{0.1856}{\metre} to \SI{0.1554}{\metre}, \SI{16}{\percent}). Randomising
	the plant does buy something on a disturbed plant. Adding measurement noise
	buys nothing further on S2 (\SI{0.1555}{\metre}, indistinguishable from
	DR-all) while costing \SI{0.1118}{\metre} on the nominal plant, the worst
	nominal figure of the four arms.

	\textbf{The structural conclusion does not depend on any of these signs.}
	Look at the last row of \cref{tab:res-dr}: untrained NMPC $N{=}1$ beats all
	three learned arms on both suites, at \SI{0.0980}{\metre} nominal and
	\SI{0.1521}{\metre} on S2. Domain randomisation asks one fixed policy to be
	adequate across a family of plants, and a policy that cannot \emph{observe}
	which plant it is on cannot do better than the compromise --- it has no
	channel through which specialisation could happen, however many plants it is
	shown. That is a limit of the method, not of the budget, and it is precisely
	the limit Increment~4 is built to remove: give the controller an estimate of
	the plant it is actually on, and the compromise becomes unnecessary.

	%-----------------------------------------------------------------------
	\section{Increment 4: adaptive scheme}
	\label{sec:res-adaptive}
	%-----------------------------------------------------------------------

	\subsection{The precondition}

	\begin{table}[htbp]
		\centering
		\caption{Training-health gate. A policy that spends training saturated against
			its input box never learns a disturbance-response manifold, so the oracle
			decomposition of \cref{sec:adaptive} is void for it. \textbf{This table is a
				precondition, read before the closed-loop results.} Source:
			\texttt{artifacts/acmpc\_adaptive/nb5\_gate.csv}.}
		\label{tab:res-trainhealth}
		\small
		\begin{tabularx}{\linewidth}{@{}X r r r r@{}}
			\toprule
			\textbf{Policy} & \textbf{Final reward} & \textbf{Collective saturation} & \textbf{Crash rate} & \textbf{Gate ($<5\%$)} \\
			\midrule
			Base & -0.0877 & 0.11\% & 0.00\% & yes \\
			Robust & -0.1627 & 0.12\% & 0.00\% & yes \\
			Oracle & -0.1628 & 0.13\% & 0.00\% & yes \\
			\bottomrule
		\end{tabularx}
	\end{table}

	\realfig{52mm}{\figTrainingGate}
	{Collective saturation through training for the three precondition policies,
		against the \SI{5}{\percent} gate. The wrench-randomisation magnitude is the
		lever: too small and the disturbance does not matter, too large and the policy
		trains against the stop.}
	{fig:gate}

	\subsection{The residual predictor in isolation}

	\begin{table}[htbp]
		\centering
		\caption{Predictor accuracy and cost on the identical $\vH\times\vFrame$ window
			and target. \textbf{Admissibility is decided by latency} against the
			\SI{20}{\milli\second} control period, not by $R^2$: a predictor at
			\SI{137}{\milli\second} with $R^2=0.88$ is not deployable, one at
			\SI{2}{\milli\second} with $R^2=0.82$ is. Source:
			\texttt{artifacts/acmpc\_adaptive/nb5\_estimators.csv}.}
		\label{tab:res-rdp}
		\small
		\begin{tabularx}{\linewidth}{@{}l r r r r r X@{}}
			\toprule
			\textbf{Encoder} & \textbf{$R^2$ all} & \textbf{$R^2$ force} & \textbf{$R^2$ moment} & \textbf{params} & \textbf{p95 [\si{\milli\second}]} & \textbf{Admissible?} \\
			\midrule
			GRU~\cite{cho2014gru} & 0.9988 & 0.9985 & 0.9992 & 97542 & 1.24 & yes \\
			LSTM~\cite{hochreiter1997lstm} & 0.9962 & 0.9945 & 0.9980 & 129926 & 1.44 & yes \\
			TCN~\cite{bai2018tcn} & 0.9951 & 0.9943 & 0.9959 & 134982 & 0.91 & yes \\
			CNN & 0.9856 & 0.9795 & 0.9917 & 90950 & 0.31 & yes \\
			\midrule
			\textbf{Shipped} & \multicolumn{6}{l}{GRU\ \ (chosen among the admissible encoders by $R^2$)} \\
			\bottomrule
		\end{tabularx}
	\end{table}

	\realfigtwo{52mm}{\figEncoderTrade}{\figChannelRtwo}
	{Residual predictor.}
	{Accuracy--latency trade across encoders, with the control period marked as a
		vertical admissibility line.}
	{Per-channel $R^2$ for the three acceleration and three rate components, with
		the break-even gate at zero.}
	{fig:rdp}

	\subsection{How much history the predictor needs}
	\label{sec:res-rdp-window}

	The window length $H$ was fixed at $64$ frames for the comparison above. It is
	not a free parameter: $H$ sets the deployment ring-buffer length, the warm-up
	after every respawn during which the estimate is unusable
	(\SI{320}{\milli\second} at $H{=}16$ to \SI{2.56}{\second} at $H{=}128$), and
	the inference latency that has to fit inside the \SI{20}{\milli\second} period.
	It also sizes two of the four architectures, since \texttt{tcn\_blocks\_for} and
	\texttt{cnn\_layers\_for} pick the smallest depth whose receptive field covers
	the whole window --- a TCN at $H{=}128$ is a deeper network, not the same one fed
	more input. The sweep generates the flights \emph{once} at the largest $H$ and
	re-windows them, so a window-length comparison is not also a different-data
	comparison.

	\begin{table}[htbp]
		\centering
		\caption[Predictor accuracy against window length]{Predictor accuracy against
			window length, every encoder on identical flights. The question is whether a
			longer window buys accuracy or the residual is essentially memoryless; the
			answer is a \emph{pair} of columns, because $H$ is paid twice --- once per
			inference and once per respawn. Source:
			\texttt{artifacts/sensitivity/rdp\_window.csv}.}
		\label{tab:res-rdp-window}
		\small
		\begin{tabularx}{\linewidth}{@{}l X r r r r@{}}
			\toprule
			\textbf{$H$} & \textbf{Warm-up} & \textbf{GRU} & \textbf{LSTM} & \textbf{TCN} & \textbf{CNN} \\
			\midrule
			\multicolumn{6}{@{}l}{\itshape overall $R^2$} \\
			\quad 16  & \SI{320}{\milli\second}  & 0.9917 & 0.9420 & 0.9890 & 0.9862 \\
			\quad 32  & \SI{640}{\milli\second}  & 0.9924 & 0.9400 & 0.9848 & 0.9785 \\
			\quad 64  & \SI{1.28}{\second}       & 0.9932 & 0.9438 & 0.9818 & 0.9531 \\
			\quad 128 & \SI{2.56}{\second}       & 0.9945 & 0.9693 & 0.9817 & 0.9568 \\
			\midrule
			\multicolumn{6}{@{}l}{\itshape p95 inference latency, one window [\si{\milli\second}]} \\
			\quad 16  & ---                      & 0.35 & 0.41 & 0.37 & 0.13 \\
			\quad 128 & ---                      & 2.46 & 2.67 & 1.20 & 0.35 \\
			\midrule
			\multicolumn{2}{@{}l}{Shortest window within \SI{5}{\percent} of the best cell}
			& \multicolumn{4}{r}{$H=16$} \\
			\bottomrule
		\end{tabularx}
	\end{table}

	\realfig{50mm}{\figRdpWindow}
	{Window length against accuracy, against the force/moment split, and against
		inference latency. Force and moment are shown separately in the middle panel
		because \cref{sec:rdp-theory} predicts they are not equally observable: a
		standing moment survives only in the motor commands once the rate integrator
		has absorbed it, so the moment channels are the ones that should need history.}
	{fig:rdp-window}

	\begin{keybox}
		\keyhead{Why the shortest adequate window wins}
		The deployment cost of $H$ is not the inference time --- all four encoders are
		far inside the period at every $H$ swept. It is the warm-up: for the first $H$
		frames after a respawn, arm or estimator reset, variants B and C are feeding a
		prediction made from a partly empty buffer into the dynamics the solver plans
		with. The right choice is therefore the \emph{shortest} window that is within
		noise of the best one, not the most accurate window; the last row of
		\cref{tab:res-rdp-window} reports that cell explicitly. If accuracy is flat in
		$H$, the honest conclusion is that the residual is close to memoryless at this
		sampling rate and the recurrent encoder is not earning its state.
	\end{keybox}

	\subsection{Consistency of the two disturbance representations}

	\Cref{prop:convert} asserts that \cref{def:residual,def:wrench} describe the
	same disturbance in different coordinates. That is a statement about the model,
	so it is checked against the plant rather than assumed: \cref{tab:res-dmod}
	converts the measured wrench through \cref{eq:convert} and compares it with the
	measured residual, on scenarios chosen so that the moment block is the dominant
	one.

	\begin{table}[htbp]
		\centering
		\caption{The conversion of \cref{eq:convert} measured against the true
			residual of \cref{def:residual}, under \texttt{mode=exact}. A ratio of $1$
			means the two representations agree, which \cref{prop:convert} predicts
			for \emph{both} blocks. The \texttt{none} ablation zeroes the moment block
			by construction and is omitted here. The \texttt{central} scenario places
			the payload on the axis, so it carries no moment and the moment ratio is
			undefined: both sides are at the float noise floor,
			\SI{5.52e-15}{\radian\per\square\second}. Source:
			\texttt{artifacts/acmpc\_adaptive/nb5\_dmod\_validity.csv}.}
		\label{tab:res-dmod}
		\small
		\begin{tabularx}{\linewidth}{@{}l r X r r r@{}}
			\toprule
			\textbf{Scenario} & \textbf{Level} & & $\|\vec\alpha_\mathrm{res}\|$ conv. & $\|\vec\alpha_\mathrm{res}\|$ true & \textbf{ratio} \\
			\midrule
			central & 0.15 & force  & \multicolumn{2}{c}{$1.4709900000000$} & $1.000000000000$ \\
			        &      & moment & \multicolumn{2}{c}{$5.52\times10^{-15}$} & --- \\
			\addlinespace[2pt]
			asym    & 0.01 & force  & \multicolumn{2}{c}{$0.0980660004983$} & $1.000000000000$ \\
			        &      & moment & $2.0857104599945$ & $2.0857104599945$ & $1.000000000000$ \\
			\addlinespace[2pt]
			asym    & 0.07 & force  & \multicolumn{2}{c}{$0.6864620030776$} & $1.000000000000$ \\
			        &      & moment & $14.599989547222$ & $14.599989547222$ & $1.000000000000$ \\
			\bottomrule
		\end{tabularx}
	\end{table}

	Both blocks agree to the last significant figure the double-precision
	arithmetic carries, on every scenario where the block is non-zero. The
	moment ratios are $0.9999999999999998$ and $1.0$ as read from the artefact;
	they are quoted to twelve figures above only to make the point that the
	agreement is exact rather than close.

	\begin{keybox}
		\keyhead{What this table can and cannot establish}
		A ratio of unity confirms that the conversion is implemented as
		\cref{eq:convert} states and that the two accessors agree on the plant. It
		does \emph{not} establish that the nominal $m$ and $\mat J_\mathrm{nom}$ in
		\cref{eq:convert} are the true ones --- under randomisation they are not, by
		construction --- nor that the model's \emph{prediction} of the disturbed
		trajectory is exact, since \cref{eq:fc-inner} is a proportional reduction of a
		PID rate loop. What it rules out is the failure mode that motivates
		\cref{sec:disturbance}: a residual channel whose two descriptions disagree, so
		that the quantity the estimator predicts is not the quantity the optimiser
		consumes.
	\end{keybox}

	\begin{cautionbox}
		\textbf{The force block was not exact until the nominal thrust models were
		reconciled.} \Cref{eq:tauext} has always resolved the nominal moment from the
		\emph{actual} rotor speeds, $\mat M_\tau^\mathrm{nom}K_T\vec\Omega^2$, while
		\cref{eq:Fext} took its thrust from the \emph{commanded} collective $T(c)$.
		The two lines of the same definition therefore disagreed about what the
		nominal model predicts. That was invisible while the control model's own
		thrust also came from the command; once \cref{eq:fc-trans} computes it as
		$K_T\sum_i\Omega_i^2$, the force block of \cref{eq:convert} was inexact by the
		rotor-lag term --- measured at $2.1\times10^{-5}$ of the collective in settled
		flight, giving the $10^{-4}$ relative disagreement visible in the force column
		before the correction. With both lines taking rotor thrust, the force ratio
		joins the moment ratio at unity.
	\end{cautionbox}

	\subsection{Closed loop}

	\begin{table}[htbp]
		\centering
		\caption[Closed-loop RMSE by scenario and level]{Closed-loop RMSE by scenario and level (moderated). \emph{Read the
				zero-level row first}: with no disturbance attached, any gap between
			controllers is pure conservatism and not a statement about disturbance
			rejection. Source: \texttt{artifacts/acmpc\_adaptive/nb5\_scenarios.csv}.}
		\label{tab:res-scenarios}
		\small
		\begin{tabularx}{\linewidth}{@{}l r r r r r r@{}}
			\toprule
			\textbf{Scenario / level} & \textbf{Base} & \textbf{Robust} & \textbf{Oracle} & \textbf{RDP-A} & \textbf{RDP-B} & \textbf{RDP-C} \\
			\midrule
			\multicolumn{7}{@{}l}{\itshape central payload (fraction of weight)} \\
			\quad 0.000 & 0.1892 & 0.1918 & 0.1919 & 0.1934 & 0.1829 & 0.1839 \\
			\quad 0.030 & 0.1922 & 0.1983 & 0.2012 & 0.1997 & 0.1898 & 0.1902 \\
			\quad 0.060 & 0.1973 & 0.2061 & 0.2119 & 0.2083 & 0.1970 & 0.1976 \\
			\quad 0.100 & 0.2038 & 0.2169 & 0.2251 & 0.2192 & 0.2072 & 0.2090 \\
			\midrule
			\multicolumn{7}{@{}l}{\itshape arm-tip payload, offset CG} \\
			\quad 0.000 & 0.1892 & 0.1918 & 0.1919 & 0.1934 & 0.1829 & 0.1839 \\
			\quad 0.004 & 0.1907 & 0.1936 & 0.1933 & 0.1946 & 0.1848 & 0.1859 \\
			\quad 0.028 & 0.2062 & 0.2098 & 0.2059 & 0.2123 & 0.2008 & 0.2040 \\
			\quad 0.044 & 0.2220 & 0.2265 & 0.2198 & 0.2286 & 0.2170 & 0.2212 \\
			\midrule
			\multicolumn{7}{@{}l}{\itshape slung load (excitation period, \si{\second})} \\
			\quad 0.0 & 0.1892 & 0.1918 & 0.1919 & 0.1934 & 0.1829 & 0.1839 \\
			\quad 1.2 & 0.3153 & 0.3273 & 0.3434 & 0.3350 & 0.3191 & 0.3255 \\
			\quad 2.0 & 0.3142 & 0.3268 & 0.3311 & 0.3289 & 0.3185 & 0.3201 \\
			\quad 4.0 & 0.3141 & 0.3292 & 0.3451 & 0.3333 & 0.3186 & 0.3223 \\
			\quad 6.0 & 0.3158 & 0.3293 & 0.3463 & 0.3343 & 0.3201 & 0.3234 \\
			\bottomrule
		\end{tabularx}
	\end{table}

	\begin{table}[htbp]
		\centering
		\caption{The two gaps the decoupled construction of \cref{sec:adaptive} exists
			to separate, as medians over levels.}
		\label{tab:res-gaps}
		\small
		\begin{tabularx}{\linewidth}{@{}X r r@{}}
			\toprule
			\textbf{Scenario} & \textbf{Value of information} (Robust $-$ Oracle) & \textbf{Cost of estimation} (best RDP $-$ Oracle) \\
			\midrule
			central & -0.0044 & -0.0132 \\
			asym & +0.0021 & -0.0068 \\
			slung & -0.0158 & -0.0243 \\
			\bottomrule
		\end{tabularx}
	\end{table}

	\realfig{54mm}{\figClosedLoop}
	{Closed-loop RMSE against disturbance level for each variant, one panel per
		scenario. The zero-level intercept is conservatism; the slope is disturbance
		rejection. The oracle is the ceiling and the gap beneath it is what estimation
		costs.}
	{fig:closedloop}

	\subsection{Analysis}

	\textbf{The precondition held.} \Cref{tab:res-trainhealth} puts collective
	saturation at \SIrange{0.11}{0.13}{\percent} for all three policies against
	a gate of \SI{5}{\percent}, with a crash rate of zero. No policy was trained
	against its input box, so \cref{tab:res-scenarios} is readable as a
	decomposition rather than as a reading of the actuator limit. That check
	comes first because it is not repairable downstream: moderating an
	evaluation cannot undo a policy trained against the stop.

	\textbf{Read the zero-disturbance row before any gap.} At the top of
	\cref{tab:res-scenarios}, with no payload at all, Base sits at
	\SI{0.1892}{\metre}, Robust at \SI{0.1918}{\metre}, Oracle at
	\SI{0.1919}{\metre} and RDP-B at \SI{0.1829}{\metre}. The disturbance-aware
	arms are not worse than Base by any margin approaching the seed spread, and
	RDP-B is slightly better. Nothing is being rejected at zero; what differences
	exist there are conservatism.

	\textbf{Both gaps come out negative, and that is the result of this
	increment.} \Cref{tab:res-gaps} reports the value of information
	(Robust~$-$~Oracle) at \num{-0.0044} on \texttt{central}, \num{+0.0021} on
	\texttt{asym} and \num{-0.0158} on \texttt{slung}; and the cost of
	estimation (best RDP~$-$~Oracle) at \num{-0.0132}, \num{-0.0068} and
	\num{-0.0243}. A negative cost of estimation says the \emph{estimated}
	residual outperformed the \emph{exact} one on all three scenarios. The
	oracle is not the ceiling, which is not what one expects and is worth
	stating plainly. Against the \SI{0.0064}{\metre} seed spread, the
	\texttt{slung} figures (\num{-0.0158}, \num{-0.0243}) and the
	\texttt{central} cost of estimation (\num{-0.0132}) are resolved; the
	\texttt{central} value of information and the whole \texttt{asym} column are
	not, and should not be read as effects.

	Two mechanisms are consistent with a negative cost of estimation, and the
	data does not separate them. The cost map was trained without the residual
	channel carrying an exact wrench, so feeding it one moves the solve off the
	distribution the map was fitted on, while the RDP's smoothed output
	(\cref{tab:res-rdp}: GRU at $R^2=\num{0.9988}$, rolling mean over
	\vSmooth{} frames) stays closer to it. And on \texttt{slung} the exact
	wrench carries pendulum content at periods of \SIrange{1.2}{6.0}{\second}
	that a one-step horizon cannot act on, so supplying it exactly is closer to
	injecting noise than to supplying information --- which is consistent with
	\texttt{slung} showing the largest negative gap of the three.

	\textbf{Variant B wins, and where it wins says where the disturbance
	belongs.} RDP-B --- residual to the model, cost map left
	disturbance-blind --- is lowest or tied-lowest in every block of
	\cref{tab:res-scenarios}, ahead of RDP-C (both channels) everywhere and
	ahead of RDP-A (observation only) by margins that grow with the disturbance:
	\num{0.0105} at zero central payload, \num{0.0120} at the largest, and
	\num{0.0159} at the \SI{1.2}{\second} slung excitation. RDP-A, which lets
	the cost map react but keeps the model nominal, is the weakest of the three
	throughout. The reading is that \textbf{a disturbance of this kind belongs
	in the model, not in the cost} --- it is a statement about what the vehicle
	will do, which is what the prediction is for, and routing it into the
	objective instead asks the cost map to compensate for a prediction it knows
	to be wrong. That RDP-C does not beat RDP-B says the second channel adds
	nothing once the first is present, and costs a little.

	In the ledger, adaptive \ACMPC{} improves on \ACMPC{} on all three suites:
	\num{0.1040}~$\to$~\num{0.1022} on S1, \num{0.1657}~$\to$~\num{0.1561} on
	S2, \num{0.2712}~$\to$~\num{0.2534} on S3. Only the S3 improvement
	(\num{0.0178}) clearly exceeds the seed spread; S2 (\num{0.0096}) is
	marginal and S1 (\num{0.0018}) is not resolved. The effect is real where the
	disturbance is large, which is the only place it was ever claimed.

	%-----------------------------------------------------------------------
	\section{The comparison}
	\label{sec:res-comparison}
	%-----------------------------------------------------------------------

	\begin{table}[htbp]
		\centering
		\caption{The ledger. One row per controller; $t_\mathrm{solve}$ is
			\emph{single-vehicle} latency in every row. The final column counts the
			quantities a human had to choose --- the axis the learned controllers are
			trying to reduce, and the one no accuracy metric shows. Source:
			\texttt{artifacts/common/ledger.csv}.}
		\label{tab:ledger}
		\small
		\begin{tabularx}{\linewidth}{@{}X r r r r r r@{}}
			\toprule
			\textbf{Controller} & \textbf{S1} & \textbf{S2} & \textbf{S3} & \textbf{p95 [\si{\milli\second}]} & \textbf{admissible} & \textbf{tuned} \\
			\midrule
			LQR + feed-forward & 0.0945 & 0.1882 & 0.3079 & 0.42 & yes & 13 \\
			NMPC $N{=}1$ & 0.0739 & 0.1275 & 0.2368 & 12.04 & yes & 13 \\
			NMPC $N{=}10$ & 0.0699 & 0.0993 & 0.1986 & 126.69 & no & 13 \\
			\LLTC{} $N{=}1$ & 1.2547 & 1.2921 & 1.3202 & 12.20 & yes & 13 \\
			\ACMPC{} $N{=}1$ & 0.1040 & 0.1657 & 0.2712 & 12.17 & yes & \textbf{0} \\
			Adaptive \ACMPC{} $N{=}1$ & 0.1022 & 0.1561 & 0.2534 & 12.56 & yes & \textbf{0} \\
			\bottomrule
		\end{tabularx}
	\end{table}

	\begin{table}[htbp]
		\centering
		\caption{Degradation factor from S1 to S3. A controller that wins on S1 and
			loses on S3 has been tuned, not designed. Moderated-S3 factors are not
			comparable with pre-moderation ones: the suite is a different, narrower
			object. Source: \texttt{artifacts/common/nb6\_degradation.csv}.}
		\label{tab:res-degradation}
		\small
		\begin{tabularx}{\linewidth}{@{}X r r r r@{}}
			\toprule
			\textbf{Controller} & \textbf{S2/S1} & \textbf{S3/S1} & \textbf{rank S1 $\to$ S3} & \textbf{change} \\
			\midrule
			LQR + feed-forward & 1.99$\times$ & 3.26$\times$ & 3 $\to$ 5 & $-2$ \\
			NMPC $N{=}1$ & 1.73$\times$ & 3.21$\times$ & 2 $\to$ 2 & 0 \\
			NMPC $N{=}10$ & 1.42$\times$ & 2.84$\times$ & 1 $\to$ 1 & 0 \\
			\LLTC{} $N{=}1$ & 1.03$\times$ & 1.05$\times$ & 6 $\to$ 6 & 0 \\
			\ACMPC{} $N{=}1$ & 1.59$\times$ & 2.61$\times$ & 5 $\to$ 4 & $+1$ \\
			Adaptive \ACMPC{} $N{=}1$ & 1.53$\times$ & 2.48$\times$ & 4 $\to$ 3 & $+1$ \\
			\bottomrule
		\end{tabularx}
	\end{table}

	\realfigtwo{54mm}{\figHeadline}{\figRadar}
	{The comparison.}
	{Accuracy against single-vehicle latency, with the \SI{20}{\milli\second} control
		period as a vertical admissibility line and marker size showing the tuning
		burden.}
	{Radar over accuracy, robustness, smoothness, effort, latency and tuning burden
		--- the axes on which the controllers genuinely differ.}
	{fig:comparison}

	\realfig{54mm}{\figByCondition}
	{Per-condition breakdown: every controller on every path under S1, S2 and S3.
		Saturation is reported beside RMSE throughout, because a cell with high
		saturation is reporting the disturbance rather than the controller.}
	{fig:bycondition}

	\realfig{48mm}{\figFragility}
	{One bar per controller per plant axis: worst RMSE over the swept range divided
		by that controller's RMSE at nominal, so $1.0$ means the sweep never moved it.
		This is the ledger's missing column --- \cref{tab:ledger} reports accuracy under
		three disturbance suites on the \emph{nominal} plant, and says nothing about
		what happens when the plant itself is wrong. Read with
		\cref{fig:mismatch-hand,fig:mismatch-learned}, from which it is computed.}
	{fig:fragility}

	\subsection{Three questions the ledger is built to answer}

	\begin{enumerate}[leftmargin=*]
		\item \textbf{Does optimisation in the loop earn its computation?} Compare LQR
		with NMPC $N=1$: same model, same weights, same feed-forward information, so
		the difference is exactly what the online solve buys. If it is small on smooth
		paths and large on the superellipse, that is the honest answer --- the solve
		pays where curvature and constraints bind, and not elsewhere.
		\item \textbf{Does the learned cost recover long-horizon behaviour at one-step
			cost?} Compare \ACMPC{} and \LLTC{} against NMPC at $N=1$ and $N=10$, on RMSE
		\emph{and} on p95 latency. This is the only place where the claim is tested
		against both ends at once.
		\item \textbf{What survives leaving the training distribution?} The S3 column and
		\cref{tab:res-degradation}. Note that a controller may improve its \emph{rank}
		on S3 while degrading in absolute terms, simply because its competitors degrade
		faster; both readings are in the table and neither alone is the answer.
	\end{enumerate}

	%-----------------------------------------------------------------------
	\section{Error analysis, limitations and interpretability}
	%-----------------------------------------------------------------------

	\subsection{Failure modes}

	\begin{table}[htbp]
		\centering
		\caption{Failure taxonomy, with the diagnostic that distinguishes each case.
			This table is the reason several results in this chapter are reported as
			negative rather than tuned away.}
		\label{tab:failures}
		\small
		\begin{tabularx}{\linewidth}{@{}l X X@{}}
			\toprule
			\textbf{Mode} & \textbf{Signature} & \textbf{Diagnostic} \\
			\midrule
			Actuator saturation & High \texttt{sat}; all controllers converge to the same RMSE; $e_{\max}$ at the bound of \cref{def:crash} & \texttt{sat} column; compare across paths --- if the path stops mattering, the disturbance is the measurement \\
			Reference infeasibility & Uniform large error across controllers on one path only & \Cref{eq:omegacap} audit, \cref{fig:feasibility} \\
			Undertraining & Learned controllers fail while hand-built ones are unaffected & The split falls exactly along the training axis; check the scale table and iteration count \\
			Linearisation corruption & \ACMPC{} collapses at high exploration noise while the MLP still learns & Exploration sweep, \cref{tab:res-acmpc-sweeps} \\
			Terminal-cost locality & \LLTC{} good in distribution, poor out of it & OOD sweep, \cref{fig:lltcood} \\
			Estimation failure & Oracle good, RDP poor & Oracle-to-estimator gap, \cref{tab:res-gaps} \\
			Untrainable oracle & Oracle poor \emph{at zero disturbance} & \Cref{tab:res-trainhealth} --- this is not an adaptation result \\
			Measurement-window artefact & A latency figure that disagrees with the same quantity measured over a longer run & Compare a warm-up-inclusive probe against a steady-state one \\
			\bottomrule
		\end{tabularx}
	\end{table}

	\subsection{Interpretability}

	The architecture retains one property no end-to-end policy offers: the learned
	object is a \emph{cost}, and a cost can be inspected. Three diagnostics exploit
	this.

	\begin{itemize}[leftmargin=*]
		\item \textbf{Cost-map slices.} $\mat S(\vec o)$ is evaluated along a
		one-dimensional slice of the observation and its eigenvalues plotted, showing
		which error directions the policy has decided to weight, and where.
		\item \textbf{Comparison with the hand-tuned weights.} The learned diagonal is
		directly comparable with the $\mat Q,\mat R$ a human chose, which turns a
		black-box claim into a legible one. The ratio of learned to hand weight is the
		quantity to report: a ratio near one would mean the learning had merely
		rediscovered the tuning, and a ratio spanning orders of magnitude means the
		learned cost is a different object rather than a rescaling.
		\item \textbf{Recovering an economic cost.} Inverting the storage-function
		modification recovers an approximation of the underlying stage cost, which is
		in general \emph{indefinite}. That is expected rather than anomalous: the
		positive-definite cost is what the solver uses for stability, and the
		indefinite one is what the task actually rewards~\cite{gros2020data}.
	\end{itemize}

	\realfig{54mm}{\figInterpretability}
	{Interpreting the learned cost: eigenvalues of $\mat S(\vec o)$ along a
		one-dimensional slice of the observation, beside the hand-tuned $\mat Q,\mat R$
		of the NMPC baseline on the same axes.}
	{fig:costmap}

	\subsection{Audit of the experimental pipeline}
	\label{sec:audit-defects}

	A systematic audit was conducted against the two source papers and against the
	manufacturer's own vehicle description, after an inconsistency was noticed
	between a summary paragraph and the table printed immediately above it. It
	identified defects in three groups: physics that did not match the airframe file,
	learning machinery that did not match the paper it cited, and narrative claims
	contradicted by their own tables. The complete register, with severity and
	whether each defect was silent, is \cref{app:audit}.


	\begin{keybox}
		\keyhead{The methodological lesson}
		Every one of the S1 defects was silent: none raised an exception, and several
		produced plausible-looking numbers. What surfaced them was not testing the code
		against itself but confronting each result with the claim written beside it,
		each equation with its implementation, and each constant with the
		manufacturer's file. The practices adopted in response --- gating on fit
		quality rather than assuming it, asserting that a control condition actually
		differs from its treatment, reporting the sign of an effect before its
		magnitude, pinning every corrected defect with a regression test, and printing
		a verdict computed from the run instead of restating a published one --- are
		transferable, and are the part of this work most likely to outlast its numbers.
	\end{keybox}
	%=======================================================================
	\chapter{Conclusion and perspectives}
	\label{ch:conclusion}
	%=======================================================================

	%-----------------------------------------------------------------------
	\section{Summary of the work}
	%-----------------------------------------------------------------------

	This project asked whether a predictive controller and a learning algorithm can
	be combined so that each supplies what the other lacks: the optimisation problem
	provides structure, constraints and a place to reason about stability; the
	learning provides the cost that no part of predictive control theory determines.
	The answer was pursued through four increments, each isolating one mechanism and
	compared against a baseline differing in exactly one respect.

	The technical contributions are:

	\begin{enumerate}[leftmargin=*]
		\item A batched, differentiable, double-precision simulation of a Holybro X500
		in which the \emph{inner} rate loop, mixer, rotor lag and the autopilot's idle
		thrust floor are part of the plant rather than abstracted away, so that model
		error exists by construction and the motor commands --- which
		\cref{sec:rdp-theory} shows are the only observable carrying a standing moment
		once the rate integrator has absorbed it --- are available to the estimator.
		\item A differentiable iLQR layer used as the policy class, with the
		reference preview supplied along the horizon, the input box set to the
		\emph{reachable} rate set, and only the final iterations differentiated,
		justified by the parametric sensitivity result of~\cite{buskens2001sensitivity}.
		\item Two ways of learning the cost --- a positive-definite terminal cost
		(\LLTC) and a full learned cost map (\ACMPC), the latter emitting both the
		quadratic and the linear term of \cref{eq:groszanon} --- evaluated against
		classical NMPC and an LQR that receives identical information.
		\item An adaptive scheme in which a recurrent predictor reconstructs the
		six-dimensional external wrench from a \SI{1.28}{\second} window of states and
		actuator commands, trained on a position hold so that the regression target is
		the disturbance rather than the controller's own transient, and decoupled into
		an oracle stage and an estimation stage so that the result is attributable.
		\item A PX4/Gazebo ROS\,2 workspace that recomputes the same wrench from onboard
		signals only, providing a sim-to-real path with a real flight-control stack in
		the loop and a specified experimental protocol
		(\cref{sec:ros-experiments}).
	\end{enumerate}

	Two of the four increments delivered on their claim, one delivered a trade,
	and one did not deliver at all. Increment~1 failed: the terminal-cost
	regression is essentially exact ($R^2=\num{0.999998}$) and the closed loop
	is eighteen times worse than the NMPC it was meant to match, because
	$\mat P_\theta$ is a local object and a closed loop does not stay local.
	Increment~2 delivered a \emph{trade} and should be described as one: every
	hand-chosen stage weight was removed --- the \texttt{tuned} column of
	\cref{tab:ledger} goes from \num{13} to \num{0} --- at a cost of
	\SI{41}{\percent} accuracy against the tuned NMPC on S1. Increment~3
	delivered a robustness gain on S2 of \SI{6.9}{\percent} and a conservatism
	premium too small to measure at two seeds; its more useful contribution was
	negative and structural, that a policy which cannot observe its plant cannot
	specialise. Increment~4 delivered: the adaptive arm improves on \ACMPC{} on
	all three suites and clearly so on S3 (\num{0.2712}~$\to$~\num{0.2534}), and
	it produced the one finding nobody designed for --- the estimated residual
	beat the exact one on every scenario tested, so the oracle is not the
	ceiling. All of this is at \texttt{medium} scale, and differences at or
	below the \SI{0.0064}{\metre} seed spread are reported as unresolved rather
	than as results.

	%-----------------------------------------------------------------------
	\section{The endpoint: an enhanced cost and an enhanced model}
	\label{sec:endpoint}
	%-----------------------------------------------------------------------

	The four increments are easiest to see backwards, from where they end up.

	The starting point is a controller that already works. NMPC with a long horizon,
	and the LQR that anchors it, track this vehicle accurately on every path in
	\cref{sec:res-baselines} --- and they charge for it three times. They charge in
	computation, because the horizon that buys the accuracy is the horizon that buys
	the latency (\cref{sec:res-horizon}). They charge in \emph{model}, because the
	solve plans several steps against a model and therefore compounds whatever is
	wrong with it, where a fixed gain merely reacts (\cref{sec:res-sensitivity}).
	And they charge in \emph{cost}, because the thirteen weights that decide how
	they behave are not determined by any part of predictive control theory ---
	someone chooses them, and \cref{sec:res-cost} measures how much of the reported
	performance is really that choice.

	Everything after that is an attempt to stop paying, without giving up the
	structure --- the constraints, the preview, the explicit optimisation --- that
	made the classical answer worth having in the first place. The strategy
	throughout is the same: \emph{shorten the horizon to one step, and put what the
	long horizon was doing into a learned object instead}.

	\begin{description}[leftmargin=*,style=nextline]
		\item[Increment 1 put it in the terminal cost, and that was not enough.] A
		positive-definite terminal cost fitted to the long-horizon value function is
		the textbook route, and it is the one the horizon-equivalence test of
		\cref{sec:res-lltc} rejects. The fit itself is good; what it cannot survive is
		the distance between the small neighbourhood where a local quadratic is valid
		and the displacements a one-step controller actually has to cover. The
		terminal cost is the right idea in the wrong place --- the last stage is not
		where a one-step horizon spends its authority.
		\item[Increment 2 put it in the whole cost map.] \ACMPC{} makes the cost a
		\emph{function of the state} produced by a network, learned by reinforcement
		against the task rather than fitted to a value function, and emitting the
		linear term as well as the quadratic one so that it can shift the optimum and
		not merely re-weight it. The claim this increment answers is not ``more
		accurate than NMPC'': it is \emph{comparable accuracy at one-step latency with
			thirteen fewer human choices}, and \cref{tab:ledger} is where that trade is
		settled in both columns at once.
		\item[Increment 3 bought generality, and paid for it in specificity.] A cost
		map learned on one plant is a cost map for that plant. Randomising mass,
		inertia, thrust and drag during training is the standard remedy and it works
		in the direction advertised --- but the same averaging that makes the policy
		indifferent to which plant it is on makes it indifferent to the plant it is
		\emph{actually} on. That is the conservatism premium of
		\cref{sec:res-dr}, evaluated deliberately on the nominal plant and reported
		signed. Robustness by randomisation is insurance, and insurance has a premium.
		\item[Increment 4 gave it the missing observation.] The premium exists because
		the policy cannot tell the plants apart --- so the answer is not to randomise
		less but to \emph{observe more}. A recurrent predictor reconstructs the
		six-dimensional external wrench from a window of states and actuator commands,
		and the estimate is handed to the policy, to the dynamics the solver plans
		with, or to both. The construction is deliberately decoupled into an oracle
		stage and an estimation stage, so that the value of the information and the
		cost of estimating it are two separate numbers (\cref{tab:res-gaps}) rather
		than one ambiguous one.
	\end{description}

	What is left at the end is a one-step optimisation with two things replaced. The
	cost is no longer a matrix a human chose but a map $\theta$ learned against the
	task: an \textbf{enhanced cost}. The dynamics are no longer the nominal model
	alone but that model plus an online estimate $\psi$ of the wrench acting on the
	vehicle: an \textbf{enhanced model}. The optimiser between them is unchanged,
	and that is the design's whole point --- the constraints are still hard, the
	preview is still explicit, and the horizon is one step.

	Whether each replacement paid for itself is a question for the numbers, and
	\cref{ch:results} answers it increment by increment rather than in aggregate.
	The architecture in \cref{fig:architecture} is the same diagram read in this
	order.

	%-----------------------------------------------------------------------
	\section{What was learned beyond the numbers}
	%-----------------------------------------------------------------------

	Four findings generalise past this platform.

	\begin{description}[leftmargin=*,style=nextline]
		\item[A saturated actuator is a measurement failure, not a hard test.] When a
		disturbance suite is severe enough to pin the actuators, every controller
		reports the error bound and the ranking becomes an artefact. The diagnostic is
		cheap and should be routine: if the \emph{path} stops affecting the result,
		the experiment is measuring the disturbance. The contraction
		\cref{eq:moderate}, which shrinks a suite while preserving its shape, is a
		reusable remedy.
		\item[A control condition must be verified to differ from its treatment.] The
		frozen-reference column intended to demonstrate the importance of the preview
		was bit-identical to the preview column for an entire campaign, because the
		mechanism that was supposed to disable the preview did not. An assertion that
		a control differs from its treatment costs one line and would have caught it
		immediately.
		\item[Interfaces between learned and physical components need dimensional
		contracts.] The most consequential defect found was a wrench in
		\si{\newton} and \si{\newton\metre} written into a slot expecting
		\si{\metre\per\square\second} and \si{\radian\per\second}. Because it was wrong
		consistently on both sides of the train/evaluate boundary, it produced no
		error, no warning and entirely plausible numbers. Magnitude assertions at such
		interfaces are cheap insurance.
		\item[Read the manufacturer's file, not the textbook model.] Three of the
		highest-severity defects were places where a standard quadrotor model had been
		written where the airframe's own configuration said something else: rotors
		that idle rather than stop, an allocator whose moment constant deliberately
		disagrees with the plant, and a rate command the autopilot clips well inside
		the nominal box. None of them is exotic, and all three changed results.
	\end{description}

	%-----------------------------------------------------------------------
	\section{Limitations}
	%-----------------------------------------------------------------------

	\begin{enumerate}[leftmargin=*]
		\item \textbf{The full-scale campaign has not been run.} This is the principal
		limitation and it bounds every other claim: \cref{ch:results} is a specified
		experiment, not a completed one. The pipeline executes end to end and its
		gates behave, but the learned-controller numbers require \texttt{full} scale
		before they can be read at all.
		\item \textbf{No hardware flight.} All results are simulated. The ROS\,2
		workspace is the bridge and its protocol is specified, but a physical airframe
		introduces effects --- aerodynamic interaction, structural flexibility, sensor
		latency, battery sag --- that no rigid-body model contains.
		\item \textbf{No formal guarantee.} The architecture is compatible with the
		stability machinery of economic MPC --- a storage-function modification of the
		cost preserves the optimal policy while rendering the stage cost positive
		definite, which yields a Lyapunov function~\cite{diehl2011lyapunov,
			angeli2012average} --- but no certificate is produced here.
		\item \textbf{Local optimality.} Policy-gradient methods on a nonlinear function
		approximator give no guarantee of a global optimum. This limitation is shared
		with essentially all deep RL and is not specific to the MPC
		parametrisation~\cite{gros2020data}.
		\item \textbf{A limited cost parametrisation.} A rich parametrisation is
		desirable in theory, but stage and terminal costs must stay quasi-convex to
		keep the inner solve tractable. Positive sums of convex functions and
		sum-of-squares constructions are the natural next step and were not explored.
		\item \textbf{The rate loop is reduced, not reproduced.} \Cref{eq:fc-inner}
		is proportional, so the model does not carry the autopilot's integrator. A
		moment that integrator has already absorbed is still predicted as a persistent
		angular acceleration, which bounds how well any residual estimate can describe
		a long-standing disturbance. Restoring $\vec I_\omega$ would close the gap at
		the cost of three further states and the steady-state gain of
		\cref{prop:dcgain}.
		\item \textbf{Computational budget.} Seed sweeps are limited to three, which is
		too few to separate some of the competing explanations \cref{sec:res-dr} has
		to consider.
	\end{enumerate}

	%-----------------------------------------------------------------------
	\section{Perspectives}
	%-----------------------------------------------------------------------

	\subsection{Short term --- completing what is specified}

	\begin{itemize}[leftmargin=*]
		\item Run the campaign at \texttt{full} scale and fill \cref{ch:results}. The
		\texttt{medium} rung exists for exactly this reason: it is roughly
		$427\times$ the \texttt{smoke} training work against $42\,667\times$ for
		\texttt{full}, and it will establish whether the learned controllers converge
		at all before a multi-week run is committed to.
		\item Execute the ROS\,2 campaign of \cref{sec:ros-experiments}, beginning with
		the zero-level correctness check, and compare the predicted steady-state wrench
		with the analytic values of \cref{eq:analytic} --- a physical prediction with
		no free parameters, and therefore the most convincing available validation.
		\item Sweep the predictor history length (experiment E\nbd D) around the
		chosen $H=\vH$, retraining on the position-hold set at each value, so that the
		choice is shown to be a maximum rather than assumed to be one.
	\end{itemize}

	\subsection{Medium term --- guarantees}

	\begin{itemize}[leftmargin=*]
		\item \textbf{Robust MPC with a scenario tree.} A scenario tree is a stochastic
		model with a discrete distribution, so the whole construction applies
		unchanged, with the branch probabilities themselves becoming learnable
		parameters~\cite{gros2020data}. This is the natural route to handling
		safety-critical constraints, which must be excluded by the formulation rather
		than penalised in the cost.
		\item \textbf{Dissipativity-based stability.} If a stage cost yields a
		stabilising optimal policy, then a scheme with a positive-definite stage cost
		can in principle be tuned to deliver that policy, with the storage function
		supplying the Lyapunov argument~\cite{diehl2011lyapunov,muller2015necessity}.
		Instantiating this for the learned cost map would convert an empirical result
		into a certified one.
		\item \textbf{Second-order and policy-gradient learning rules.} Temporal-difference
		learning on a nonlinear parametrisation is a proxy for policy quality: fitting
		$Q_\theta\approx Q_\star$ does not guarantee $\pi_\theta\approx\pi_\star$.
		Least-squares TD and deterministic policy gradient~\cite{silver2014dpg}
		address this directly.
	\end{itemize}

	\subsection{Longer term --- deployment}

	\begin{itemize}[leftmargin=*]
		\item \textbf{Online adaptation of the predictor.} The Neural-Fly
		decomposition~\cite{oconnell2022neuralfly} --- a slowly learned basis with
		rapidly adapted linear coefficients --- fits this architecture directly and
		would let the estimator track conditions it never saw in training.
		\item \textbf{Embedded execution.} The admissibility criterion is a
		\SI{20}{\milli\second} period on the flight computer, not $R^2$. The exported
		predictors are already framework-neutral NumPy with an asserted parity to the
		trained model; the remaining work is quantised inference and a
		bounded-latency solve. The three-kernel-per-iteration structure reported for
		CUDA-accelerated differentiable MPC is a promising route for the solve itself.
		\item \textbf{Hardware validation.} Flights on a physical X500 with instrumented
		payloads, comparing the estimated wrench against a force--torque measurement
		and closing the loop on it.
	\end{itemize}

	\begin{keybox}
		\keyhead{Closing remark}
		The premise this project set out to test --- that model error can be absorbed
		into the cost, so an MPC on a wrong model can still act optimally --- is a
		theorem, not a hope~\cite{gros2020data}. What this work contributes is an
		account of what it takes to realise it on a real platform: the interfaces that
		must be dimensionally consistent, the constants that must be read from the
		airframe rather than from a textbook, the disturbances that must be severe
		enough to matter but not so severe that the actuators decide the outcome, the
		control conditions that must be verified to differ from their treatments, and
		the discipline of reporting what the tables say rather than what the method was
		expected to show.
	\end{keybox}

	%=======================================================================
	%  BIBLIOGRAPHY
	%=======================================================================
	\cleardoublepage
	\addcontentsline{toc}{chapter}{Bibliography}
	\begin{thebibliography}{99}
		\small

		\bibitem{gros2020data}
		S.~Gros and M.~Zanon.
		\newblock Data-driven economic NMPC using reinforcement learning.
		\newblock \emph{IEEE Transactions on Automatic Control}, 65(2):636--648, 2020.
		\newblock arXiv:1904.04152.

		\bibitem{romero2024acmpc}
		A.~Romero, Y.~Song, and D.~Scaramuzza.
		\newblock Actor-critic model predictive control.
		\newblock In \emph{IEEE International Conference on Robotics and Automation
			(ICRA)}, 2024. arXiv:2306.09852.

		\bibitem{zanon2021safe}
		M.~Zanon and S.~Gros.
		\newblock Safe reinforcement learning using robust MPC.
		\newblock \emph{IEEE Transactions on Automatic Control}, 66(8):3638--3652, 2021.

		\bibitem{amos2018differentiable}
		B.~Amos, I.~D.~J. Rodriguez, J.~Sacks, B.~Boots, and J.~Z. Kolter.
		\newblock Differentiable MPC for end-to-end planning and control.
		\newblock In \emph{Advances in Neural Information Processing Systems (NeurIPS)},
		2018. arXiv:1810.13400.

		\bibitem{li2004ilqr}
		W.~Li and E.~Todorov.
		\newblock Iterative linear quadratic regulator design for nonlinear biological
		movement systems.
		\newblock In \emph{International Conference on Informatics in Control,
			Automation and Robotics (ICINCO)}, 2004.

		\bibitem{tassa2012synthesis}
		Y.~Tassa, T.~Erez, and E.~Todorov.
		\newblock Synthesis and stabilization of complex behaviors through online
		trajectory optimization.
		\newblock In \emph{IEEE/RSJ International Conference on Intelligent Robots and
			Systems (IROS)}, 2012.

		\bibitem{lowrey2019polo}
		K.~Lowrey, A.~Rajeswaran, S.~Kakade, E.~Todorov, and I.~Mordatch.
		\newblock Plan online, learn offline: Efficient learning and exploration via
		model-based control.
		\newblock In \emph{International Conference on Learning Representations (ICLR)},
		2019.

		\bibitem{bhardwaj2021blending}
		M.~Bhardwaj, S.~Choudhury, and B.~Boots.
		\newblock Blending MPC and value function approximation for efficient
		reinforcement learning.
		\newblock In \emph{International Conference on Learning Representations (ICLR)},
		2021.

		\bibitem{rawlings2017mpc}
		J.~B. Rawlings, D.~Q. Mayne, and M.~M. Diehl.
		\newblock \emph{Model Predictive Control: Theory, Computation, and Design}.
		\newblock Nob Hill Publishing, 2nd edition, 2017.

		\bibitem{mayne2000constrained}
		D.~Q. Mayne, J.~B. Rawlings, C.~V. Rao, and P.~O.~M. Scokaert.
		\newblock Constrained model predictive control: Stability and optimality.
		\newblock \emph{Automatica}, 36(6):789--814, 2000.

		\bibitem{scokaert1999feasibility}
		P.~O.~M. Scokaert and J.~B. Rawlings.
		\newblock Feasibility issues in linear model predictive control.
		\newblock \emph{AIChE Journal}, 45(8):1649--1659, 1999.

		\bibitem{buskens2001sensitivity}
		C.~B\"uskens and H.~Maurer.
		\newblock Sensitivity analysis and real-time optimization of parametric
		nonlinear programming problems.
		\newblock In \emph{Online Optimization of Large Scale Systems}, pages 3--16.
		Springer, 2001.

		\bibitem{diehl2011lyapunov}
		M.~Diehl, R.~Amrit, and J.~B. Rawlings.
		\newblock A Lyapunov function for economic optimizing model predictive control.
		\newblock \emph{IEEE Transactions on Automatic Control}, 56(3):703--707, 2011.

		\bibitem{angeli2012average}
		D.~Angeli, R.~Amrit, and J.~B. Rawlings.
		\newblock On average performance and stability of economic model predictive
		control.
		\newblock \emph{IEEE Transactions on Automatic Control}, 57(7):1615--1626, 2012.

		\bibitem{muller2015necessity}
		M.~A. M\"uller, D.~Angeli, and F.~Allg\"ower.
		\newblock On necessity and robustness of dissipativity in economic model
		predictive control.
		\newblock \emph{IEEE Transactions on Automatic Control}, 60(6):1671--1676, 2015.

		\bibitem{marchetti2009modifier}
		A.~Marchetti, B.~Chachuat, and D.~Bonvin.
		\newblock Modifier-adaptation methodology for real-time optimization.
		\newblock \emph{Industrial \& Engineering Chemistry Research},
		48(13):6022--6033, 2009.

		\bibitem{schulman2017ppo}
		J.~Schulman, F.~Wolski, P.~Dhariwal, A.~Radford, and O.~Klimov.
		\newblock Proximal policy optimization algorithms.
		\newblock arXiv:1707.06347, 2017.

		\bibitem{schulman2016gae}
		J.~Schulman, P.~Moritz, S.~Levine, M.~I. Jordan, and P.~Abbeel.
		\newblock High-dimensional continuous control using generalized advantage
		estimation.
		\newblock In \emph{International Conference on Learning Representations (ICLR)},
		2016. arXiv:1506.02438.

		\bibitem{silver2014dpg}
		D.~Silver, G.~Lever, N.~Heess, T.~Degris, D.~Wierstra, and M.~Riedmiller.
		\newblock Deterministic policy gradient algorithms.
		\newblock In \emph{International Conference on Machine Learning (ICML)}, 2014.

		\bibitem{sutton2018rl}
		R.~S. Sutton and A.~G. Barto.
		\newblock \emph{Reinforcement Learning: An Introduction}.
		\newblock MIT Press, 2nd edition, 2018.

		\bibitem{kingma2015adam}
		D.~P. Kingma and J.~Ba.
		\newblock Adam: A method for stochastic optimization.
		\newblock In \emph{International Conference on Learning Representations (ICLR)},
		2015. arXiv:1412.6980.

		\bibitem{garcia2015safe}
		J.~Garc\'ia and F.~Fern\'andez.
		\newblock A comprehensive survey on safe reinforcement learning.
		\newblock \emph{Journal of Machine Learning Research}, 16:1437--1480, 2015.

		\bibitem{berkenkamp2017safe}
		F.~Berkenkamp, M.~Turchetta, A.~P. Schoellig, and A.~Krause.
		\newblock Safe model-based reinforcement learning with stability guarantees.
		\newblock In \emph{Advances in Neural Information Processing Systems (NeurIPS)},
		2017.

		\bibitem{mahony2012multirotor}
		R.~Mahony, V.~Kumar, and P.~Corke.
		\newblock Multirotor aerial vehicles: Modeling, estimation, and control of
		quadrotor.
		\newblock \emph{IEEE Robotics \& Automation Magazine}, 19(3):20--32, 2012.

		\bibitem{mellinger2011minimum}
		D.~Mellinger and V.~Kumar.
		\newblock Minimum snap trajectory generation and control for quadrotors.
		\newblock In \emph{IEEE International Conference on Robotics and Automation
			(ICRA)}, 2011.

		\bibitem{faessler2018differential}
		M.~Faessler, A.~Franchi, and D.~Scaramuzza.
		\newblock Differential flatness of quadrotor dynamics subject to rotor drag for
		accurate tracking of high-speed trajectories.
		\newblock \emph{IEEE Robotics and Automation Letters}, 3(2):620--626, 2018.

		\bibitem{hwangbo2017quadrl}
		J.~Hwangbo, I.~Sa, R.~Siegwart, and M.~Hutter.
		\newblock Control of a quadrotor with reinforcement learning.
		\newblock \emph{IEEE Robotics and Automation Letters}, 2(4):2096--2103, 2017.

		\bibitem{kaufmann2023champion}
		E.~Kaufmann, L.~Bauersfeld, A.~Loquercio, M.~M\"uller, V.~Koltun, and
		D.~Scaramuzza.
		\newblock Champion-level drone racing using deep reinforcement learning.
		\newblock \emph{Nature}, 620:982--987, 2023.

		\bibitem{shi2019neurallander}
		G.~Shi, X.~Shi, M.~O'Connell, R.~Yu, K.~Azizzadenesheli, A.~Anandkumar, Y.~Yue,
		and S.-J. Chung.
		\newblock Neural lander: Stable drone landing control using learned dynamics.
		\newblock In \emph{IEEE International Conference on Robotics and Automation
			(ICRA)}, 2019.

		\bibitem{oconnell2022neuralfly}
		M.~O'Connell, G.~Shi, X.~Shi, K.~Azizzadenesheli, A.~Anandkumar, Y.~Yue, and
		S.-J. Chung.
		\newblock Neural-Fly enables rapid learning for agile flight in strong winds.
		\newblock \emph{Science Robotics}, 7(66):eabm6597, 2022.

		\bibitem{torrente2021gpmpc}
		G.~Torrente, E.~Kaufmann, P.~F\"ohn, and D.~Scaramuzza.
		\newblock Data-driven MPC for quadrotors.
		\newblock \emph{IEEE Robotics and Automation Letters}, 6(2):3769--3776, 2021.

		\bibitem{saviolo2022physics}
		A.~Saviolo, G.~Li, and G.~Loianno.
		\newblock Physics-inspired temporal learning of quadrotor dynamics for accurate
		model predictive trajectory tracking.
		\newblock \emph{IEEE Robotics and Automation Letters}, 7(4):10256--10263, 2022.

		\bibitem{cho2014gru}
		K.~Cho, B.~van Merri\"enboer, C.~Gulcehre, D.~Bahdanau, F.~Bougares, H.~Schwenk,
		and Y.~Bengio.
		\newblock Learning phrase representations using RNN encoder--decoder for
		statistical machine translation.
		\newblock In \emph{EMNLP}, 2014. arXiv:1406.1078.

		\bibitem{hochreiter1997lstm}
		S.~Hochreiter and J.~Schmidhuber.
		\newblock Long short-term memory.
		\newblock \emph{Neural Computation}, 9(8):1735--1780, 1997.

		\bibitem{bai2018tcn}
		S.~Bai, J.~Z. Kolter, and V.~Koltun.
		\newblock An empirical evaluation of generic convolutional and recurrent
		networks for sequence modeling.
		\newblock arXiv:1803.01271, 2018.

		\bibitem{tobin2017domain}
		J.~Tobin, R.~Fong, A.~Ray, J.~Schneider, W.~Zaremba, and P.~Abbeel.
		\newblock Domain randomization for transferring deep neural networks from
		simulation to the real world.
		\newblock In \emph{IEEE/RSJ International Conference on Intelligent Robots and
			Systems (IROS)}, 2017.

		\bibitem{peng2018simtoreal}
		X.~B. Peng, M.~Andrychowicz, W.~Zaremba, and P.~Abbeel.
		\newblock Sim-to-real transfer of robotic control with dynamics randomization.
		\newblock In \emph{IEEE International Conference on Robotics and Automation
			(ICRA)}, 2018.

		\bibitem{meier2015px4}
		L.~Meier, D.~Honegger, and M.~Pollefeys.
		\newblock PX4: A node-based multithreaded open source robotics framework for
		deeply embedded platforms.
		\newblock In \emph{IEEE International Conference on Robotics and Automation
			(ICRA)}, 2015.

		\bibitem{koenig2004gazebo}
		N.~Koenig and A.~Howard.
		\newblock Design and use paradigms for Gazebo, an open-source multi-robot
		simulator.
		\newblock In \emph{IEEE/RSJ International Conference on Intelligent Robots and
			Systems (IROS)}, 2004.

		\bibitem{furrer2016rotors}
		F.~Furrer, M.~Burri, M.~Achtelik, and R.~Siegwart.
		\newblock RotorS --- a modular Gazebo MAV simulator framework.
		\newblock In \emph{Robot Operating System (ROS): The Complete Reference,
			Volume~1}, pages 595--625. Springer, 2016.

		\bibitem{jax2018github}
		J.~Bradbury, R.~Frostig, P.~Hawkins, M.~J. Johnson, C.~Leary, D.~Maclaurin,
		G.~Necula, A.~Paszke, J.~VanderPlas, S.~Wanderman-Milne, and Q.~Zhang.
		\newblock JAX: Composable transformations of Python+NumPy programs, 2018.
		\newblock \url{http://github.com/jax-ml/jax}.

	\end{thebibliography}

	%=======================================================================
	%  APPENDICES
	%=======================================================================
	\cleardoublepage
	\appendix
	\renewcommand{\chaptermark}[1]{\markboth{\appendixname\ \thechapter.\ #1}{}}

	%-----------------------------------------------------------------------
	\chapter{Supplementary derivations}
	\label{app:derivations}
	%-----------------------------------------------------------------------

	\section{Hover linearisation with the idle-floor thrust map}

	Let $\delta\inp=\inp-\inp_\refsym$ with
	$\inp_\refsym=[u_\mathrm{hover},0,0,0]$. From \cref{eq:thrustmap},
	$T(c)=T_{\max}(r+c(1-r))^2$ with $r=\Omega_{\min}/\Omega_{\max}$, so hover
	requires $r+u_\mathrm{hover}(1-r)=\sqrt{mg/T_{\max}}$, giving
	\begin{equation}
		u_\mathrm{hover}=\frac{\sqrt{mg/T_{\max}}-r}{1-r}=\vUhover,
	\end{equation}
	and the collective authority follows by the chain rule,
	\begin{equation}
		\frac{\partial a_z}{\partial\delta u_0}
		=\frac{2T_{\max}(1-r)\big(r+u_\mathrm{hover}(1-r)\big)}{m}
		=\frac{2(1-r)}{m}\sqrt{mg\,T_{\max}}
		=\vDadc\ \si{\per\square\second}.
	\end{equation}
	The pure-square assumption $r=0$ gives $2g/\sqrt{mg/T_{\max}}=\SI{25.49}{\per\square\second}$
	instead, an overstatement of \SI{17.6}{\percent}; the linear-thrust assumption
	$T=T_{\max}c$ gives $T_{\max}/m=\SI{16.56}{\per\square\second}$, an
	understatement of \SI{23.6}{\percent}. Both errors propagate into $\mat P_T$ and
	the LQR gain.

	Continuous hover linearisation in the error coordinates of \cref{eq:err}:
	\begin{equation}
		\mat A_c=\begin{bmatrix}
			\mat 0 & \mat I_3 & \mat 0 & \mat 0 & \mat 0\\
			\mat 0 & \mat 0 & +g\,\mat\Xi & \mat 0 & \ast\\
			\mat 0 & \mat 0 & \mat 0 & \mat I_3 & \mat 0\\
			\mat 0 & \mat 0 & \mat 0 & \mat 0 & \ast\\
			\mat 0 & \mat 0 & \mat 0 & \ast & \ast
		\end{bmatrix},
		\qquad
		\mat B_c=\begin{bmatrix}\mat 0\\ \mat 0\\ \mat 0\\ \mat 0\\ \ast\end{bmatrix},
		\qquad
		\mat\Xi=\begin{bmatrix}0&1&0\\-1&0&0\\0&0&0\end{bmatrix},
	\end{equation}
	where $\ast$ denotes a nonzero block and the row order follows \cref{eq:err}.
	Two structural facts are worth reading off. The velocity/attitude block carries
	$+g\,\mat\Xi$: with $\vec\delta=2\vec q_{e,v}$ one has
	$\vec z_b\approx\vec e_3+\vec\delta\times\vec e_3$, and the opposite sign yields
	an LQR gain whose closed-loop spectral radius on the true plant exceeds unity.
	And $\dot{\vec\delta}=\bodyrate$ places the attitude/rate coupling in
	$\mat A_c$, so by \cref{prop:cascade} $\mat B_c$ is zero on every row but the
	rotor block --- the input reaches attitude only through two integrations. The
	gains that a direct attitude input would have stood for survive as the
	steady-state statements of \cref{prop:dcgain}, and it is those that are asserted
	numerically at import time, together with the $\mat A_c$ blocks above.

	The system is discretised exactly by the matrix exponential at $\Delta t_c$,
	after which $\mat P_T$ and $\mat K_\mathrm{LQR}$ follow from the discrete
	algebraic Riccati equation.

	\section{Peak demand of the superellipse}

	With $r(\theta)=f(\theta)^{-1/8}$ and $f=\cos^8\theta+\sin^8\theta$,
	\begin{align}
		f' &= 8\,(\sin^7\theta\cos\theta-\cos^7\theta\sin\theta),\\
		r' &= -f^{-9/8}\,(\sin^7\theta\cos\theta-\cos^7\theta\sin\theta),\\
		r'' &= 9f^{-17/8}(\cdot)^2-f^{-9/8}(\cdot)',
	\end{align}
	and $f\in[2^{-3},1]$ never vanishes, so $r$ is smooth everywhere --- which is why
	the superellipse is used instead of a true square, whose unbounded corner
	curvature would make the reference attitude \cref{eq:qref} discontinuous.
	Position, velocity and acceleration in the plane follow as
	\begin{align}
		\vec p_\refsym &= R\,r\,(\cos\theta,\ \sin\theta,\ 0),\\
		\vec v_\refsym &= R\omega\big((r'\cos\theta-r\sin\theta),\ (r'\sin\theta+r\cos\theta),\ 0\big),\\
		\vec a_\refsym &= R\omega^2\big((r''\cos\theta-2r'\sin\theta-r\cos\theta),\
		(r''\sin\theta+2r'\cos\theta-r\sin\theta),\ 0\big),
	\end{align}
	whose peak at $R=\omega=1$ is $\kappa_a=\num{9.0779}$ by dense sampling, a value
	the physics self-test asserts.

	\section{Exactness of the \texorpdfstring{$\ell_1$}{l1} relaxation}

	\begin{proposition}
		Let $\vec\mu^\star$ be the optimal multipliers of the hard-constrained problem
		\cref{eq:ocp} at a feasible initial state. If the penalty weights satisfy
		$w_j>\mu_j^\star$ for all $j$, then the relaxed problem \cref{eq:l1relax} has
		$\vec\sigma^\star=\vec 0$ and the two problems share a minimiser.
	\end{proposition}

	\begin{proof}[Sketch]
		This is the standard exact-penalty argument~\cite{scokaert1999feasibility}.
		The relaxed problem's optimality conditions in $\vec\sigma$ give
		$0\in\partial(w_j\sigma_j)-\mu_j$ at $\sigma_j>0$, which requires
		$\mu_j=w_j$; since $\mu_j\le\mu_j^\star<w_j$ by assumption, no positive
		$\sigma_j$ is optimal. Hence $\vec\sigma^\star=\vec 0$, the constraint is
		satisfied exactly, and the objectives coincide there.
	\end{proof}

	\begin{remark}
		The relaxation is nevertheless necessary for the \emph{learning} to be defined,
		independently of exactness: an infinite value function cannot enter a
		temporal-difference update. Safety-critical constraints must therefore be
		excluded by the formulation --- for instance by a robust or tube-based
		scheme~\cite{zanon2021safe} --- rather than represented as a large penalty.
	\end{remark}

	%-----------------------------------------------------------------------
	\chapter{Audit register}
	\label{app:audit}
	%-----------------------------------------------------------------------

	Severity S1 changes results, S2 changes reported numbers, S3 concerns
	correctness of the record. ``Silent'' indicates that the defect raised no
	exception and produced plausible output. Every correction listed here is pinned
	by a regression test in \texttt{study/tests/}.

	\begin{table}[htbp]
		\centering
		\caption{Physics against the manufacturer's airframe description.}
		\label{tab:audit-physics}
		\scriptsize
		\begin{tabularx}{\linewidth}{@{}l c c X@{}}
			\toprule
			\textbf{Id} & \textbf{Sev.} & \textbf{Silent} & \textbf{Description and resolution} \\
			\midrule
			A-2 & S1 & yes & The PX4 idle floor was ignored: $T=T_{\max}c^2$ rather than \cref{eq:thrustmap}. Hover moved from \num{0.769} to \vUhover, collective authority from \SI{25.49}{} to \vDadc\,\si{\per\square\second}, and the model asserted a zero-thrust state the vehicle cannot reach. Affine-then-quadratic map throughout; both constants asserted by the self-test. \\
			A-3 & S1 & yes & One $k_m$ served both the allocator and the plant, deleting a disagreement that exists on the real vehicle. Separate $k_m^\mathrm{ctrl}=\vKmCtrl$ from the SDF rotor-drag coefficient; $\mathrm{cond}(\mat M)=\vCondPlant$ against $\mathrm{cond}(\mat M_\mathrm{ctrl})=\vCondCtrl$, and roughly a third of a commanded yaw moment is realised. \\
			A-4 & S2 & yes & PWM normalisation did not match \texttt{actuator\_motors}. Derived from one place and asserted across the study/workspace boundary. \\
			A-5 & S2 & yes & SDF rotor drag and rolling moment were absent from the plant. Added to \crefrange{eq:fp-trans}{eq:fp-rot}. \\
			A-6 & S2 & yes & Composite inertia used body-only values. Recomputed with the rotor masses; $(\vJxx,\vJyy,\vJzz)$. \\
			D-7 & S1 & yes & The optimiser's input box was the nominal rate box, \SI{62}{\percent} of whose roll/pitch extent the autopilot clips away. Box shrunk to the reachable set. Applying the clip inside the dynamics instead leaves zero derivative beyond the knee and singularises $\mat Q_{uu}$; measured, that produced \SI{47}{\percent} saturation in an untrained policy. \\
			P-1 & S3 & yes & The helix peak-speed constant ignored the climb term, \SI{3}{\percent} low. Corrected and its validity range annotated. \\
			\bottomrule
		\end{tabularx}
	\end{table}

	\begin{table}[htbp]
		\centering
		\caption{Learning machinery against the papers it cites.}
		\label{tab:audit-learning}
		\scriptsize
		\begin{tabularx}{\linewidth}{@{}l c c X@{}}
			\toprule
			\textbf{Id} & \textbf{Sev.} & \textbf{Silent} & \textbf{Description and resolution} \\
			\midrule
			B-1 & S1 & yes & Exploration noise never reached the action, so the PPO ratio \cref{eq:ppo} compared two densities that generated nothing. Sample $a=\mu_\theta+\sigma\varepsilon$ on the action. This alone changed the character of the results: with the noise where it belongs, a short training budget genuinely fails rather than appearing to succeed. \\
			B-2 & S1 & yes & The linear cost term $\vec c_k$ was never learned, discarding the part of \cref{eq:groszanon} that absorbs a model \emph{bias}. Both $(\mat S_k,\vec c_k)$ are now emitted. \\
			B-4 & S2 & partly & The observation normaliser had no variance floor, and near-constant channels amplified the first moving sample by $\sim10^4$, reaching \texttt{NaN} within three iterations. The model-free arm was additionally unnormalised while the cost map was, so the exploration comparison measured the normalisation rather than the architecture. Floor, clip, and normalise both arms identically. \\
			B-5 & S2 & no & The exploration standard deviation was fixed rather than learned. $\log\sigma$ is now a parameter. \\
			--- & S1 & yes & Because the policy mean is the solution of an optimisation problem, a single clipped step moved it far more than the parameter change suggested; the sampled KL reached order $10$. A KL trip that reverts the update, plus an adaptive learning rate, was added. \\
			\bottomrule
		\end{tabularx}
	\end{table}

	\begin{table}[htbp]
		\centering
		\caption{Disturbance interface, experiment design and the record.}
		\label{tab:audit-rest}
		\scriptsize
		\begin{tabularx}{\linewidth}{@{}l c c X@{}}
			\toprule
			\textbf{Id} & \textbf{Sev.} & \textbf{Silent} & \textbf{Description and resolution} \\
			\midrule
			D-1 & S1 & yes & The wrench (\si{\newton},\si{\newton\metre}) was supplied to the residual slot (\si{\metre\per\square\second},\si{\radian\per\second}); force high by the factor $m$, moment dimensionally unrelated. Wrong consistently on both sides of the train/evaluate boundary, hence no error message. Route through \cref{prop:convert}; assert magnitude bounds. \\
			D-3 & S1 & yes & The conversion discarded the moment block entirely. Adopt \cref{eq:convert}, which is exact in both blocks because \cref{eq:fc-rot} carries a torque input, and measure it against the true residual (\cref{tab:res-dmod}). \\
			D-5 & S2 & yes & The helix reference was inconsistent when the feasibility cap bound. Demand and realisation now agree by construction. \\
			M-1 & S1 & yes & The allocator was rebuilt from the shifted CG, giving the inner loop exact payload knowledge unlike a real allocator. Separate plant truth from controller belief. \\
			M-2 & S1 & partly & Disturbance randomisation at \SI{32}{\percent} of weight saturated training. Reduced, and the saturation fraction became a gate. \\
			A-1 & S1 & yes & The ``frozen reference'' control was not frozen; its column was bit-identical to the treatment. Build a controller that cannot bind the environment; assert the columns differ. \\
			C-1 & S2 & no & A summary asserted horizon equivalence while the table above it showed the opposite. Gate on fit quality and print the ratio. \\
			C-2 & S2 & no & A summary restated a published parametrisation finding the run did not reproduce. Print the verdict computed from the run. \\
			C-4 & S2 & no & The conservatism premium was reported with the wrong sign. Report the sign before the magnitude. \\
			C-5 & S2 & yes & The predictor window and hidden sizes did not match the source work, and the prescribed output smoothing was absent. $H=\vH$, two layers of $64$, rolling mean over $\vSmooth$ frames. \\
			C-6 & S2 & yes & Position hold was a frozen clock rather than a fixed point, so the reference still carried velocity and acceleration. $\vec v_\refsym=\vec a_\refsym=\vec 0$ under hold. \\
			E-1 & S2 & yes & The attitude-error inverse used a first-order approximation; \SI{-2.7}{\degree} at \SI{60}{\degree} tilt. Apply \cref{lem:inverse}. \\
			E-3 & S3 & yes & \texttt{EnvCfg.task} was a dead field. It now selects the training objective. \\
			E-4 & S3 & no & The residual target of the oracle channel was never exercised. Both targets are now reachable and tabled. \\
			B-1' & S2 & yes & Batched throughput and single-vehicle latency were reported in one column, implying a large spurious speed-up. One definition everywhere. \\
			S-1 & S3 & no & A specification claimed full disjointness of the in- and out-of-distribution suites; they overlap on four channels. Reworded, and the per-channel table is printed at run time. \\
			V-1 & S3 & no & Clip RMSE used the second half while the maximum used the whole clip. Whole clip for both. \\
			V-2 & S3 & partly & Wind statistics differed between clips sharing a spec and seed, consistent with a mid-clip respawn. Respawns are recorded per clip. \\
			\bottomrule
		\end{tabularx}
	\end{table}

	%-----------------------------------------------------------------------
	\chapter{Reproduction}
	\label{app:repro}
	%-----------------------------------------------------------------------

	\section{Environment}

	\begin{itemize}[nosep]
		\item \textbf{Study half:} Python~3.11, JAX (CPU or CUDA backend), NumPy,
		pandas, Matplotlib, \texttt{ffmpeg}. \texttt{float64} throughout. Verified on
		jax/jaxlib \texttt{0.10.2} on both a CPU-only build and a CUDA build.
		\item \textbf{Robotics half:} ROS\,2, PX4-Autopilot SITL, Gazebo,
		\texttt{px4\_msgs}. Builds without JAX and without PyTorch; the deployed
		predictor is pure NumPy.
	\end{itemize}

	\section{Commands}

	\begin{lstlisting}[language=bash, caption={End-to-end reproduction.}]
# ---- study: the full campaign
cd study/notebooks
export X500_SCALE=full        # smoke = wiring check only; medium = intermediate
for n in nb?_*.py; do python "$n" || break; done

# ---- verification, no GPU required
cd study
python -m pytest tests/ -q          # 86 tests
python check_consistency.py         # study <-> workspace constants and models

# ---- export the predictor for deployment (asserts JAX/NumPy parity)
python export_estimator.py --arch GRU --out ../rdp_acmpc_ws/models/rdp_gru.npz

# ---- robotics, no PX4 required for the self-tests
cd rdp_acmpc_ws/src/acmpc_controller
python -m acmpc_controller.check_ctbr      # action mapping
python -m acmpc_controller.check_glue      # frames, B_d bridge, exported models

# ---- the disturbance campaign (134 planned runs; E-0 first)
python -m experiment_manager.runner --plan full
	\end{lstlisting}

	\section{Artefacts}

	Every increment writes versioned artefacts so that any figure or table can be
	regenerated without rerunning the training that produced it: checkpoints under
	\texttt{artifacts/} in one subdirectory per increment (\texttt{lltc},
	\texttt{acmpc}, \texttt{domrand}, \texttt{acmpc\_adaptive}); tabular results as
	CSV under \texttt{artifacts/common/}; figures under
	\texttt{artifacts/figures/}; and flight clips with a manifest under
	\texttt{artifacts/videos/}.

	This document reads figures directly from \texttt{artifacts/figures/} by path,
	so re-running the study updates them in place with no edit here. The tables in
	\cref{ch:results} name their source CSV in the caption, so each blank cell has
	an unambiguous origin.

	\begin{cautionbox}
		\textbf{Scale discipline.} The \texttt{smoke} configuration exists to verify
		that the pipeline executes; it completes a campaign in minutes and produces
		undertrained policies. Its learned-controller numbers are not results and must
		never be quoted as such. A diagnostic in the reporting stage detects the
		characteristic signature --- hand-built controllers unaffected while every
		learned controller diverges --- and prints a warning naming the scale table.
		\Cref{ch:results} is left blank rather than filled from \texttt{smoke} for
		exactly this reason.
	\end{cautionbox}

\end{document}
